% ============================================================
%  R. Nielsen, "Philosophisk Propædeutik i Grundtræk."
%  Kjøbenhavn: Forlagt af den Gyldendalske Boghandling (F. Hegel);
%  Thieles Bogtrykkeri, 1857.
%  TRANSCRIPTION (Danish) — from the Google Books scan.
%
%  Scan: ~/bibliotek/Nielsen, Rasmus/PPG-1857.pdf (220 PDF pp., Google Books;
%        an identical-title twin, Philosophisk_Propædeutik_i_Grundtræk.pdf, also
%        exists — same 220-pp Google scan).
%  TYPE: Fraktur (Danish body); Latin phrases in antiqua -> \textit{}.
%        Letterspacing (Sperrsatz) -> \emph{}; image-verify per page by zoom.
%
%  OFFSET (VERIFIED): PDF = printed + 8.
%        printed p.4 = PDF 12; p.5 = PDF 13; ... p.8 = PDF 16; p.200 = PDF 208;
%        p.204 = PDF 212. Title page = PDF 9 (printed p.1); Indledning opens on
%        printed p.3 (PDF 11; page number suppressed on the opening leaf). The
%        Indhold lists the Indledning as "1--5" (loose); the running text is 3--5.
%
%  EXTENT: body printed pp. 3--204 (PDF 11--212). Back matter in the scan:
%        Bacon motto (PDF 213), "Rettelser" errata (PDF 214), "Indhold" (PDF 215).
%
%  STRUCTURE (from the Indhold, PDF 215):
%        Indledning — I. Erkjendelseslære (A. Sensualisme, B. Idealisme,
%        C. Sagerkjendelse) — II. Videnskabslære (A. Videnskabernes Videnskab,
%        B. De særskilte Videnskaber, C. Philosophien og de særskilte Videnskaber)
%        — III. Ideelære (A. Ideen i dens Væsenhed, B. Virkeligheden mod Ideen,
%        C. Ideens Gyldighed) — Slutning. Fine structure: numbered §§ ("§ 1." …).
%
%  ERRATA (from the "Rettelser" page; apply inline when transcribing those pages):
%        p.12  l.18 fra oven : Gjenstandenes  -> Gjenstandens
%        p.15  l.9  fra neden: Continuitet    -> Contiguitet
%        p.16  l.1  fra oven : Continuitet    -> Contiguitet
%        p.42  l.14 fra neden: det vilkaarlige-> det virkelige
%        p.48  l.16 fra neden: Punkt;         -> Punkt;"
%        p.48  l.15 fra neden: opløse."       -> opløse
%        p.106 l.14 fra neden: Ilter          -> Metalilter
%
%  STATUS: IN PROGRESS. Done & image-verified: front matter (title + Bacon motto)
%        + Indledning (pp.3--5) + I. Erkjendelseslære §1 (pp.5--7) + A. Sensualisme
%        §2 opening (pp.7--8). Resume at printed p.9 (PDF 17). See RESUME-NOTES.md.
% ============================================================
\documentclass[12pt]{book}
\usepackage[utf8]{inputenc}
\usepackage[T1]{fontenc}
\usepackage[danish]{babel}
\usepackage[protrusion=true,expansion=true]{microtype}
\usepackage[margin=1.4in]{geometry}
\usepackage{setspace}
\usepackage{amsmath}
\usepackage{libertinus}
\usepackage{libertinust1math}
\usepackage{textalpha}   % polytonic Greek typed directly, if it occurs later
\usepackage{fancyhdr}
\usepackage[colorlinks=true,linkcolor=black,urlcolor=black,
  pdftitle={Philosophisk Propædeutik i Grundtræk},
  pdfauthor={Rasmus Nielsen}]{hyperref}
\setlength{\headheight}{15pt}
\pagestyle{fancy}
\fancyhf{}
\fancyhead[LE,RO]{\thepage}
\fancyhead[RE]{\textit{Philosophisk Propædeutik}}
\fancyhead[LO]{\textit{\leftmark}}
\setlength{\parindent}{1.5em}
\setlength{\parskip}{0pt}
\onehalfspacing

\begin{document}

% ============================================================
% TITLE PAGE — PDF 9 (printed p.1)
% ============================================================
\begin{titlepage}
\centering
\vspace*{2.2cm}
{\Huge\bfseries Philosophisk Propædeutik\par}
\vspace{1.1cm}
{\LARGE i\par}
\vspace{0.9cm}
{\LARGE\bfseries Grundtræk\par}
\vspace{1.1cm}
{\large af\par}
\vspace{0.9cm}
{\LARGE\bfseries R.\ Nielsen,\par}
\vspace{0.2cm}
{\normalsize Professor i Philosophien.\par}
\vspace{2.2cm}
\rule{4cm}{0.4pt}\par
\vspace{1.6cm}
{\large\bfseries Kjøbenhavn.\par}
\vspace{0.4cm}
{\normalsize Forlagt af den Gyldendalske Boghandling (F.\ Hegel).\par}
\vspace{0.3cm}
{\small\emph{Thieles Bogtrykkeri.}\par}
\vspace{0.6cm}
{\large 1857.\par}
\end{titlepage}

\frontmatter

% MOTTO — in the Google scan this leaf sits at the BACK (PDF 213), but a book
% motto conventionally faces the text opening; reproduced here at the front.
\thispagestyle{empty}
\vspace*{0.38\textheight}
\begin{quote}
\itshape\noindent Adamant homines scientias et contemplationes particulares, aut
quia autores et inventores se earum credunt, aut quia plurimum in illis operæ
posuerunt, iisque maxime assueverunt. Hujusmodi vero homines, si ad philosophiam
et contemplationes universales se contulerint, illas ex prioribus phantasiis
detorquent et corrumpunt.
\par\vspace{8pt}\normalfont\hfill\textsc{Franc.\ Baco.}
\end{quote}

\tableofcontents
\mainmatter

% ============================================================
% BODY BEGINS — printed p.3 (PDF 11)
% ============================================================

% ---- printed p.3 (PDF 11): Indledning ----
\begin{center}
  {\Large\bfseries Indledning.}\par\vspace{4pt}
  \rule{2.5cm}{0.4pt}
\end{center}
\phantomsection
\addcontentsline{toc}{section}{Indledning}
\markboth{Indledning}{Indledning}

\noindent Den philosophiske Propædeutik lægger an paa at besvare to
Grundspørgsmaal: \emph{hvad er Philosophie?} \emph{hvilken Gyldighed kan der
tillægges Philosophien som Videnskab?}

Til Besvarelse af disse Spørgsmaal maa Propædeutiken charakterisere Philosophien
ved at paavise de Skikkelser, denne paa sine forskjellige Udviklingstrin har
antaget; det almindelige Grundbegreb vil da blive nærmere bestemt, efterhaanden
som det gjennemløber betegnende Stadier. Propædeutiken maa, idet den paa sin Viis
lader Philosophien blive til for Betragtningen, overalt optage et historisk Givet.
Men, endskjøndt det Historiske afbenyttes, kan her dog hverken gives et Udtog af
Philosophiens Historie i dens Heelhed, eller en særskilt udførlig Behandling af
enkelte vilkaarlig valgte Hovedpartier for sig; det Historiske maa opfattes og
benyttes saaledes, at det i den ved Sagens Natur bestemte Sammenhæng belyser de
mødende Opgaver, og tjener Totaliteten.

For at charakterisere Philosophien som Videnskab i dens Forhold til de øvrige
Videnskaber, er det vistnok nødvendigt at gaae nærmere ind paa en Undersøgelse af
det, der medfører de forskjellige Videnskabers Eiendommelighed, og saaledes
paavise det Delingsprincip, hvorpaa deres Afstillelse
% ---- printed p.4 (PDF 12) ----
væsentlig beroer; men da Bestemmelsen af det Almeenvidenskabelige er og bliver
Hovedsagen, kan Behandlingen af de forskjelligartede Videnskabers indbyrdes
Forhold ikke i den Grad optage Forgreningens Enkeltheder, at her skulde kunne være
Tale om en Videnskabernes Encyklopædie.

Forsaavidt Philosophien selv er en i bestemte Led articuleret Videnskab, vil det
for Grundbegrebets Skyld ligeledes være nødvendigt, indenfor denne Videnskabs eget
Omraade at paavise de større Afdelingers indre Sammenhæng med det articulerende
Princip; men, da enhver Oversigt, hvor ikke de enkelte Dele tillige behandles med
nogen Udførlighed, nødvendigviis bliver overfladisk, maa Propædeutiken, istedetfor
at sætte sig en encyklopædisk Fremstilling af Philosophiens Discipliner til
Opgave, meget mere anlægge sine Hovedpartier saaledes, at den ved udførligere
Behandling af hver Afdeling for sig kan faae Leilighed til at belyse den
systematiske Bygning fra en eller anden bestemt Side, og derved efterhaanden aabne
et Indblik i det Heles architektoniske Forhold.

For at indbefatte alle særskilte Hensyn under et omfattende Synspunkt, begynder
Propædeutiken med den Begrebsbestemmelse, at Philosophien er en videnskabelig
Bevidsthedslære.

Bevidsthedslæren viser, hvad det vil sige at blive sig Noget bevidst, og maa
forsaavidt skjelne imellem Bevidstheden selv og det, der kommer til Bevidsthed.
Denne Adskillelse fordrer igjen de adskilte Siders Forening. En tom Bevidsthed, en
Bevidsthed, der Intet er sig bevidst, er en Modsigelse: Bevidstheden maa have et
Indhold.

Uagtet Bevidstheden og dens Indhold forsaavidt ere væsentlig det Samme, som
Bevidsthedens Liden og Virken altid bestemmes i Lighed med det, der under denne
Liden og Virken optages i Bevidstheden, ere de dog i Virkeligheden saa
forskjellige, at den videnskabelige Retning afhænger af, hvorvidt Opmærksomheden
særlig henvendes enten paa
% ---- printed p.5 (PDF 13) ----
den ved Indholdet bestemte Bevidsthed, eller paa Indholdets mod Bevidstheden
stillede Gjenstændighed.

Da Bevidsthedssiden overalt er og maa være det i Philosophien Overveiende, er det
nødvendigt at undersøge Grundforholdet a) imellem Bevidstheden og dens Indhold i
Almindelighed, b) imellem Bevidstheden og det særskilte Indhold, c) imellem
Bevidstheden og det absolute Indhold. Propædeutikens hertil svarende Hoveddele ere
da I) Erkjendelseslære, II) Videnskabslære og III) Ideelære.

% ---- printed p.5 (PDF 13): I. Erkjendelseslære ----
\begin{center}
  {\large\bfseries I.}\par\vspace{4pt}
  {\Large\bfseries Erkjendelseslære.}\par\vspace{4pt}
  \rule{2.5cm}{0.4pt}
\end{center}
\phantomsection
\addcontentsline{toc}{section}{\texorpdfstring{I.\quad Erkjendelseslære}{I. Erkjendelseslære}}
\markboth{Erkjendelseslære}{Erkjendelseslære}

\begin{center}{\bfseries \S\,1.}\end{center}

\noindent Erkjendelsen forudsætter Identitet af Bevidsthed og Indhold; som
Gjenstanden er, maa ogsaa Bevidstheden være: „det Lige forstaaes af det Lige.“
Forsaavidt Bevidstheden er sandsende, er Gjenstanden sandselig, forsaavidt den er
tænkende, tænkelig.

Den væsentlige Forskjel mellem Sandsning og Tænkning kan kjendes paa Forskjellen
mellem Positivt og Negativt. Sandsningens Gjenstand er, uden Hensyn til, om
Sandsen er udvortes eller indvortes, altid et umiddelbart, i Øieblikket Værende,
og forsaavidt et Positivt: det Negative kan ikke sandses. Ved at ophæve
Umiddelbarheden og sætte Negationen fremkommer Tanken; man kan tænke over det
Umiddelbare, men dets Umiddelbarhed kan ikke umiddelbart tænkes. Forsaavidt Tanken
begynder med det Umiddelbare, viser Forstanden sig som den intellectuelle Sands.
Den sandsende Forstand er ikke en Tænkning, men Tænkningens umiddelbare
Forudsætning, en Tænkning i Mulighed: den actuelle Tænkning begynder med det
Negative.
% ---- printed p.6 (PDF 14) ----
I sin rene, qvalitative Modsætning til Sandsningens Umiddelbarhed er Tænkningens
Negation selv umiddelbar. Det umiddelbart Positive og det umiddelbart Negative
begrændse hinanden; ved denne Begrændsning bestemme de hinanden gjensidig. Naar
der spørges, hvad det Positive er, maa der svares ved Hjælp af et Negativt,
\textit{omnis determinatio est negatio}; ligeledes maa paa den anden Side det
Positive forudsættes, naar Negationen og det Negative skal bestemmes.

Som sandsende er Bevidstheden væsentlig lidende; her gjælder det, at „al vor
Erkjenden er en Finden, Modtagen og Skuen (\textit{scimus, quia accepimus})“; som
tænkende er Bevidstheden væsentlig selvvirksom. Hvad Selvvirksomheden angaaer, er
al vor Erkjenden (Jfr. Sibbern: Om Erkjendelse og Granskning) „en Gjøren og
Frembringen (\textit{scimus, quia facimus})“. Det Sandselige objectiveres i
Forestilling og Billede, det Tænkelige i Indseen og Begriben.

Sandsningen er ikke blot ueensartet med Tænkningen: hvad der kan sandses, kan ikke
tænkes, og omvendt: hvad der kan tænkes, kan ikke sandses; men de enkelte Sandser
ere tillige indbyrdes ueensartede: hvad der kan sees, kan ikke høres, hvad der kan
høres, kan ikke føles. Ikke desto mindre bestemme disse Modsatte dog ogsaa
hinanden gjensidig. Det Hørlige kan bestemmes ved det Synlige (Klangfigurer), det
Sandselige ved det Tænkelige. I den sandsende Bevidstheds Liden er der en Virken,
i den tænkende Virksomhed er der en Liden: Forestillingen faaes kun, idet
Bevidstheden selv danner den; Tænkningen bestemmer kun, forsaavidt den ledes,
tvinges og bestemmes ved det, der tænkes; den sandsende Bevidsthed er tænkende,
den tænkende sandsende; det handlende Jeg er lidende, det lidende handlende
(\textit{scientia et potentia in idem concidunt}).

Da det nu er een og samme Bevidsthed, der baade er lidende og virkende, baade
sandsende og tænkende, og da een og samme Gjenstand kan have baade sandselige og
% ---- printed p.7 (PDF 15) ----
tænkelige Egenskaber; saa følger heraf, at der ikke blot maa være et uendeligt
Vexelforhold mellem Sandsning og Tænkning, men at dette Vexelforhold endog maa
kunne stille sig saa forskjelligt, at det snart er Tænkningen, der gaaer op i
Sandsningen, snart omvendt Sandsningen, der forsvinder i Tænkningen.

Forsaavidt Tænkningen gaaer saaledes op i Sandsningen, at kun det i Gjenstanden
Sandselige tillægges Gyldighed, faae vi en eensidig Sandselighedsretning:
\emph{Sensualisme}; den modsvarende Eensidighed bestaaer da deri, at Sandsningen
opfattes som Tænkningens flygtige Forudsætning, saa at kun det i Gjenstanden
Tænkelige, det negative Væsen, det rene Begreb anerkjendes som det Sande:
\emph{Idealisme}. Det uendelige Vexelforhold af Sandsning og Tænkning vil altsaa
først blive indlysende, naar Sensualisme og Idealisme have ydet hver sine Bidrag
til en nærmere Bestemmelse af Sandheden i en dobbeltsidig \emph{Sagerkjendelse}.

\begin{center}
  {\large\bfseries A.}\par\vspace{2pt}
  {\bfseries Sensualisme.}\par\vspace{4pt}
  \rule{2cm}{0.4pt}
\end{center}
\phantomsection
\addcontentsline{toc}{subsection}{\texorpdfstring{A.\quad Sensualisme}{A. Sensualisme}}

\begin{center}{\bfseries \S\,2.}\end{center}

\noindent Med en umiddelbar Eenhed af Sandsning og Tænkning begynder Individets
Bevidsthed; men Erkjendelsens Fremskridt er betinget af, at Tænkningen fremhæves
for sig i dens Modsætning til Sandsningen. Saalænge nu Ligevægten mellem de to
Bevidsthedsyttringer bevares, saalænge Overeensstemmelsen mellem det Sandselige og
det Intellectuelle endnu vedligeholdes, er Individualiteten lykkelig, og Dannelsen
skjøn. Men med den voxende Abstraction, idet Indbildningskraften ret begynder at
bearbeide og ombanne
% ---- printed p.8 (PDF 16) ----
Sandsningselementerne, bliver Tanken urolig, en Spænding indtræder, og Ligevægten
gaaer tabt. Den tabte Ligevægt, hvis Kjendemærke er, at snart det Intellectuelle,
snart det Sandselige, snart en overspændt Aandelighed, snart en djærv Naturlighed
har faaet Overhaand, vil da igjen oprettes, og den saaledes oprettede Ligevægt
igjen tabes.

Tilsvarende Overgange sees i Philosophiens Historie.

Den græske Philosophie begynder hos Ionerne (Thales, Anaximander, Anaximenes) med
en naiv Erkjendelse af det Usandselige i det Sandselige; Tanken søger
Tilværelsens Princip, aner ogsaa, at Principet maa være „et Ubegrændset“, men
finder det dog nærmest repræsenteret ved et givet sandseligt Stof: „Vandet“,
„Luften“. Denne Tænkningens Gaaen i Eet med Sandseligheden er endnu ikke
Sensualisme. Tanken er her saa langt fra at føle Trykket af Sandsernes Overvægt,
at den netop er i Begreb med at hæve sig over det Sandseliges Udstykning; det er
Aleenheden, Tanken her har grebet med ungdommeligt Mod, idet den første Gang
prøver et Opsving mod det Almene.

Først, naar Tænkningen har forladt Umiddelbarheden, har fortæret Indholdet, og er
mat af at speile sig selv i sig selv, fremkommer Sensualismen. Sensualismen er en
bevidst Modsætning, en udtrykkelig Modvægt mod det tomme Begrebsvæsen: først
efterat den græske Tænkning (Pythagoræer, Eleater, Plato) havde forflygtiget
Sandsningsstoffet, fastholdt den rene Væsenhed, og forsøgt sig i allehaande
Abstractioner, hævder Aristoteles Empirien, og den tilbagetrængte Sensualisme
kommer saa egentlig frem i Stoikernes og Epikuræernes Erkjendelseslære.

De Stadier, der i denne Henseende charakterisere Bevægelsens Retninger i den
græske Philosophie, gjentage sig efter en endnu større Maalestok i Philosophiens
Universalhistorie.

I Modsætning til Middelalderen fremstiller Oldtiden i det Hele Sandsningens og
Tænkningens skjønne Ligevægt og naturlige Eenhed; men i Middelalderen gjør den
abstracte

% ---- printed p.9 (PDF 17) ----
Aandelighed et saadant Overgreb, at en derved fremkaldt Reaction har kunnet føre
den nyere Tid til en Aandsfornegtelse, hvis Omfang, Skarphed og indre
Gjennemsigtighed er ukjendt hos de Gamle.

I Middelalderen, da det Naturlige kun er et Forkrænkeligt, det Sandselige kun et
Syndigt, det Nærværende kun en Skuffelse, er det naturligviis umuligt at erkjende
Tingene, som de ere. Ængstet af Phantasie og Følelse gaaer Videnskabens Tænkning i
Troens Tjeneste, og kues af det Overnaturlige; den forkuede Tænkning finder kun en
formel Frihed i det Logiske. Herfra den Tomhed, Spidsfindighed, det Ordkløveri og
den videnskabelige Ufrugtbarhed, hvori Scholastiken ender. Istedetfor at gribe en
Virkelighed, og ernære sig ved Bearbeidelsen af et fordøieligt Stof, tærede
Reflexionen paa sig selv, og piintes af Hunger.

De Bestræbelser, den videnskabelige Hunger fremkaldte, droge Aanderne til Naturen
i to, hinanden modsatte Retninger, en speculativ og en empirisk; men dog saaledes,
at Noget af den ene kan gjenfindes i den anden.

\emph{Den speculative Retning.} Forsaavidt der kan være Tale om en speculativ
Sensualisme, finde vi en saadan hos de italienske Philosopher: Telesius, Cardanus,
Bruno o. Fl. Med afgjort Forkjærlighed for den antike Philosophies Former
begeistredes disse Tænkere ved Forvisningen om Naturens Identitet med Aanden, ved
Beskuelsen af det Guddomsliv, der besjæler Verdensorganismen, den levende Eenhed,
der gjennemvirker Alt.

At aristoteliske Kategorier anvendes til Udtydning af Naturphænomenerne, er i og
for sig ikke Sensualisme; det Samme gjorde jo Scholastikerne i god Forstaaelse med
Hierarchiet. Det Sensualistiske kommer derimod tilsyne i den Kjækhed, Freidighed
og Ubetingethed, hvormed et allestedsnærværende Naturliv fortrænger det Hiinsidige,
og sætter sig i det Overnaturliges Sted. Det Fornuftens Rige, hvorom Giordano Bruno
handler (\textit{De la causa, principio ed uno}),
% ---- printed p.10 (PDF 18) ----
er tillige et Sandsernes Rige, og de Vidundere, hvortil Vanini (\textit{De
admirandis Naturæ, reginæ Deæque mortalium, arcanis}) sigter, skyldes ikke „den
evige Konge“, men „de Dødeliges Dronning“.

\emph{Den empiriske Vei.} Ved ædruelig Forskning, ved omhyggelig Iagttagelse af
Phænomenerne (Tycho Brahe) og tilsvarende Beregninger blev den Vei endelig
betegnet, som ene kunde føre til sand Naturerkjendelse. Faldloven paa Jorden
(Galilei) og Planetbanerne paa Himlen (Keppler) viste tydelig nok, at Veien var
funden. De Forskere, der saaledes forstode, ikke blot at aabne Øiet, men ogsaa at
væbne Øiet (Kikkerten), gjorde da i Videnskabens Navn Sandsernes Ret gjældende, og
satte Kjendsgjerninger i Kraft, hvis Uimodsigelighed omsider maatte betvinge
Middelalderens Indbildninger, og ryste den scholastiske Metaphysik i dens
Grundvolde.

Skulde slige Bestræbelser opfattes, bedømmes og fremstilles i en paa
Naturvidenskabens Tarv beregnet Erkjendelseslære, da maatte det vistnok især komme
an paa at indskjærpe to Hovedfordringer: Tankens Frigjørelse fra abstracte Theorier
(fra \textit{anticipationes naturæ}), og Videns Tilbageførelse til den sande
Erfaring (\textit{interpretatio naturæ}). Ved at modarbeide Formalismen, fremhæve
„Inductionen“, og bringe Naturlovens Begreb i Samklang med Physikens Fordring har
F. Baco af Verulam († 1626; \textit{Novum organum scientiarum}), skjøndt uden
naturvidenskabelig Indvielse, dog i aandrige Træk paaviist et nyt Grundlag, og
betegnet Overgangen til en ny Tid.

\begin{center}{\bfseries \S\,3.}\end{center}

Den fremrykkende Empirisme opdagede snart den Mistydning af det Aprioriske, der
skjulte sig i Læren om „de medfødte Ideer“. Imod denne Mistydning har J. Locke
(† 1704) rettet sit „philosophiske Forsøg over den menneskelige Forstand“
(\textit{An essay concerning human understanding}
% ---- printed p.11 (PDF 19) ----
\textit{in four books}), hvis Indhold vi her kunne betragte under følgende
Synspunkter: „Ideernes“ Oprindelse, „Ideernes“ Væsen og Tingenes Erkjendelse.

\emph{a) „Ideernes“ Oprindelse.} Der gives ingen medfødte Ideer. Selv om det kunde
godtgjøres, at der gaves Sætninger, om hvis Almindelighed Alle maatte være enige,
vilde dette dog Intet bevise, naar disse Sætningers almindelige Antagelse kunde
forklares paa en anden Maade. Men det er saa langt fra, at slige Sætninger findes,
at Nationer paa forskjellige Dannelsestrin endog ere uenige om de allervigtigste
praktiske Grundsætninger. Fremsætter man det Bud: „Hvad Du vil, at Andre skulle
gjøre imod Dig, gjør det ogsaa imod Dem,“ da vil Enhver, som ikke har hørt Budet
før, men tænker efter, naturligviis spørge: hvorfor? Men at spørge om Begrundelsen
af et Bud, der er medfødt, er en Modsigelse.

Hvad de theoretiske Sætninger, saasom: „Alt hvad der er, det er; det er umuligt, at
een og samme Ting kan til samme Tid baade være og ikke være“, angaaer, da ere Børn,
Idioter og Udannede aldeles ubekjendte med deres Betydning. Men, dersom slige
Grundsætninger virkelig vare Forstanden medfødte, saa maatte jo Alle fra den
tidligste Barndom være vidende om dem, thi at være i Forstanden er det Samme som at
være i Viden.

Ligeledes er det en Udflugt, naar man, for at forsvare de medfødte Ideer, beraaber
sig paa, at disse Ideer umiddelbart indlyse for Bevidstheden, saasnart Mennesket
begynder at bruge sin Fornuft. Forholdet er just det omvendte, at det nemlig ikke
er almindelige, men particulære Sætninger, saasom: „Sødt er ikke bittert“, man
allerførst bliver sig bevidst, naar man begynder at tænke. At theoretiske og
praktiske Grundsætninger skulde være medfødte, er i det Hele ligesaa urimeligt, som
at Kunster og Videnskaber skulde være medfødte. Det er altsaa umuligt at fødes med
Ideer;
% ---- printed p.12 (PDF 20) ----
Barnets Sjæl er som et reent Papiir, saa tom, som en \textit{tabula rasa}.

Hvorfra faaer Forstanden da sine Ideer? Fra Erfaringen. Men Erfaringen er af
dobbelt Art: forsaavidt den gjøres ved en gjennem Sandserne anstillet Iagttagelse
af ydre Gjenstande, kalde vi den Fornemmelse, \emph{Sensation}; forsaavidt den
beroer paa Iagttagelse af vor egen Forstandsvirksomhed, kalde vi den indvortes
Sands, \emph{Reflexion}. Fornemmelse og Reflexion ere at betragte som „de eneste
Kanaler, der tilføre Forstanden Kundskaber, de eneste Vinduer, gjennem hvilke
Ideernes Lys falder ind i Sjælens dunkle Værelse.“

\emph{b) „Ideernes“ Væsen.} De fra Erfaringen kommende Ideer ere deels enkelte,
deels sammensatte.

De enkelte Ideer afgive Erkjendelsens Materiale, ere, saa at sige, Bogstaverne i
Tilværelsens Alphabet. Idet disse Enkeltideer udenfra paatrænge sig Bevidstheden,
er Sjælen passiv, og optager Omverdenens Indtryk, som Speilet optager
% ERRATA (Rettelser p.12 l.18 fra oven): print „Gjenstandenes“; read „Gjenstandens“.
Gjenstandens Billede; faaes de derimod ved Reflexionen, saaledes som Tænkningens
Idee, Villiens Idee, er Sjælen activ. Enkeltideerne ere deels saadanne, der komme
til Forstanden gjennem en enkelt Sands: Farverne gjennem Øiet, Tonerne gjennem
Øret, Uigjennemtrængeligheden gjennem Berøringssandsen; deels saadanne, der tilføres
den gjennem flere Sandser, som naar Rummet og Bevægelsen opfattes baade ved
Berøringen og Synet.

Ved de enkelte Ideers Forbindelser og Sammensætninger kan nu Sjælen erkjende
Omverdenens Gjenstande, „ligesom man i Sproget ved mangfoldige Sammensætninger af
Bogstaver kan danne Stavelser og Ord“; men, idet de sammensatte, complexe Ideer
henføres til tre Hovedklasser: Modi, Substanser og Forhold, bliver det især
Opgaven, at gjøre sig Rede for Substansbegrebet.

Det nominelle Væsen angiver Forholdet imellem Slægt og Arter; det reelle Væsen
betegner derimod den tilværende
% ---- printed p.13 (PDF 21) ----
Tings indre Eiendommelighed. I Enkeltideerne og Modi ere det reelle og det
nominelle Væsen Eet og det Samme; med Hensyn paa „Substanserne“ ere de
forskjellige. At en Flade, som begrændses af tre rette Linier, er et Triangel,
udtrykker Eenheden af det reelle og det nominelle Væsen. Anderledes med den Deel af
Materien, der danner Ringen paa min Finger; af de usandselige Deles reelle
Eiendommelighed afhænge Egenskaberne ved dens Farve, Vægt, Blødhed, Haardhed o. s.
v. Efter de iagttagne Egenskaber dannes Benævnelsen Guld, et Udtryk, der betegner
det nominelle, men ikke det reelle Væsen.

Ved Sandsning og Reflexion erfare vi, hvorledes et Antal Enkeltideer træde sammen i
en eiendommelig Forbindelse, forudsætte desaarsag en nødvendig Forbindelsesgrund,
og henføre denne Grund til et for sig bestaaende Substrat, nemlig \emph{Substansen};
men er det nu muligt at erkjende Substansen?

\emph{c) Tingenes Erkjendelse.} Som Bærer af de Qvaliteter, ved hvilke
Enkeltideerne frembringes hos os, kan Substansen kun erkjendes derved, at der gjøres
Forskjel imellem de Egenskaber, der nødvendigviis tilkomme Tingen i sig, og bevise
Enkeltideernes Overeensstemmelse med Gjenstandene (\textit{qualitates primariæ}), og
saadanne Egenskaber: Farve, Tone, Smag, Lugt, Varme, Kulde, som vel hidrøre fra
Tingenes Indvirkning paa Sjælen, men dog ere aldeles forskjellige fra Tingenes
objective Væsen (\textit{qualitates secundariæ}).

Paa denne Maade vilde Locke bestemme Modsætningen mellem den reale (\textit{real
truth}) og den mentale Sandhed (\textit{mental truth}); men uagtet Substansen just
adskiller sig fra alle andre complexe Ideer derved, at den har sin Archetyp uden for
os --- man forbinder ikke et Faars Brægen med Forestillingen om en Hest ---, og just
deri tilkjendegiver sin objective Realitet; er det os dog alligevel umuligt at
trænge ind i Substantialiteten selv.
% ---- printed p.14 (PDF 22) ----
Hvad vi kalde det Væsentlige, gaaer ikke paa Individet. „Det er vistnok nødvendigt,
at jeg er, hvad jeg er: Gud og Naturen har dannet mig saaledes; men jeg veed dog
Intet, som er mig, Individet, væsentligt som Individ: en Hændelse, en Sygdom kan
hidføre utrolige Forandringer.“ Saa gaaer det Væsentlige altsaa paa Arten; men det
kan umuligt være det reelle Væsen, som bestemmer Arten, da det reelle Væsen er os
ubekjendt. Vi ere henviste til en Samling af sandselige Qvaliteter, som vi kunne
iagttage; men selv, naar vore Iagttagelser ere anstillede med den størst mulige
Nøiagtighed, er vor Erkjendelse ligesaa langt fra at svare til den Tingens indre
Constitution, hvoraf disse Qvaliteter følge, „som en Bondes Forestilling om Uhret i
Straßborg er langt fra at svare til denne mærkværdige Maskines indre, kunstfærdige
Bygning.“

Hvad Locke havde begyndt, fortsatte Condillac († 1780; \textit{Essai sur l'origine
des connaissances humaines}). Ved at lade Erkjendelsens tvende Kilder: Sandsningen
og Reflexionen flyde sammen saaledes, at denne forsvinder i hiin, og ved at lægge et
forstærket Eftertryk paa Fornemmelsen, kom Condillac til det Resultat, at den
saakaldte indre Sands kun er en Reflex af det Ydre, og at følgelig alle
Aandsyttringer kun forestille modificerede Sandseindtryk. Har Sensualismen, saaledes
gjennemført, først viist, at det Sande kun er et Sandseligt, da ligger det nær, at
vende Sætningen om, og sige: kun det Sandselige er det Sande. Her er en
Materialisme, for hvilken alt Aandeligt er en Indbildning (La Mettrie).

\begin{center}{\bfseries \S\,4.}\end{center}

Af Forudsætninger, hvorved et eensidigt Standpunkt udtrykkelig betegnes, kunne
Conseqvenser udvikles deels i sideordnet, deels i overordnet Retning. Til
Conseqvenser i en med Standpunktet sideordnet Retning hører Alt, hvad der udfolder
dets Charakteer, og hæver dets Betydning
% ---- printed p.15 (PDF 23) ----
stedse videre Omfang. I denne Forstand have Condillac og La Mettrie drevet
Sensualismens Conseqvenser til det Yderste.

Til Conseqvenser i en for Standpunktet overordnet Retning høre de, der føre ud over
Standpunktet selv, og altsaa enten fastsætte eller dog ibetmindste indlede et nyt
Standpunkt. I denne Forstand er Humes († 1776) Skepticisme (\textit{A treatise of
human nature}) at opfatte som en høiere Conseqvens af den lockeske Sensualisme.

For at prøve Begrebernes Gyldighed, maa man --- heri er Hume enig med Locke ---
holde sig til Erfaringen. De Begreber, som kunne bevise deres Oprindelse fra
Sandseindtryk (\textit{impressions}) eller dog fra de til disse svarende
Reflexioner, bør anerkjendes; de, som ikke kunne det, ere af uægte Herkomst, og bør
forkastes.

Ifølge Locke opstaaer „Aarsagens og Virkningens Idee“, naar vi ved at betragte
Tingenes stabige Vexel iagttage, hvorledes visse Ting, det være sig nu Substanser
eller Qvaliteter, ved et Andets lovbestemte Virksomhed begynde at existere. Men
hvilken Gyldighed har nu denne Iagttagelse? hvoraf vide vi, at to Ting forholde sig
til hinanden som Aarsag til Virkning? \textit{A priori} kunne vi ikke vide det; thi
den aprioriske Erkjendelse er analytisk, men Erkjendelsen af Aarsagsbegrebet maa
være synthetisk, da Virkningen virkelig er noget Andet end Aarsagen. Vi vide det
heller ikke af Erfaring; thi Erfaringen viser os kun, at den ene Kjendsgjerning
deels i Tiden følger efter den anden (Succession),
% ERRATA (Rettelser p.15 l.9 fra neden): print „Continuitet“; read „Contiguitet“.
deels i Rummet berører den anden (Contiguitet); men ingen Erfaring kan lære os, at
det Ene væsentlig følger af det Andet, som Virkningen af sin Aarsag.

For at hævde Substansens Gyldighed havde Locke allerede henviist til Vanen; denne
Anviisning følger Hume: „Vore Slutninger ud af Erfaringen grunde sig alene paa Vane;
vi ere vante til (\textit{the mind is carried by habit}) at see den ene Ting følge
deels med, deels efter den anden, og vænne os da derved til den Forestilling, at det
Ene
% ---- printed p.16 (PDF 24) ----
% ERRATA (Rettelser p.16 l.1 fra oven): print „Continuitets-“; read „Contiguitets-“.
maa være en Virkning af det Andet. Af Contiguitets- og Successionsforholdet danner
vor Indbildning Causalitetsforholdet: \textit{cum hoc aut post hoc, ergo propter
hoc}.“

At Undersøgelsens Traade her løbe sammen i Spørgsmaalet om Causalitetsbegrebet som i
deres Knudepunkt, viser den Sikkerhed, hvormed Opgaven er stillet.

Er Causalitetsbegrebet først faldet, saa falde Kategorierne: Kraft, Substans,
Nødvendighed, Sandhed, som af sig selv; de Væsensbaand, der skulde sammenholde
Tingene, maae nu alle briste.

Som hidrørende fra en Samvirken af Fornemmelse og Reflexion kunde Begrebet „Kraft“
efter Lockes Mening vel have sin Gyldighed; men hvorledes er det nu muligt at
erkjende Kraften? Det kunde ved første Øiekast synes, som om Forholdet mellem Sjæl
og Legeme i denne Henseende afgav en paalidelig Forudsætning. Idet visse
Legemsorganer uvilkaarlig adlyde vor Villie, virker jo Aanden som Kraft, og vi
behøve altsaa kun at blive os vor Vilkaarlighed bevidst, for at blive os en
Kraftyttring bevidst. Dette er imidlertid en skuffende Reflexion; deels kjende vi
ikke de Midler, hvorved Aanden virker paa Legemet; deels er det jo heller ikke alle
Legemets Organer, hvis Virksomhed er Villien underlagt. Skulde Kraften være
Gjenstand for Erfaring, da maatte, hvad der er umuligt, Aarsag og Virkning ogsaa
være Gjenstand for Erfaring.

Ved at fremstille Substansen som et Ubekjendt havde Locke allerede indledet Tvivlen
om dens Tilværelse. Hvorledes skulle vi nu erkjende Substansen? Skulde
Substantialitetens Gyldighed lade sig bevise, da maatte jo den indre nødvendige
Sammenhæng af de Egenskaber, der i Substansen ere forbundne til Eenhed, lade sig
bevise; men ved Causaliteten er den nødvendige Sammenhæng (\textit{this idea of a
necessary connexion}) bortfalden. At henvise til Forskjellen imellem oprindelige og
afledede Egenskaber hjælper ikke. De afledede Egenskaber ere alle subjective, og
have altsaa kun
% ---- printed p.17 (PDF 25) ----
Tilværen i vor Opfatning; hvad de oprindelige angaaer, da lade de sig tilsidst
opløse i een eneste, nemlig „Soliditeten,“ men Soliditeten er et tomt Ord, eftersom
vi, naar alle øvrige Egenskaber ere ophævede, umulig kunne gjøre det Solide til
Gjenstand for nogen Erfaring.

I Tillid til den „indre Sands“ havde Locke endnu fastholdt „Villiens Idee,“ og
anerkjendt Selvet. Men hvorledes skulde det nu være muligt at forvisse sig om Jegets
Realitet? Vore Følelser, Tanker og Forestillinger kunne umulig betragtes som
Kraftyttringer af et Jeg, naar Kategorien Kraftyttring ingen Gyldighed har; de
kunne umulig henføres til et virkende Selv, som til deres sande Aarsag, naar
Aarsagsbegrebet ingen Sandhed har.

Den Rest af Realitet, der endnu hos Locke hæftede ved Sandhedsbegrebet
(\textit{real truth}), er altsaa forsvunden: Sandheden er borte, Sandsynligheden er
traadt i dens Sted. Jeget, Selvet, er da i Virkeligheden ikke Andet end en
Indbildning, som fremkommer derved, at et Indbegreb af mange hurtig paa hinanden
følgende Forestillinger underlægges et Substrat.

At en Philosophie, som underordner Tænkningen under Sandsningen, endog i den Grad,
at Begreberne maae tabe deres Almeengyldighed, at Tingenes Love maae forsvinde for
Ideeassociationens Regler, Fornuftgrunden for Vanen, Sandheden for Sandsynligheden,
Jegets Selverkjendelse for en bedaarende Indbildning, i sig hverken er holdbar eller
tilfredsstillende, forstaaer sig; ikke desto mindre er Humes Skepticisme i
Modsætning til den dogmatiske Materialisme en høiere Potens, en aandig Magt, hvis
Angreb paa Tanken kan forvandle sig til et Tankens Forsvar.

\begin{center}\rule{2.5cm}{0.4pt}\end{center}
% ---- printed p.18 (PDF 26) ----
\begin{center}
  {\large\bfseries B.}\par\vspace{2pt}
  {\bfseries Idealisme.}\par\vspace{4pt}
  \rule{2cm}{0.4pt}
\end{center}
\phantomsection
\addcontentsline{toc}{subsection}{\texorpdfstring{B.\quad Idealisme}{B. Idealisme}}

\begin{center}{\bfseries \S\,5.}\end{center}

\noindent Kan Sensualismen kjendes derpaa, at den afhører Sandserne for at dømme
Tanken, saa kan Idealismen igjen kjendes derpaa, at den gjør Tanken til Sandsernes
Dommer; sætter Sensualismen Tingene over Bevidstheden, saa vender Idealismen
Forholdet om, og sætter Bevidstheden over Tingene.

Idet Cartesius († 1650; \textit{Principia philosophiæ}), begynder med Tvivlen, vil
han ligesom Baco afkaste Fordomme, og frigjøre Sindet fra blind Autoritetstroe;
men, idet han ved at tvivle om Alt (\textit{de omnibus dubitandum est}) ogsaa kan
drage Naturtingenes og Erfaringens Gyldighed i Tvivl, gaaer han videre end Baco.

Da Tvivlen er en Tankens Gjerning, er Tanken midt i Tvivlen hævet over Tvivlen; der
kan tvivles om Alt, men den tvivlende Tankes Gyldighed kan ikke drages i Tvivl; at
tvivle om Tvivlen er at stadfæste Tvivlen.

Empirismen, der begynder med en ubetinget Stolen paa det umiddelbare Sandseindtryk,
ender conseqvent i Skepticisme; thi Skepticisme er en Aandsvirksomhed, hvori Tanken,
efterat have stirret sig blind paa det Sandselige, uophørlig kommer fremmed til sig
selv, nedriver sig selv, og ikke kan kjende sit eget Væsen.

I den idealistiske Tvivlen om Alt begynder Tanken just med at finde sig selv, idet
den begynder med at fjerne det Fremmede, og afryste Overleveringens Paahæng. Denne
Tvivlen er da saa langt fra at tabe sig i Skepticisme, at den Tænkende meget mere
vil forvisse sig om Gyldigheden af Alt, hvad der udholder Tænkningens Prøvelse.

% ---- printed p.19 (PDF 27) ----
Den Tænkende indseer da ikke allene sin Tankens Eenhed med Væren (\textit{cogito
ergo sum}); men faaer tillige Vished om, at Bevidstheden bæres, gjennemtrænges og
opfyldes af en absolut Væren, af det Ubetingede, hvis Billede er i Sjælen, hvis
Former (\textit{ideæ innatæ}) indlyse med en Klarhed, som kun kan udgaae fra det
Fuldkomne.

Hos Cartesius er Idealismen dog endnu indskrænket derved, at der i Modsætning til
den Væren, som er i Tanken, staaer en materiel Væren ligeoverfor Tanken; at hæve
denne substantielle Adskillelse, og forvandle Tænken og Væren til Attributer i een
og samme Substans er et Fremskridt, der skyldes Spinoza († 1677). Forsaavidt her
spørges om Tankens Gyldighed og dens Identitet med Væren, er Spinozismen afgjort
Idealisme.

Hvad Sjælens aprioriske Selvvirksomhed angaaer, gik Leibnitz († 1716) et Skridt
videre. I sine „Nouveaux essais“ modarbeider han den lockeske Sensualisme, og
befrier Antagelsen af „de medfødte Ideer“ fra den Misforstaaelse, som kunde
retfærdiggjøre Lockes Indvendinger. Imod Sensualisternes: \textit{nihil est in
intellectu, quod non antea fuerit in sensu}, tilføies da et: \textit{nisi
intellectus ipse}. At „Ideerne“ ere medfødte, maa ikke mistydes, som om de ved
Fødselen laae fuldfærdige i Sjælen, da de tvertimod kun ere i den efter Anlæget.
Denne Synsmaade er i væsentlig Sammenhæng med Monadelæren; som Monade er Sjælen
oprindelig selvvirksom, kan ingen Indtryk modtage udenfra, men derimod udvikle en
Verden i sit Indre.

Medens Locke sammenlignede Sjælen med en ubeskreven Tavle, sammenligner Leibnitz
den med „en Marmorblok, hvori Aarerne præformere Billedstøttens Skikkelse.“
Modsigelsens Grundsætning (\textit{principium contradictionis}), Sætningen om den
tilstrækkelige Grund (\textit{principium rationis sufficientis}) og det saakaldte
\textit{principium indiscernibilium} høre blandt „de medfødte Ideer.“
% ---- printed p.20 (PDF 28) ----
Hvorledes Idealismen ved at overtænke Forholdet mellem Sandserne og det Sandsede,
og paavise de Modsigelser, der overalt ledsage Erkjendelsen af et Udvortes, kan
forvandle den hele Yderverden til en Aandeverden, er i Særdeleshed viist af G.
Berkeley († 1754; \textit{A treatise concerning the principles of human
knowledge}; og \textit{Three dialogues between Hylas and Philonous}). ---
Overbeviist om, at det er Antagelsen af Materiens ubetingede Existens, der især
begunstiger Skepticisme og Atheisme, henregner Berkeley denne Antagelse til det
meest Meningsløse og Forvirrede, der nogensinde kan opkomme i et Menneskes Hoved.

Sandselige Ting ere saadanne, som umiddelbart (\textit{immediately}) opfattes ved
Sandserne, hvorimod Tingenes Aarsag, der ikke kan sandses, udtrykkelig maa tænkes,
og, idet den tænkes, være i Bevidstheden. Seer jeg, at den ene Deel af Himlen er
rød, den anden blaa, saa slutter jeg, at der maa være en Aarsag til denne
Farveforskjel; men Aarsagen kan ikke siges at være en sandselig Ting (\textit{or
perceived by the Sense of Seeing}). Heller ikke kan jeg, naar jeg hører
forskjellige Toner, siges at høre Tonernes Aarsag; ligesaa lidt som jeg, naar jeg
ved Følelsen erfarer, at en Ting er hed eller tung, kan siges at føle Hedens eller
Tyngdens Aarsag.

Sandselige Ting ere Indbegreb af sandselige Egenskaber (\textit{Combination of
sensible Qualities}); disse sandselige Egenskaber ere ikke udenfor Sjælen, men i
Sjælen. Saalidt som man kan tillægge Naalen en Fornemmelse af den Smerte, dens Stik
foraarsager os; saalidt kan man tillægge Ilden Varme. At have en høi Varmegrad, at
være hed, er at være i Smerte; skulde Ilden nu virkelig være hed, da maatte den
brænde sig selv, som den brænder os, den maatte da være i Smerte; men det er
urimeligt.

Ikke blot om de afledede Egenskaber, ogsaa om de oprindelige gjælder det, at de kun
ere til i vor Bevidsthed. Skulde Egenskaber som: Udstrækning, Figur, Tyngde,
Tæthed,
% ---- printed p.21 (PDF 29) ----
Bevægelse, være Noget udenfor Bevidstheden, da maatte en og samme Ting være baade
stor (naar den nemlig sees nærved) og lille (forsaavidt den sees i lang Frastand),
baade glat og knubret, baade lige og kroget, og den samme Bevægelse baade langsom
og hurtig; men det er urimeligt.

Ligesaa urimelig er den Mening, at de nævnte Egenskaber skulde have et Substrat;
det Substrat, hvoraf Udstrækningen og de øvrige Egenskaber skulde afhænge, maatte
jo, for at bevirke Udstrækning, allerede forud have en Udstrækning, og saaledes i
det Uendelige.

Materiens uendelige Deelbarhed er almindelig antaget. Der er da i hver endelig
Materie et uendeligt Antal Dele, som Sandserne endnu ikke opfatte; men jo mere Øiet
skjærpes, desto flere Dele seer det, og desto større bliver hver Deel: kunde Øiet
skjærpes i det Uendelige, vilde hver mindste Deel blive uendelig stor. Denne
Materialisme modsiger sig selv; thi heraf følger, at ethvert Legeme er i sig
uendelig udstrakt, det vil sige: uden al Skikkelse eller Figur. Det er altsaa ikke
et Materiens Væsen, men Bevidstheden selv, der danner alle Forandringer i de
Legemer, hvoraf den synlige Verden (\textit{the visible World}) bestaaer: Materien
er en Idee.

Hvad den mathematiske Naturphilosophie angaaer, da glemmer Mathematiken at gjøre
tilstrækkeligt Rede for sine Principer. Arithmetiken regner med Enere
(\textit{Unities}), men er ikke istand til at vise, hvad „Eneren“ egentlig er; i
Geometrien spiller den endelige Udstræknings uendelige Deelbarhed en Hovedrolle;
men den endelige Udstrækning er, som allerede viist, kun en Idee i Sjælen. Da det
nu er indlysende, at Ingen er istand til at adskille uendelig mange Dele enten i en
endelig Linie eller i en endelig Flade eller i et endeligt Legeme, saa kan Enhver,
ved at tydeliggjøre sig Udstrækningens Idee, selv overbevise sig om, at den
foregivne Uendelighed ikke er tilstede.
% ---- printed p.22 (PDF 30) ----
Hvad vi i Hverdagssproget kalde Ting, maae vi da, ifølge de givne Oplysninger,
kalde Ideer. Vistnok er det en haard Tale, om vi til dagligt Brug skulde æde Ideer,
drikke Ideer, og klæde os i Ideer; men Eet er Sprogbrugen, et Andet Sandheden. At
Tingene udenfor os forvandles til Ideer i vor egen Sjæl, er jo ikke at forstaae,
som om de sandselige Forestillinger beroede paa vor Vilkaarlighed. Vi kunne
vilkaarlig aabne og lukke Øinene, men de Skikkelser og Farver, der vise sig for
vore opladte Øine, paavirke os med Nødvendighed.

Vore sandselige Fornemmelser og Forestillinger maae vistnok henføres til en Aarsag,
der er uafhængig af os. Men, da Aarsagen og Virkningen maae være ligeartede, og da
en Ting umulig kan meddele os, hvad den ikke selv har; saa er det ligesaa urimeligt
at antage, at en Ting, der ikke selv har Fornemmelse, skulde kunne vække
Fornemmelse, som, at et Hjerneindtryk skulde kunne blive til en Forestilling. Der
existerer altsaa ingen Legemer, men kun Aander; det er ikke en materiel Aarsag, men
en almægtig Aand, der igjennem alle Sandsebilleder, Forestillinger, Fornemmelser og
Indtryk lader de endelige Aander meddele sig til hverandre.

Ved saaledes at forklare det Sandselige og afskaffe det Materielle har denne
Idealisme forvandlet sig til en Spiritualisme, der tilintetgjør Naturen, appellerer
til det Overnaturlige, og culminerer i Theologie. \textit{It is therefore plain,
that nothing can be more evident to any one, that is capable of the least
Reflexion, than the Existence of GOD, or a Spirit who is intimately present to our
Minds, producing in them all that variety of Ideas or Sensations, which continually
affect us, on whom we have an absolute and intire Dependence, in short, in whom we
live and move, and have our Being.} (\textit{Principles of hum. knowl. London 1734.
S. 166}).
% ---- printed p.23 (PDF 31) ----
\begin{center}{\bfseries \S\,6.}\end{center}

Iblandt de mange hinanden krydsende Tilløb til en Erkjendelseslære mødes
Skepticismen og Dogmatismen som stridige Resultater. I Skepticismen (Hume) er
Tanken splidagtig med sig selv, og Erfaringen staaer som et uforklaret Factum; i
Dogmatismen (Chr. Wolf) er Værens Eenhed med Tænkningen opfattet saa umiddelbart,
at det stive Forstandsschema vel blænder ved sine Beviser, men mangler
Begrundelsen. Imod disse mægtige Eensidigheder optræder I. Kant († 1804;
\textit{Kritik der reinen Vernunft}).

Cartesius's Idealisme, der gjør Gjenstandenes Tilværelse i Rummet ubeviislig og
tvivlsom, kalder Kant den problematiske; Berkeley's, der fremstiller Tingenes
udvortes Tilværelse som falsk og umulig, kalder han den dogmatiske. Den dogmatiske
Idealisme er uundgaaelig, naar Rummet skal ansees for en Egenskab, der tilkommer
Tingene i og for sig; thi da er enhver Ting, for hvilken et saadant Rum skulde
afgive Betingelse, en Uting. Grunden til denne Idealisme ophæves i „den
transcendentale Æsthetik.“ Den problematiske Idealisme bliver derimod først
gjendreven, idet der føres Beviis for, at selv vor indre, af Cartesius ikke
betvivlede Erfaring, kun er mulig under Forudsætning af ydre Erfaring.

Endskjøndt den kantiske Fornuftkritik, som et i sig systematisk Hele, ingenlunde er
bleven til af endeligt Hensyn til denne eller hiin Synsmaade, kan den dog ikke
desto mindre, naar man saa vil, opfattes saaledes, at „den transcendentale
Elementærlæres“ tvende Dele, nemlig Æsthetiken og Logiken, fornemmelig maae siges
at være rettede imod Skepticismen, dens tredie Deel, nemlig Dialektiken, imod
Dogmatismen, og „den transcendentale Methodelære“ imod begge. Vi behandle da i
denne Sammenhæng den kantiske Erkjendelseslære, forsaavidt den ved sin Modsætning
til de nævnte Yderligheder charakteriserer sig som en „\emph{transcendental}“
Idealisme. Altsaa:
% ---- printed p.24 (PDF 32) ----
\begin{center}{\bfseries 1)\quad Mod Skepticismen.}\end{center}

At al vor Erkjendelse begynder med Erfaring, er udenfor al Tvivl; men deraf følger
endnu ikke, at Erkjendelsen i Eet og Alt har sit Udspring fra Erfaringen. Der gives
aprioriske Sandseanskuelser, aprioriske Forstandskategorier og aprioriske
Fornuftideer.

a) „Den transcendentale Æsthetik“ omhandler de rene Sandseanskuelser: Rum og Tid.
Rummet er Formen for den udvortes Sands, Tiden for den indvortes. Abstrahere vi fra
Gjenstandene udenfor os, faae vi Rummet som den rene, almindelige Form tilbage;
abstrahere vi fra de Forandringer, der iagttages ved den indvortes Sands, bliver
den tomme Tid tilbage.

At Rums- og Tidsanskuelser ere aprioriske, kan oplyses deels fra et metaphysisk,
deels fra et transcendentalt Synspunkt. Den metaphysiske Overveielse godtgjør, at
Anskuelserne: Rum og Tid umulig kunne være laante fra Erfaring, da man netop maa
forudsætte dem for at kunne gjøre en Erfaring; den godtgjør endvidere, at Rum og
Tid ikke ere Begreber, men Anskuelser: de enkelte Rum indeholdes som Dele i det
almindelige Rum, ligeledes de enkelte Tider i den almindelige Tid, medens de almene
Begreber derimod altid indbefatte de særlige som Led under sig, aldrig som Dele i
sig. Den transcendentale Overveielse beviser Rummets og Tidens aprioriske
Oprindelse paa en indirect Maade, idet den viser, at visse Videnskaber, som f. Ex.
den rene Mathematik, vilde være umulige, dersom der ikke gaves almindelige og
nødvendige Sætninger om Rum og Tid; men almindelige og nødvendige Sætninger haves
aldrig fra Erfaringen, de haves kun \textit{a priori}.

b) „Den transcendentale Analytik“ behandler de rene Forstandsbegreber. Forstanden
er, negativt bestemt, en ikke sandselig Erkjendelsesevne; udenfor Anskuelsen gives
% print: line begins „n Erkjendelse" (a faint/broken „en"); read „en".
der kun en Erkjendelse ved Begreber. Til Anskuelsen som Receptivitet
% ---- printed p.25 (PDF 33) ----
svarer Tænkningens Spontaneitet: Anskuelser beroe paa „Affectioner,“ Begreber paa
„Functioner.“ Anskuelsen gaaer umiddelbart paa Gjenstanden, Begrebet derimod
middelbart formedelst en anden Forestilling. Forstanden dømmer; Dommen er en
middelbar Erkjendelse, en Forestilling om en Forestilling.

Naar vi da abstrahere fra en Doms Indhold, og kun see paa den blotte Forstandsform,
finde vi, at Tænkningens Function lader sig bringe under fire Titler, af hvilke
enhver igjen indeholder tre Momenter. I Dommene ligge nemlig, forsaavidt de
inddeles, efter Qvantiteten, i almindelige, særskilte og enkelte, Kategorierne:
Alhed, Fleerhed, Eenhed; forsaavidt de inddeles, efter Qvaliteten, i bekræftende,
benegtende, (limiterende), Kategorierne: Realitet, Negation, Limitation; forsaavidt
de, efter Relationen, deles i kategoriske, hypothetiske og disjunctive,
Kategorierne: Substans og Inhærens, Aarsag og Virkning, Vexelvirkning; og endelig,
naar de adskille sig, efter Modaliteten, i problematiske, assertoriske og
apodiktiske, Kategorierne: Mulighed, Virkelighed og Nødvendighed.

At disse tolv Kategorier, hvoraf de øvrige lade sig udlede, ikke skyldes
Erfaringen, men ere aprioriske Begreber, sees deraf, at ogsaa disse ere ligesaa
nødvendige Betingelser for en Erfarings Mulighed, som de rene Anskuelser. I og for
sig ere Kategorierne tomme Former, der kun faae Materialet tilført gjennem
Anskuelsen; først idet Anskuelsens Sandseindhold gaaer op i Forstandens Form,
kommer Erfaringen. Dog møder os her en Vanskelighed. Gjenstandene, der ere af
sandselig Natur, ere jo ueensartede med Begreberne: hvorledes gaaer det da til, at
aposterioriske Sandsegjenstande kunne subsumeres under rene Forstandsbegreber, og
Domme \textit{a priori} deraf dannes? Dette skeer derved, at et Tredie, som paa den
ene Side er apriorisk, paa den anden Side sandseligt, træder imellem, nemlig, de
rene Anskuelser: Rum og Tid. For i denne Henseende at bringe Begrebet til inderlig
Forening med Anskuelsen, danner Indbildningskraften „det tran-
% ---- printed p.26 (PDF 34) ----
scendentale Schema.“ Saaledes har Qvantiteten til sit almindelige Schema Tidsrækken
eller Tallet, saa at Størrelsens rene Forstandsbegreb møder Anskuelsen, idet vi i
Indbildningen frembringe Enere, og derved gaae fra Eenhed gjennem Fleerhed til
Alhed. Qvalitetens Schema er Tidsindholdet: det Reale udfylder en given Tid, og,
naar Realitetens Indhold er borte, anskues Negationen i Tidens Tomhed.
Relationskategoriernes Schema er Tidsordenen, thi Substantialitet er det Reales
Vedvaren i Tiden, Causalitetens Led følge med Nødvendighed paa hinanden i Tiden, og
Vexelvirkningen er to eller flere Substansyttringers regelmæssige Sammenværen i
Tiden. Modalitetens Schema er Tidsindbegrebet: Mulighedens Schema betegner en
Forestillings Samstemmen med Tidens Betingelser i Almindelighed, Virkelighedens
Schema en Gjenstands Tilværen i en given Tid, og Nødvendighedens en Gjenstands Væren
til alle Tider.

Da de Gjenstande, vi kunne iagttage, alle tilsammen falde i Tiden, maae de
nødvendigviis ogsaa falde under hine Tidens Schemata. Dette bevises i
„Grundsætningernes Analytik.“ I Spidsen for de analytiske Domme staaer Modsigelsens
Grundsætning; men de synthetiske Dommes øverste Grundsætning fastsætter, at enhver
mulig Erfaring „maa staae under de Betingelser, der ere nødvendige for at bringe
Anskuelsens mangfoldige Indtryk til synthetisk Eenhed.“ α) Alle Phænomener ere da,
eftersom de kun kunne opfattes under Rummets og Tidens Former, extensive
Størrelser; „Anskuelsens Axiomer.“ β) Alle Phænomener ere, efterdi enhver
Fornemmelse kan forstærkes og svækkes i en bestemt Grad, intensive Størrelser;
„Iagttagelsens Anticipationer.“ γ) Alle Phænomener ere ifølge
Relationskategoriernes Schema nødvendige Tidsbestemmelser undergivne; det
Substantielle bliver i Tiden; det Accidentelle vexler i Tiden; Accidensernes
timelige Vexel udtrykker et Forhold mellem Aarsag og Virkning; „Erfaringens
Analogier.“ δ) Endelig ere alle Phænomener, i Overeensstemmelse med
Modalitetskategoriernes
% ---- printed p.27 (PDF 35) ----
Schema, bestemte saaledes, at hvad der samstemmer med Erfaringens formelle
Betingelser, er muligt; hvad der samstemmer med dens materielle Betingelser,
virkeligt; hvad der udviser almeengyldige Erfaringsbetingelser, nødvendigt; ---
„den empiriske Tænknings Postulater.“

c) „Den transcendentale Dialektik“ undersøger Fornuftens Ideer.

Vor Erkjendelse begynder med Sandserne, gaaer derfra over til Forstanden, og ender
i Fornuften som i sit Høieste. Er Forstanden Evnen til at bringe Phænomenerne under
en Eenhed formedelst Regler, saa er Fornuften igjen Evnen til at bringe
Forstandsreglerne til Eenhed under Principer. Forstanden udvikler det Betingede,
Fornuften paaviser det Ubetingede; Forstanden afgiver Præmisserne, Fornuften
udbrager Slutningen. Af de tre logiske Fornuftslutninger, nemlig: den kategoriske,
den hypothetiske og den disjunctive, udledes Fornuftens tre speculative Ideer:
Ideen om Sjælen som en tænkende Substans (den psychologiske Idee), Ideen om Verden
som alle Phænomeners Indbegreb (den kosmologiske Idee), Ideen om Gud, som den
øverste Betingelse for Muligheden af Alt (den theologiske Idee).

At Fornuftideerne ikke ere abstraherede af Erfaringer, sees umiddelbart deraf, at
der i hele Erfaringsverdenen ikke kan paavises nogen Gjenstand, der svarer til
nogen af hine Ideer. At disse Ideer ikke destomindre have Gyldighed, nemlig
apriorisk Gyldighed, forstaaes deraf, at de ere Forstandens uundværlige
Ledestjerner, Erkjendelsens regulative Principer. Vi ordne vore subjective Evner
ved Hjælp af Sjælens Idee, Betingelsernes uendelige Rækker ved Hjælp af Verdens
Idee, Aarsagernes opadstigende Trin ved Hjælp af Guds Idee.

\begin{center}{\bfseries 2)\quad Mod Dogmatismen.}\end{center}

Det Aprioriske er, ifølge Kant, alene i vor Bevidsthed, og har Intet med Tingenes
Væsen at gjøre. Allerede „den
% ---- printed p.28 (PDF 36) ----
transcendentale Æsthetik“ advarer os imod, at betragte Rum og Tid som værende
udenfor os. Vilde man opstille den Sætning, at Tingene ere i Rum og Tid, da vilde
det være en Skuffelse; ikke i sig, men kun for os ere Tingene i Rum og Tid, og det
saaledes, at Phænomenerne for den udvortes Sands ere saavel i Rummet som i Tiden,
Phænomenerne for den indvortes Sands kun i Tiden. Ligesaa eftertrykkelig
indskjærper „den transcendentale Analytik,“ at Forstandens Kategorier kun have
Almeengyldighed, forsaavidt de gjælde for os, idet vi anvende dem paa Tingene som
Erfaringsgjenstande, men ikke paa Tingene i og for sig. Idet vi see Objecterne
gjennem subjective Former, og opfatte dem i subjective Begreber, ere vi som de, der
see Tingenes Farver gjennem malede Briller.

Den Splid mellem Tænken og Væren, Kant saaledes hidførte ved fra Først til Sidst at
indskjærpe Adskillelsen af Tingen som Erfaringsgjenstand og Tingen i sig
(\textit{das Ding an sich}), en Adskillelse, der endvidere belyses ved den Maade,
hvorpaa han med Hensyn paa \textit{Phænomena} og \textit{Noumena} efterviser
„Reflexionsbegrebernes Amphibolie,“ er ret egnet til at bibringe den ukritiske
Dogmatisme et Grundstød; dette Stød bliver især føleligt i den Kritik,
Fornuftideerne maae gjennemgaae i „den transcendentale Dialektik.“

Dialektiken, „\textit{eine Logik des Scheins},“ skal opklare de Illusioner, hvori
Fornuften ved at anvende Forstandens Kategorier paa det Ubetingede, nødvendigviis
maa hilde sig. Dette „transcendentale Skin,“ dette Fornuftens optiske Selvbedrag
viser sig i de forskjellige Fornuftideer fra en forskjellig Side.

a) Den rene Fornufts Paralogismer. Den rationelle Psychologi gjør Sjælen til „en
Sjæleting med Ulegemligheden (Immaterialitetens) Attribut; til en Enkeltsubstans
med Uforkrænkeligheden (Incorruptibilitetens) Attribut; til en numerisk-identisk
intellectuel Substans med Personlighedens
% ---- printed p.29 (PDF 37) ----
Prædicat; til et rumløst, tænkende Væsen med Udødeligheden Prædicat.“ Disse den
rationelle Psychologies Sætninger ere lutter Tilsnigelser. Til Grund for de fire
herhen hørende Paralogismer kan ikke ligge Andet end den simple, i og for sig tomme
Forestilling: „Jeg tænker“; men dette: „jeg tænker“ er jo hverken Anskuelse eller
Begreb, det er kun en Bevidsthedsact, der ledsager, bærer og forbinder alle mine
Forestillinger og Begreber. At Jeg, der tænker, altid kan betragtes som Subject, et
Noget, der ikke blot som et Prædicat hænger ved Tænkningen, er en apodiktisk, ja
endog identisk Sætning; men den betyder ikke, at Jeg som Object er et for mig selv
bestaaende Væsen eller Substans. Beviserne for Udødelighed hvile paa skuffende
Slutninger. „Jeg kan vel adskille min rene Tænken ideelt fra mit Legeme; men deraf
følger ingenlunde, at min Tænken ogsaa kan vedvare reelt, efterat være adskilt fra
Legemet.“

b) Den rene Fornufts Antinomier. Til at bestemme de Hovedsider, fra hvilke
Verdensideen viser sig, veiledes man ved de i Tavlen angivne Kategorier. I Henseende
til det Qvantitative, bliver her da Noget at fastsætte om Totaliteten af Verdens
Tider og Rum, og hvad Qvaliteten angaaer, Noget om Materiens Deelbarhed; i Henseende
til Relationen maae de i Verden givne Virkningers Causalitet paavises, og, hvad
Modaliteten angaaer, det Tilfældiges Betingelser forfølges, indtil de tilsidst gaae
op i det Heles Nødvendighed.

Ved ethvert af disse Punkter staaer Paastand mod Paastand (Antithesis mod Thesis)
med Beviis af lige Kraft. Imod Thesis: „Verden har en Begyndelse i Tiden og er
indesluttet inden Rummets Grændse“, gjælder Antithesis: „Verden er saavel i
Henseende til Tiden som til Rummet uendelig“; imod Thesis: „Enhver sammensat
Substans i Verden bestaaer af enkelte Dele“, Antithesis: „Ingen sammensat Ting i
Verden bestaaer af enkelte Dele“; imod Thesis: „Foruden Causaliteten efter
Naturlove gjælder i Verden en

% ---- printed p.30 (PDF 38) ----
Frihedens Causalitet“, Antithesis: „der gives ingen Frihed, men Alt i Verden skeer
alene efter Naturlove“; imod Thesis: „Til Verden hører Noget, der enten som en Deel
af den eller som dens Aarsag er et ligefrem nødvendigt Væsen“, Antithesis: „Der
existerer intet ligefrem nødvendigt Væsen som Verdensaarsag, hverken i Verden eller
udenfor den.“

Disse kosmologiske Antinomier vise os da den dialektiske Kamp, den Selvmodsigelse,
hvori Fornuften hilder sig, naar den ved Verdens Idee vil tænke sig Verdens
objective Væsen.

c) Den rene Fornufts Ideal. „Kategorien“ er tom, men kan dog finde sin Gjenstand i
Erfaringen; „Ideen“, der slet ikke har nogen Erfaringsgjenstand, er altsaa fjernere
fra den objective Realitet end Kategorien. „Men endnu mere end Ideen synes det at
være fjernet fra objectiv Realitet, som jeg kalder Idealet, og hvorunder jeg
forstaaer Ideen, ikke blot \textit{in concreto}, men \textit{in individuo} d. e. som
en enkelt ved Ideen alene bestemmelig og bestemt Ting.“ Efter at have viist,
hvorledes Fornuften bestemmer sit Ideal paa forskjellig Maade, nemlig som Urvæsen
(\textit{ens originarium}), forsaavidt den har det i sig, som det høieste Væsen
(\textit{ens summum}), forsaavidt den har det over sig, og som Væsnernes Væsen
(\textit{ens entium}), forsaavidt Alt er betinget af det, gaaer Kant over til et
Angreb paa Theologiens „speculative Fornuft“, og driver i sin Kritik af de
dogmatiske Beviser for Guds Tilværelse Adskillelsen af Tænken og Væren til det
Yderste.

Blandt de tre Beviser, det ontologiske, det kosmologiske og det
physikotheologiske, maa Eftertrykket vel især falde paa det første. „Om
Umuligheden af et ontologisk Beviis for Guds Tilværelse“ bør da herved fremhæves
Følgende. Beviset er i Korthed dette: „Det allerrealeste Væsen er muligt; under de
mange Realiteter hører ogsaa Tilværelse; det ligger altsaa i det allerrealeste
Væsens Begreb, at dette Væsen selv nødvendigviis maa være til.“ Man oplyser dette
ved
% ---- printed p.31 (PDF 39) ----
Exempler og blænder ved Exempler. Naar jeg ophæver Prædicatet i en identisk Dom, og
beholder Subjectet, saa fremkommer ganske rigtig en Modsigelse, og derfor siger
man, at Prædicatet tilkommer Subjectet med Nødvendighed. Men ophæver jeg Subjectet
tilligemed Prædicatet, saa fremkommer ingen Modsigelse; „\textit{denn es ist nichts
mehr, welchem widersprochen werden könnte}.“ At sætte et Triangel og dog ophæve
dets tre Vinkler, er en Modsigelse; men at ophæve et Triangel tilligemed dets tre
Vinkler, er ingen Modsigelse. Saaledes med det absolut nødvendige Væsen. Gud er
almægtig; Almagten kan ikke ophæves, naar I sætte Gud; men, naar I sige: Gud er
ikke, er der hverken Almagt eller andet Prædicat. Væren er, som Copula i Dommen,
intet reelt Prædicat. Det Virkelige indeholder i Begrebet ikke mere end det blot
mulige. Hundrede virkelige Daler indeholde ikke det mindste mere end hundrede
mulige; thi i modsat Fald vilde Begrebet ikke være adæqvat. Et Begreb mister ikke
en eneste af sine Egenskaber, fordi dets Gjenstand ikke er til, og vinder ikke en
eneste Egenskab derved, at dets Gjenstand findes existerende.

„Paa det saa berømte ontologiske (cartesianske) Beviis er altsaa al Møie og Arbeide
spildt; et Menneske vilde ved blotte Ideer ligesaa lidt blive rigere paa Indsigt,
som en Kjøbmand paa Formue ved at føie nogle Nuller til sin Kassebeholdning.“ Fra
Tænken til Tilværen gives ingen speculativ Overgang.

\begin{center}{\bfseries 3)\quad Den transcendentale Idealisme.}\end{center}

Da Kant ved „en transcendental Erkjendelse“ udtrykkelig forstaaer en saadan, „som
ikke egentlig har at gjøre med Gjenstandene, men kun med vor Erkjendelse af
Gjenstandene, forsaavidt denne \textit{a priori} skal være mulig“, saa sees allerede
heraf, hvorfor han selv kalder sin Kriticisme en Transcendentalphilosophie. Denne
Philosophie vil nu vistnok paa sin Viis hævde Erfaringsverdenens Gyldighed; men det
vilde
% ---- printed p.32 (PDF 40) ----
Berkeleys Spiritualisme jo ogsaa paa sin Viis. Den Maade, hvorpaa
Subjectivitetsmomentet bliver det i al vor Erkjendelse Overgribende, charakteriserer
den kantiske Transcendentalphilosophie som Idealisme.

I hvilket Forhold denne transcendentale Idealisme udtrykkelig har sat sig saavel til
Skepticismen som til Dogmatismen, udtales især i „Methodelæren“, der behandler den
rene Fornufts \emph{Disciplin}, \emph{Kanon}, \emph{Architektonik} og
\emph{Historie}.

„Vilde man,“ hedder det, „spørge den koldblodige, til Dommens Ligevægt egentlig
skabte, David Hume, hvad der da bevægede ham til ved Betænkeligheder, som kun en
møisommelig Grublen kunde udfinde, at undergrave den for Menneskene saa trøstelige
og nyttige Indbildning, at deres Fornuftindsigt slaaer til, naar Begrebet af det
høieste Væsen skal bestemmes, da vilde han vel svare: Intet, uden den Hensigt at
bringe Fornuften videre i Selverkjendelse, og derhos en vis Uvillie over den Tvang,
man vil paalægge Fornuften, idet man gjør sig til af den, og dog hindrer den i at
aflægge en frimodig Bekjendelse om de Svagheder, som den kun ved alvorlig Prøvelse
kan opdage hos sig selv.“

Saa utilfredsstillende Skepticismen end er i og for sig, saa berettiget er den
ligeover for en ukritisk Dogmatisme, der ikke veed, hvad Forstanden formaaer, og
ikke har prøvet Grændserne for Erkjendelsens Mulighed efter Principer. „\textit{Und
so führt der Skeptiker, der Zuchtmeister des dogmatischen Vernünftlers, auf eine
gesunde Kritik des Verstandes, ja der Vernunft selbst}.“

I den rene Erkjendelse som Videnskab maa Fornuften derimod altid gaae dogmatisk
tilværks, d. e. føre strenge Beviser af sikkre Principer \textit{a priori}; men at
føre strenge Beviser af sikkre Principer \textit{a priori} er ikke det samme som
(\textit{Wolffs methodus mathematica}) at efterligne Mathematiken. Mathematikens
Grundighed beroer paa Definitioner, Axiomer og Demonstrationer; Philosophiens paa
Explica-
% ---- printed p.33 (PDF 41) ----
tioner, Deductioner og akroamatiske Beviser. Det sømmer sig ikke for Philosophien
„\textit{mit einem dogmatischen Gange zu strotzen und sich mit den Titeln und
Bändern der Mathematik auszuschmücken, in deren Orden sie doch nicht gehört}.“ De
apodiktiske Sætninger, (de være nu beviislige, eller ogsaa umiddelbart visse), dele
sig i Dogmata og Mathemata. En direct synthetisk Sætning af Begreber er et
\emph{Dogma}, hvorimod en lignende Sætning ved Begrebernes Construction er et
\emph{Mathema}. Gives der nu i den rene Fornufts speculative Anvendelse ikke et
eneste Dogma, saa er ogsaa enhver dogmatisk Methode, den maatte nu være laant hos
Mathematikeren, eller paa anden Maade bleven til eiendommelig Maneer, i sig
upassende. „\textit{Denn sie verbirgt nur die Fehler und Irrthümer, und täuscht die
Philosophie, deren eigentliche Absicht ist, alle Schritte der Vernunft in ihrem
klärsten Lichte sehen zu lassen}.“

\begin{center}{\bfseries \S\,7.}\end{center}

„I enhver Erfaring haves to Factorer satte mod hinanden: Jeget og Tingen,
Intelligensen og dens Gjenstand. Abstraherer man fra Jeget, saa faaer man en Ting
\emph{i og for sig}, og maa da opfatte Forestillingen som et Product af Gjenstanden;
abstraherer man fra Gjenstanden, faaer man et Jeg \emph{i og for sig}, og maa da
holde sig til den selvvirksomme, sit Indhold frembringende Bevidsthed. Hiint er
\emph{Dogmatismen}, dette \emph{Idealismen}. Disse To ere indbyrdes uforenelige; et
Tredie gives ikke: her maa vælges. For at dømme Dogmatismen og Idealismen imellem,
mærke man sig Følgende: Jeget forekommer i Bevidstheden, Tingen i sig er en Digtning
af Bevidstheden. Idet Dogmatismen skal forklare Forestillingens Oprindelse fra
Tingen i sig, øver den det mislige Kunststykke at gaae bagved Bevidstheden, og
paakalde Væren; men Væren virker kun Væren, ikke Bevidsthed; altsaa kan kun det
Standpunkt være sandt, der vælges, ikke i Væren, men i Bevidstheden,
% ---- printed p.34 (PDF 42) ----
ikke i Gjenstanden, men i Jeget selv.“ I denne Aand har I. G. Fichte († 1814)
draget en Kantianismens Conseqvens og (\textit{Erste Einleitung 1797}) indledet sin
„\textit{Wissenschaftslehre}“.

Medens Fichte opløser „Tingen i sig“, fastholder Jeget, og gjennemfører den
\emph{subjective Idealisme}, bliver den kantiske Lære belyst fra en modsat Side,
idet I. F. Herbart (Metaphysik) fæster Opmærksomheden paa „Tingen“, opfatter
endogsaa Jeget som en Ting for sig, et „Reale med mange Egenskaber“ (Psychologie),
og paa sin Viis uddanner en \emph{kritisk Realisme}. „Erfaringsstoffet er
Philosophiens Grundlag; men som et blot Givet staaer det endnu udenfor Philosophien.
Hvad er da Philosophiens første Skridt, den Handling, hvormed den begynder?
Erfaringen har paatvunget os Begreber, der ikke lade sig tænke. Vistnok tænkes de
paa en Maade af Hverdagsforstanden; men det er en dunkel og forvirret Tænken, der
ikke gjennemskuer de modstridende Særkjender. Den skjærpede Tænkning, Metaphysikens
logiske Analyse, opdager i Erfaringsbegreberne (Rum, Tid, Bevægelse o. s. v.)
Modsigelser, modstridende, hinanden udelukkende Kjendemærker. Hvad er herved at
gjøre? Bortkastes kunne disse Begreber ikke; de ere os givne, og vi have kun det
Givne at holde os til. Beholdes kunne de heller ikke, da de ere utænkelige, i sig
umulige. Den eneste Udvei er: vi maae omarbeide dem. \emph{Erfaringsbegrebernes
Omarbeidelse er Speculationens Opgave.}“

Vende vi, efter denne Antydning af den herbartske Retning tilbage til Fichte, da
maae vi, for at forstaae den subjective Idealisme, vel indprænte os
„Videnskabslærens“ øverste Grundsætninger; thi i disse Grundsætninger er
„Videnskabslæren“ Erkjendelseslære.

α) \emph{Thesis.} Indrømmer man, at Sætningen: $A = A$ er indlysende ved sig selv,
da indrømmer man med det Samme Bevidsthedens Evne til at sætte Noget. Man sætter
ikke, at $A$ er, men kun, at, dersom $A$ er, saa er det $A$.
% ---- printed p.35 (PDF 43) ----
Ville vi nu, istedetfor at see paa det Satte, see paa den Sættende, eller rettere:
paa Identiteten af den Sættende og det Satte, da forvandler Sætningen: $A = A$ sig
til Jeg $=$ Jeg. Medens vi istedetfor $A = A$ ikke kunde sige, at $A$ er, saa kunne
vi nu, istedetfor: Jeg $=$ Jeg, sige: Jeg er. Dette paa sig selv Grundede er Grund
til al Handling i den menneskelige Aand, og altsaa det oprindelige Særkjende paa
Bevidsthedsvirksomheden i og for sig.

β) \emph{Antithesis.} Sætningen Non-$A$ ikke $=A$ er som en fri Urhandling efter sin
Form ligeledes ubetinget, men derimod efter sit Indhold, sin Materie, betinget,
fordi Antithesis forudsætter Thesis, fordi der, naar et Non-$A$ skal sættes,
allerede forud maa være sat et $A$. Ved den første Grundsætning: $A = A$, var
Urhandlingens Form en Sætten, ved den anden er den en Modsætten. Hvad Non-$A$ er,
vides endnu ikke; kun dette vides, at Non-$A$ er Modsætning til $A$. Men nu er $A$
sat ved Jeget; her er kun forudsat et Jeg: det Jeget Modsatte er altsaa Ikke-Jeg.

γ) \emph{Synthesis.} Den tredie Grundsætning er, efterdi den bestemmes ved de to
foregaaende, næsten heelt igjennem istand til at modtage et Beviis. Opgaven for den
Handling, der udtrykkes ved den tredie Grundsætning, er given med de to foregaaende,
men ikke Opgavens Løsning. Paa den ene Side ophæves Jeget fuldstændig ved
Ikke-Jeget: Jeget kan ikke være sat, forsaavidt Ikke-Jeget er sat. Paa den anden
Side er Ikke-Jeget, sat af Jeget, selv et Jeg, sat i Bevidstheden. Jeget ophæves
altsaa ikke ved Ikke-Jeget; det er ophævet, og dog alligevel ikke ophævet. Dette er
Modsigelsen. Til Løsningen af denne Modsigelse maae vi see at finde et $X$, hvori de
hinanden udelukkende Grundsætninger kunne forenes. Men hvorledes kunne Væren og
Ikke-Væren, Realitet og Negation forenes med hinanden, uden at tilintetgjøre
hinanden? Svar: ved gjensidig Indskrænkning! Det omspurgte $X$ er altsaa en Skranke,
og Selv-indskrænkning Jegets Urhandling.
% ---- printed p.36 (PDF 44) ----
Forsaavidt Jeghedens Væsen, det rene Jeg, opfattes som et Princip, der udvikler sig
i det empiriske Jegs Vexelvirkning med en Verden af Skranker, har Subjectiviteten i
den Grad faaet Overvægt, at heri indeholdes en naturlig Opfordring til at lægge et
Idealitetens Eftertryk ogsaa paa det Objective.

I sit Skrift: „\textit{Von der Weltseele}“ har F. V. J. Schelling († 1854), endnu
med Fichtes Forudsætninger, søgt at paavise, hvorledes Materiens Begreb nødvendig
fremgaaer af den menneskelige Aands Natur og Anskuelse. „Sindet (\textit{das
Gemüth}) er Eenheden af en uindskrænket og indskrænkende Kraft. Fraværelse af
Skranker vilde ligesaavel som den absolute Indskrænkethed umuliggjøre Bevidstheden.
Kun idet den Kraft, der stræber ud i det Uindskrænkede, indskrænkes ved den
modsatte, medens den indskrænkede selv omvendt frigjøres fra sine Skranker, er
Følen, Iagttagelse, Erkjendelse tænkelig. Et Tilsvarende sees i Naturen. Ikke
Materien som Materie, men de Kræfter, hvis Eenhed Materien udgjør, er her det
Første. Materien er Attractionens og Repulsionens stedse vordende Product. Materien
er, hvad den er, ved Kraften; men Kraft er det i Materien Immaterielle, Kraft i
Naturen er det, der kan sammenlignes med Aanden. Naturen og Aanden maae oprindelig
være forenede i en høiere Identitet. Aandens Organ for Opfattelsen af Naturen, er
Anskuelsen; det er Anskuelsen, der tager det ved tiltrækkende og frastødende Kræfter
begrændsede og opfyldte Rum i Besiddelse.“

Med disse Tanker har Schelling allerede betraadt Naturphilosophiens objective Grund.
„Al Viden“, hedder det (\textit{System des transcendentalen Idealismus}) „beroer paa
et Subjects Overeensstemmelse med et Object. Indbegrebet af det blot Objective er
Natur; Indbegrebet af det blot Subjective er Jeget eller Intelligensen. Til
Foreningen af begge Sider ere to Veie mulige; enten gjør man Naturen til det Første,
og spørger, hvorledes det Intelligente kommer til den,
% ---- printed p.37 (PDF 45) ----
d. e. man stræber at udlægge den i rene Tankeformer --- Naturphilosophie; eller man
gjør Subjectet til det Første, og spørger, hvorledes Objecterne fremgaae af
Subjectet --- Transcendentalphilosophie. Al Philosophie maa gaae ud paa, enten af
Naturen at gjøre en Intelligens, eller af Intelligensen at gjøre en Natur. Ligesom
Transcendentalphilosophien har at underordne det Reelle under det Ideelle, saaledes
maa Naturphilosophien søge at forklare det Ideelle af det Reelle. Men begge ere kun
een og samme Videns to Poler, der gjensidig søge hinanden; derfor maa man ogsaa,
naar man gaaer ud fra den ene Pol, nødvendig komme til den anden.“

„Alle Philosophiens Dele maae forebrages i een Continuitet, og den hele Philosophie
maa forebrages som det, den er, nemlig som Bevidsthedens fortløbende Historie, for
hvilken det i Erfaringen Nedlagte kun tjener saa at sige som Mindesmærke og som
Document.“

Schellings transcendentale Idealisme kan nu betragtes som Forløber for Hegels
„\textit{Phänomenologie des Geistes}“; hvad Schelling lovede, opfyldte Hegel
(† 1831).

„Bevidstheden har i Almindelighed, eftersom Gjenstanden viser sig, tre Trin.
Gjenstanden er nemlig enten et ligeoverfor Jeget staaende Object, eller den er Jeget
selv, eller noget Gjenstændigt, der dog ligesaa meget tilhører Jeget, Tanken“. De
tre Trin ere altsaa: Bevidsthed, Selvbevidsthed, Fornuft.

a) \emph{Bevidsthed.} Paa det første Trin er Bevidstheden sandsende, iagttagende,
forstandig.

α) \emph{Den sandsende Bevidsthed er umiddelbar Vished om en udvortes Gjenstand.}
Udtrykket for en saadan Gjenstands Almindelighed er, \emph{at den er, og er denne,
nu efter Tiden, her efter Rummet}. Gjenstanden er altsaa derved aldeles forskjellig
fra alle andre Gjenstande, og for sig tilstrækkelig bestemt. Men dette Nu er ikke
mere, et andet er traadt i dets Sted, men som dog ogsaa er et Nu: Nuet
% ---- printed p.38 (PDF 46) ----
er i sin Forsvinden et Blivende, nemlig det Almene. Dette Her, jeg mener og
paapeger, har et Høire, et Venstre, et Foroven, et Forneden, et Bagved, et Foran i
det Uendelige. Dette Her er altsaa ikke det enkelte, umiddelbart bestemte, men: nu
her, nu her, nu her! et Indbegreb af Mange, og i disse Mange dog det Ene, altsaa det
Almene.

β) Det Almene, forsaavidt det endnu er svævende mellem Sandselighed og Reflexion, er
Gjenstand for den \emph{iagttagende} Bevidsthed; \emph{her er Tingen med dens
Egenskaber}. De sandselige Egenskaber ere paa een Gang umiddelbart hver for sig, og
dog i Forhold til hverandre, tilhøre Tingen som denne enkelte Ting, og ere dog saa
almindelige, at de, adskilte fra hverandre, kunne findes i mangfoldige andre Ting.
Ting og Egenskaber forandre sig, og have dog i denne Forandring deres Bestaaen: det
Foranderlige er et Ydre, det Bestaaende et Indre.

γ) Forsaavidt Gjenstanden er et Indre med sit Ydre, en Kraft i sin Yttring, er den
til for den \emph{forstandige} Bevidsthed. Forstanden seer i Tingenes vexlende Ydre
en Række Phænomener, hvis Indre bestaaer i et System af blivende Forhold: Lovene.
Tingene ere under Lovene, og Lovene ere Tingenes Love; Tingene og Lovene ere to
Sider af det Samme: det Sande, Tingenes Væsen, er Tanken, Begrebet.

b) \emph{Selvbevidsthed.} Idet Bevidstheden vender sig mod Tanken som Tanke, mod
Væsenet som Væsen, mod det Indre som Indre, faaer den sin egen Reflexion og derved
sig selv til Gjenstand. Selvbevidstheden har i sin Dannelse eller Bevægelse tre
Trin: Attraaen, Forholdet af Hersken og Tjenen, den almene Selvbevidsthed.

α) \emph{I Attraaen forholder Selvbevidstheden sig til sig som enkelt, alene værende
Selvbevidsthed}; det er en selvløs Gjenstand, der skal tilfredsstille
Begjærligheden; men, idet Gjenstanden fortæres, og Tilfredsstillelsen indtræder, har
% ---- printed p.39 (PDF 47) ----
Bevidstheden i sin Selvfølelse forstaaet, at den som subjectiv tillige væsentlig maa
være objectiv.

β) Selvbevidsthedens Begreb som et Subject, der tillige er et Object, finder sit
fuldstændige Udtryk deri, at Selvbevidstheden i sit Object seer en anden
Selvbevidsthed. Idet de to ligeoverfor hinanden stillede Selvbevidstheder hver efter
sit Begreb ere selvstændige, opkommer Spændingen mellem det at være Herre og det at
være Tjener. Tjeneren er selvløs, forsaavidt hans Selv er et Moment i Herrens
Bevidsthed; men Herren er jo ogsaa selvløs, forsaavidt hans Væren Herre beroer paa,
at han har en Tjener; han er kun Herre formedelst Tjeneren. Med denne Vexelvirkning
gjør Selvbevidstheden Overgang til articuleret Almeenvillie, til positiv Frihed.

γ) \emph{I den almene Selvbevidsthed veed den Enkelte sig} som tjenende de Manges
Villie, idet han væsentlig tjener sig selv. Her er ikke en Sum af Villier, men de
mange Villier ere Momenter i Almeenvillien. Denne Selvbevidsthed er Grundlaget for
alle Dyder: Kjærlighedens, Ærens, Venskabets, Tapperhedens, Opoffrelsens, Hæderens
o. s. v.

c) \emph{Fornuft.} Fornuften er den høieste Forening af Bevidsthed og
Selvbevidsthed i Viden af Gjenstanden og Viden af sig selv. Fornuften er den Vished,
at dens Bestemmelser ligesaameget ere gjenstændige, Bestemmelser ved Tingenes
Væsen, som at de ere Subjectets egne Tanker. Fichtes „Jeg“ og Herbarts „Reale“ ere
herved blevne Momenter; Fornuftens Viden er ikke blot subjectiv Vished, men Sandhed;
men Sandheden er Eenhed af Vished og Gjenstændighed, af Viden og Væren.

„Maalet, den absolute Viden, eller den sig som Aand vidende Aand har til sin Vei
Erindringen om Aanderne, som de i sig selv ere, og fuldføre deres Riges
Organisation. Dens Opbevaren i Henseende til deres frie, i Tilfældighedens Form
fremtrædende Tilværelse er Historien, men i Henseende til deres begrebne
Organisation den phænomenale
% ---- printed p.40 (PDF 48) ----
Videns Videnskab; begge tilsammen, den begrebne Historie, danne den absolute Aands
Erindring og Domsted, den Thrones Virkelighed, Sandhed og Vished, uden hvilken den
vilde være det Livløse, Eensomme; kun af dette Aanderiges Kalk fremsprudler
Uendeligheden for den.“

\begin{center}\rule{2.5cm}{0.4pt}\end{center}

\begin{center}
  {\large\bfseries C.}\par\vspace{2pt}
  {\bfseries Sagerkjendelse.}\par\vspace{4pt}
  \rule{2cm}{0.4pt}
\end{center}
\phantomsection
\addcontentsline{toc}{subsection}{\texorpdfstring{C.\quad Sagerkjendelse}{C. Sagerkjendelse}}

\begin{center}{\bfseries \S\,8.}\end{center}

\noindent See vi nu tilbage paa de her omhandlede Erkjendelsestheorier, de
sensualistiske saavel som de idealistiske; da er det os ret paafaldende, at ikke en
eneste iblandt dem alle i sine Conseqvenser tilfredsstiller „sund Forstand“, medens
de dog i deres Samvirken, forsaavidt enhver for sig paa det Kraftigste hævder sin
Side af Sagen, ligesom forene sig om at give „den sunde Menneskeforstand“ Ret.

Heraf kan nu ikke sluttes, at „den sunde Forstand“ er klogere end Philosophien, men
kun, at den med Instinctets Sikkerhed fastholder de Led, hvis indbyrdes Forhold
Philosophien vil forklare. Naar da Philosophien, istedetfor at opklare Leddenes
Forhold, bortforklarer eller forqvakler et væsentligt Led, mærker „den sunde
Forstand“ Uraad, og protesterer.

„Den sunde Forstand“ indrømmer gjerne, at man ikke kan blive sig en Omverden bevidst
uden at optage dens Indtryk og tilsvarende Billeder i sin Forestillingskreds; men
den paastaaer tillige, at den Verden, vi blive os bevidste, maa, saasandt den er en
virkelig Verden, have en virkelig, af vor Bevidsthed uafhængig, Tilværen.
Idealisten kan nu anstrenge sig for at bevise Verdens og Bevidsthedens uopløselige
Eenhed, saameget han vil: „den sunde Forstand“ prote-

% ---- printed p.41 (PDF 49) ----
sterer. At vore Fornemmelser og Forestillinger hidrøre fra umiddelbare
Sandseindtryk, vil Forstanden gjerne indrømme; men, naar Skeptikeren heraf uddrager
den Conseqvens, at Aarsag og Virkning kun ere til i vor Indbildning, at Lovene ikke
ere Love, staaer „den sunde Forstand“ stille, og begynder at protestere.

En Philosophie, der er i aabenbar Modsigelse med „den sunde Forstand“, er falsk; en
Philosophie, der umiddelbart retter sig efter „den sunde Forstand“, er overfladisk.
Philosophien vil ikke vilkaarlig modsige den naturlige Forstand, men kun vise, at den
saakaldte sunde Sands, uden at have nogen Anelse derom, uophørlig modsiger sig selv.

„Den sunde Forstand“ fordrer, at jeg umiddelbart skal blive mig bevidst, at der er
Noget udenfor min Bevidsthed; men, idet jeg bliver mig Noget bevidst, er det jo i min
Bevidsthed; forsaavidt det er udenfor, kan jeg ikke blive mig det bevidst: altsaa skal
jeg blive mig bevidst, at det, der er i min Bevidsthed, ikke er der; hvilken
Modsigelse! Naar Solen skinner, og Træerne speile sig i Aaen, mener „den sunde
Forstand“, at Lyset er Klarhed, enten der er Nogen, der seer det, eller ikke, og at
Træerne virkelig speile sig i Aaen, enten saa denne Speiling afspeiler sig i et seende
Øie, eller ikke. Men Spørgsmaalet er, hvorledes man skal overbevise sig herom. Ved
Synet? Det er en Modsigelse, da her jo handles om Phænomenet, ikke forsaavidt det
sees, men forsaavidt det ikke sees. Ved Tanken? Det er en Modsigelse; Phænomenet
lader sig som Phænomen ikke tænke, thi i saa Fald kunde den Blindfødte jo ogsaa tænke
det. Ved Andres Vidnesbyrd? En tom Udflugt; da det jo her maa gaae den Ene, som den
Anden. Kun for den Seende er Lyset Klarhed; kun, idet Bækken og Træet speile sig i
Øiet, kan Træet selv speile sig i Bækken.

Det falder naturligt, at En tænker tilbage paa den Tid, da han endnu ikke var; og
ligesaa naturligt, at han tænker sig det Øieblik, da han igjen vil have ophørt at være.

% ---- printed p.42 (PDF 50) ----
Ikke desto mindre er her i begge Dele en forunderlig Modsigelse. At Bevidstheden kan
være sig sin egen Tilværen bevidst, er indlysende; men, hvorledes kan den tilværende
Bevidsthed blive sig bevidst som ikkeværende? I Virkeligheden? Nei, det er en
Modsigelse, thi naar Bevidstheden virkelig er, er dens Ikke-Væren jo virkelig
udelukket; og, naar Bevidstheden ikke er, kan den ikke blive sig Noget bevidst. I
Phantasien? Men, naar Bevidstheden kun seer sin Ikke-Væren som en Phantasie,
hvorledes kan den da overbevise sig om, at det Virkelige enten har svaret eller vil
svare til denne Phantasie? „Den sunde Forstand“ kan gjerne rykke frem med sin
Forklaring, men Sagen er, at Forklaringen indeholder ligesaamange Modsigelser, som
det, den skal forklare.

Hvad er Sandhed? Overeensstemmelse mellem det, Bevidstheden antager for virkeligt, og
dette Virkelige selv. Og Usandhed? Naturligviis Strid imellem det, der sættes i
Bevidstheden, og det der er i Virkeligheden. Dersom nu alt Virkeligt med Eet blev
borte, og den tomme Bevidsthed med sine Phantasier og Tanker stod ene tilbage, vilde
den kunne drive den allervilkaarligste Leg med Indbildninger. At Sandheden bliver
borte, naar det virkelige Verdensindhold bliver borte, vil sund Forstand vist gjerne
% ERRATA (Rettelser, p.42 l.14 fra neden): print „det vilkaarlige“ → „det virkelige“.
indrømme; men det Virkelige som blot virkeligt er jo, naar Bevidstheden selv bliver
borte, heller ikke det Sande. En Steen, en Stjerne, en hvilkensomhelst virkelig
Gjenstand er i og for sig hverken noget Sandt eller noget Usandt. For at Forskjellen
mellem Sandt og Usandt kan fremkomme, maa der fældes en Dom om den virkelige
Gjenstand; men for at der kan fældes en Dom, maa der absolut være en dømmende
Bevidsthed. Men er Bevidstheden et uundværligt Reqvisit for Sandheden, kan det
Virkelige for sig umulig være det Sande; hvorledes kan „den sunde Forstand“ da mene
at kunne tænke sig en Tid, da der i \emph{Sandhed} vare virkelige Atomer, virkelige
Elementer, en virkelig Jord, men endnu

% ---- printed p.43 (PDF 51) ----
ingen, slet ingen Bevidsthed? Virkelige Objecter, som umulig kunde objectiveres;
hvilken Modsigelse!

Den ydre, virkelige Verden er i sin rene ubevidste Virkelighed et Intet; vil Nogen
paastaae det Modsatte, maa han jo paastaae det i Kraft af Bevidsthed og altsaa modsige
sig selv.

Hvad er et Evigt? Naar en almeengyldig Sætning erkjendes, seer det slet ikke ud, som
om det Sande først blev til i Erkjendelsens Øieblik; det Sande viser sig tværtimod for
Bevidstheden som det, der altid har været. At enhver Ting er, hvad den er; at to
Størrelser, af hvilke enhver for sig er liig med een og samme tredie, ere indbyrdes
lige; at en Grund nødvendigviis maa have sin Følge, en Virkning sin Aarsag, alle disse
saakaldte „evige Sandheder“ fremstille sig for os som de, der have været forud for al
endelig Viden; og dog ere de i sig Intet, naar man tager Viden bort: hvorledes kan
den i Tiden vordende Bevidsthed blive sig det Ikke-Vordende, det Evige, bevidst?

Hvad betyder Modsigelsen? Et tilværende Intet er en Modsigelse; en rund Fiirkant er en
Modsigelse; en flyvende Hest er ogsaa en Modsigelse. Ikke desto mindre ere disse
Modsigelser af forskjellig Art. En flyvende Hest kan man forestille sig, men ikke
tænke. Det Utænkelige ligger ikke ligefrem i det Logiske, at Prædicatet „flyvende“
modsiger Begrebet „Hest“, thi dette Begreb kunde jo udvides; men deri, at det Hele
saaledes udvidede Begreb staaer i Strid med Erfaringens Betingelser og Love.
Anderledes med den runde Fiirkant, her sees Modsigelsen umiddelbart i Forestillingen
selv. At forestille sig Fiirkanten rund, er umuligt: at tænke sig en rund Fiirkant er
derimod muligt, forsaavidt det er muligt at tænke Prædicatets \emph{Strid} med
Subjectet. Modsigelsen er altsaa ikke blot den, at Prædicatets Strid med Subjectet
gjør Dommen meningsløs; men Modsigelsen er i høiere Forstand denne, at Modsigelsen
selv, dette Utænkelige, dog alligevel lader sig tænke. Dette

% ---- printed p.44 (PDF 52) ----
bliver især klart, naar Opgaven er at tænke det tilværende Intet. Det er jo Intet,
men idet Bevidstheden sætter denne Negation, har Intet netop Tilværelse i Tanken,
altsaa en tænkt Tilværelse.

Ved at grunde lidt paa disse mange Modsigelser vil „den sunde Forstand“ bedst lære at
forstaae, at den ivrige Strid, Philosopherne føre om „Modsigelsens Grundsætning“, er
mere end en Ordstrid.

Herbart og Hegel ere begge enige i, at Erfaringsbegreberne indeholde en Modsigelse; de
ere fremdeles enige i, at Modsigelsen skal ophæves; men om Modsigelsens Forhold til
Tingenes Væsen, ja om Modsigelsens eget Væsen ere de i høieste Grad uenige. Idet
Herbart fastholder Modsigelsens Grundsætning, seer han i Modsigelsen en Vildfarelse,
man skal undgaae, en Skuffelse, man skal fjerne. Indeholde Erfaringsbegreberne
Modsigelser, da ligger Skylden ikke i Tingenes objective Natur, men i den erkjendende
Subjectivitet. Den tænkende Bevidsthed maa altsaa, ved at omarbeide
Erfaringsbegreberne og bortskaffe Modsigelsen, see at fjerne den falske Opfattelse.
Ivrende mod Modsigelsen beskylder Herbart endog den hegelske Philosophie for
Empirisme, da den „af Erfaringen optager Begreberne uforandrede, og trods Indsigten i
deres modsigende Natur, just derved, at de ere empirisk givne, betragter dem som
retfærdiggjorte, ja endog for deres Skyld omskaber Logiken.“ Hegel med sit:
„\textit{Alles ist ein Widerspruch}“ lægger derimod Modsigelsen ind i Tingenes Væsen,
saa at en Væren i Endelighed nødvendigviis er en Væren i Modsigelse (Løsning af de
kantiske Antinomier). Paa Modsigelsen beroer ikke blot Livets Smerte, men ogsaa den
Drift, der driver Tilværelsen frem: Modsigelsen fremkalder Bevægelse, og Bevægelse er
igjen Modsigelsens Løsning. Saalidt som Tilværelsen kan udelukke Modsigelsen, saa lidt
kan Tanken udelukke den: Tankens Puls er Dialektik, Modsigelsen Dialektikens Nerve.

% ---- printed p.45 (PDF 53) ----
\begin{center}{\bfseries \S\,9.}\end{center}

Hvor inderlig Modsigelsen er forenet med Tanken, sees i de logiske Domme. Ordet „Blad“
kan som Ord i Sproget vise hen paa et sandselig Givet i dettes umiddelbare
Singularitet, men ogsaa som Begrebsudtryk betegne et Almindeligt. Et Udraab som:
„dette Blad!“ er i Sprogform ufuldkomment, og giver kun en bestemt Mening, naar En med
det Samme peger paa et virkeligt Blad. Fæstes nu Øiet paa „dette Blad“, kan Tanken
forenes med det Sandselige, idet en Dom udtales: „dette Blad“ er grønt. Men „dette
Blad“ er jo ikke grønt, thi det er guult; „dette Blad“ er altsaa grønt og dog ikke
grønt; „dette Blad“ er baade grønt og guult; her er en Vrimmel af Modsigelser. For det
Første er Prædicatet „grønt“ i Forhold til „dette Blad“ et Almindeligt, medens „dette
Blad“ selv er et Enkelt. Dommen udsiger altsaa, at det Enkelte er et Almindeligt; men
det Enkelte er jo dog alligevel ikke et Almindeligt, men et Enkelt: her er en
Modsigelse. Hvad er dog „dette“? „Dette“ er et Blad, „dette“ er et Bord, „dette“ er en
Hest; „dette“ er altsaa ikke „dette“, men et Andet, end „dette“: lutter Modsigelser.

Gjennem Philosophiens Historie drager sig en Række af Stridigheder, der alle vise hen
til den i Sagen værende Modsigelse. I Stilpons Yttring: „den Kaal, som her falbydes,
er ikke; thi Kaal var allerede for 1000 Aar siden“; ligger den Mening, at det
Individuelle i dets sandselige Umiddelbarhed ikke kan tænkes. Scholastikernes Strid om
Forholdet mellem \textit{res} og \textit{universalia} peger hen paa den samme
Modsigelse, der ogsaa, hvad den nyere Philosophie angaaer, gjentager sig i
Modsætningen mellem Sensualisme og Idealisme.

Locke og Hume ere begge opmærksomme paa den Skuffelse, som kan fremkomme derved, at
det Almene for sig fastholdes i Ordet; og, da den objective Idealisme var rykket
saavidt frem, at Hegels Dialektik antoges at have forvandlet

% ---- printed p.46 (PDF 54) ----
det Sandselige til et i Begrebet svindende Moment, optraadte Feuerbach
(\textit{Philosophie der Zukunft}) med en Kritik, som fordrede ubetinget Sensualisme.

Til Forstaaelse af Idealismens Eensidighed fremhæve vi i denne Sammenhæng Forholdet
imellem det Sandselige og det Tænkelige, imellem det Individuelle og det Almene,
imellem det Virkelige og det Væsentlige.

\emph{a) Det Sandselige og det Tænkelige.}

At det Sandselige er et Tankens Andet, en Reflexionens Skranke, maa Idealismen
naturligviis indrømme; men, idet Reflexionen viser, at Tanken just behøver denne
Skranke, bliver Andetheden, Skranken selv, jo dog alligevel tænkt, og det Sandselige,
det Negative, som Tanken tænker, idet den tænker sig selv, derved omsmeltet til Moment
i Begrebet. Denne Tvetydighed kan for en Deel tilhylles i den rene, objective Logik,
hvor „Materien“, dette Ulogiske, dog ved Dialektikens Hjælp tilsidst bliver ligesaa
logisk som „Formen“; men i „Naturphilosophien“, hvor Reflexionen støder an imod
Kjendsgjerningen, blottes Vildfarelsen.

„Ved den gamle almindelige Tanke, at ethvert Legeme bestaaer af fire Elementer, eller
den nyere paracelsiske, at det bestaaer af Mercur eller Flydenhed, Svovl eller Olie,
og Salt (Jacob Bøhm kaldte det den store Treenighed), og ved mange andre Tanker af
denne Art er for det Første Gjendrivelsen bleven let, idet man under hine Navne vilde
forstaae de enkelte, empiriske Stoffer, der nærmest betegnes ved disse Navne. Men det
er ikke at miskjende, at de meget væsentligere skulde indeholde og udtrykke
Begrebsbestemmelserne. Man maa derfor meget mere beundre
% [print: „algemeine“ for „allgemeine“ — reproduced as printed]
„\textit{die Gewaltsamkeit, mit welcher der Gedanke, der noch nicht frei war, in
solchen sinnlichen besondern Existenzen nur seine eigene Bestimmung und die algemeine
Bedeutung erkannte und festhielt; er darf darum auch nicht auf experimentirende Weise
widerlegt werden:}“ (Hegel, Naturphil. S. 271).

% ---- printed p.47 (PDF 55) ----
Denne Opfatning er ganske i Idealismens Aand, og at Hegel heller ikke vil have den
indskrænket til en Bedømmelse af hiin Tanke, „\textit{der noch nicht frei war}“, men
endog udtrykkelig anvendt til Fordeel for sin egen Udtydning af Phænomenerne, er
indlysende af den Maade, hvorpaa han selv t.\,Ex.\ indfører det „lunariske“ og det
„cometariske“ Princip; thi netop i dette „Lunariske“ og dette „Cometariske“ sees
„\textit{die Gewaltsamkeit, mit welcher der Gedanke in solchen sinnlichen besondern
Existenzen}“ kun erkjender og fastholder „\textit{seine eigene Bestimmung und die
allgemeine Bedeutung.}“

Grundvildfarelsen (πρῶτον ψεῦδος) ligger i Miskjendelsen af det Sandselige. Det
Sandselige er ingenlunde ophævet, fordi det er opfattet som Tankens Andet;
Sandselighedselementet er mere end en logisk Negation, det er for Tanken den reelle
Skranke, det i sig Utænkelige, som derfor skal opfattes og bestemmes i Forhold til den
umiddelbare Sands.

„Efter Momentet af dens „Væren-i-sig“ vilde Materien være det enkelte Punkt; efter
Momentet af dens „Væren udenfor sig“ vilde den nærmest være en Mængde hinanden
udelukkende Atomer. Men, idet disse ved at udelukke hinanden forholde sig til
hinanden, har Atomet ingen Virkelighed, og det Atomistiske saavelsom den absolute
Continuitet eller den uendelige Deelbarhed er kun efter Mulighed til i den.“ I denne
hegelske Opfatning seer det jo ved første Øiekast rigtignok ud, som om Tanken
gjennemtrængte Materien. Ved den logiske Modsætning af „\textit{Insichseyn}“ og
„\textit{Außersichseyn}“ og dennes dialektiske Behandling slipper man jo vistnok for
de Vanskeligheder, som Atomistiken, denne Speculationens Anstødssteen, ellers
medfører, naar den skal oplyses paa en de Sagkyndige tilfredsstillende Maade; men
Vanskeligheden er kun omgaaet, ikke hævet: idet Sandselighedselementet hævder sin Ret,
optager Forholdet mellem Materiens Discretion og Continuitet et Utænkeligt, et for

% ---- printed p.48 (PDF 56) ----
Reflexionen Uigjennemtrængeligt; thi det Sandselige er Tankens Skranke.

For endnu fra en anden Side at overbevise os om, hvor Utænkeligt det Sandselige er,
behøver man blot at kaste et Blik paa den herbartske Synsmaade. Med den objective
Idealisme har Herbarts Intellectual-Realisme nemlig det tilfælleds, at den, skjøndt
paa sin Viis, vil gjøre det Sandselige tænkeligt. Umiddelbart opfattet er det
Sandselige kun et Skin; men dette Skin forudsætter et i sig Værende. Dette Værende, en
selvstændig, paa sig selv hvilende Realitet, kan kun erkjendes af Tanken. Det ellers
materielle Atom er her forvandlet til et Tankens Atom. Af slige eiendommelige,
positive, uforanderlige Tanke-Atomer, Qvalia, gives der en stor Mangfoldighed.
„Ligesaa godt som et mathematisk Punkt kan“ — til Trods for \textit{principium
indiscernibilium}! — „falde sammen med et andet paa samme Sted, ligesaa godt kunne jo
rigtignok disse absolut positive Tanke-Atomer gjennemtrænge hinanden, og falde sammen
i eet og samme
% ERRATA (Rettelser, p.48): l.16 fra neden add closing quote after „Punkt;“, and
% l.15 fra neden remove the stray closing quote after „opløse.“ — the Herbart
% quotation closes at „…samme Punkt;“, and Nielsen's own clause („men slige
% Tanke-Atomer … lader sig ikke opløse.“) stands outside the quotation.
Punkt;“ men slige Tanke-Atomer ere da runde Fiirkanter: det Sandselige lader sig ikke
opløse.

\emph{b) Det Individuelle og det Almene.}

I „\textit{Phänomenologie des Geistes}“ blev „dette Her“, og „dette Nu“, der
uvilkaarlig betegnede det Punktuelle, af Reflexionen omdannet til „det Almene“, og det
Punktuelle ophævet. Men, idet der ved dette „Nu“! dette „Her“! dette „Der“! peges paa
et Sandseligt, ophæves igjen det Almene, og det Punktuelle træder i Kraft. „Dette
her“! er naturligviis ikke det samme som „dette der!“; det Almene er i Ordet, men
Ordet viser, ved denne Pegen, udover sig selv paa det Individuelle. Det er med Hensyn
hertil af Betydning, at Locke og Hume have undersøgt Begrebernes Oprindelse.

I det almene Begreb forsvinder Forestillingen; med Bestemmelsen af det Individuelle
vender den igjen tilbage:

% ---- printed p.49 (PDF 57) ----
at miskjende Forestillingen er at ringeagte Kjendsgjerningen, og miskjende det
Individuelle.

I Bearbeidelsen af Forestillingselementet ere to Factorer virksomme: den productive
Indbildning og den reflecterende Tænkning. Den uvilkaarlige Forestillen danner
Billedet i Overeensstemmelse med Gjenstanden; herved betinges den factiske
Overeensstemmelse imellem Forestillingskredsen og den givne Omverden. Den frie
Indbildningskraft tilveiebringer derimod Uoverensstemmelse mellem Virkeligheden og
Bevidstheden, idet den paa forskjellig Maade omdanner Forestillingerne. Spørge vi nu,
hvorfra den paa sine Begrebers Dialektik rastløs arbeidende Tænkning faaer sit
bøieligste, aandigste, og altsaa beqvemmeste Materiale: fra de factiske
Individualbestemmelser, eller de ved Indbildningen omdannede Forestillinger og frit
svævende Almeenbilleder? kan Svaret ikke være tvivlsomt.

I hvilken Grad den objective Idealisme har miskjendt Forestillingselementet ved at
ombytte Facticitetens Billeder med svævende Anskuelser, kan sees af mange Exempler.
„Den første qvalificerede Materie er Materien som den rene Identitet med sig, som
Eenhed af Reflexionen i sig; forsaavidt er den kun den første, selv endnu abstracte
Manifestation. Tilværende i Naturen er den Relationen til sig som selvstændig imod
Totalitetens andre Bestemmelser. Dette Materiens existerende, almene \emph{Selv} er
\emph{Lyset}, som Individualitet \emph{Stjernen}, og denne som Moment i en Totalitet,
\emph{Solen}.“ (Hegel Naturphil. S. 130).

Saalænge denne Anskuelse af Lysets Natur holder sig i reen Almindelighed fjern fra
alle endelige Facticitetsbestemmelser, kan man gjerne indrømme den en sindrig
Anvendelse af logiske Kategorier; men fordres derimod en Tilknytning til bestemte
Phænomener, vil den hele Opfatning vise sig saa huul og svævende, at her end ikke paa
noget Punkt er Mulighed til en Indtrængen i det Factiske.

% ---- printed p.50 (PDF 58) ----
Hvor Speculationen gaaer nærmere ind paa en Charakteriseren af givne Objecter, og
saaledes indlader sig med det Individuelle, bliver Misforstaaelsen endnu føleligere.
„Leerjorden er ligesom Kiselen det umiddelbare, eenfolde, uoplukkede Begreb, saaledes
det første differente Jordige — Brændbarhedens Mulighed\ldots, Talk eller Bitterjord
er Saltets Subject; deraf kommer Havets Bitterhed. Det er en Mellemsmag, der er bleven
til Ildprincip\ldots.“ (Naturphilosophie S. 416).

Istedetfor at udvikle Begreberne ved skarp Analyse af det Individuelle, og reducere
det Particulære saaledes, at det indholdsrige Almene kan finde sine Bestemmelser
bekræftede i det Givne, begaaer Idealismen det Misgreb at lade Almeenformerne bevæge
sig i en saadan Afstand fra Kjendsgjerningerne, at disse liig Phænomener i den fjerne
Horizont antage svævende Omrids, og, eftersom „Begrebet“ finder det tjenligt, blive
enten til Klipper, eller til Dyr, eller til Skyer.

\emph{c) Det Virkelige og det Væsentlige.}

At En kan have Ret i et Væsentligt og dog i Virkeligheden Uret, er ligesaa fatteligt,
som at Noget kan være sandt i Almindelighed, og dog i Anvendelsen paa det Enkelte blive
usandt. At Tænken er identisk med Væren, er væsentlig sandt; men at vor bevidste Tænken
og Tingenes endelige Existens skulde være identiske, er virkelig usandt; at det, der
existerer, maa udtrykke Existensens Kategorie, er væsentlig sandt; men, at det kun er
Existensens Kategorie, der existerer, er virkelig usandt. De, der tale om Væsenhed, og
de, der tale om det Virkelige, tale et forskjelligt Sprog.

„Universet er at tænke under Billedet af en Linie, i hvis Midtpunkt falder $A = A$, ved
hvis Ender falder paa den ene Side $\overset{+}{A} = B$, d.\,e.\ en Overgriben af det
Subjective, paa den anden Side $A = \overset{+}{B}$, d.\,e.\ en Overgriben af det
Objective, dog saaledes, at relativ Identitet finder Sted

% ---- printed p.51 (PDF 59) ----
selv i disse Extremer. Den ene Side er det Reale eller Naturen, den anden det Ideale.
Den reale Side udvikler sig efter tre Potenser, idet Potens betegner en bestemt
qvalitativ Differens af Subjectivitet og Objectivitet: den første Potens, $A$, er
Materien og Tyngdekraften — Objectivitetens største Overvægt; den anden Potens, $A^2$,
er Lyset, en indre — ligesom Tyngden en ydre — Naturanskuen. Lyset er den høiere
Yttring af det Subjective. Det er selve den absolute Identitet. Den tredie Potens er
Lysets og Tyngdekraftens fælleds Product, Organismen, $A^3$. Organismen er ligesaa
oprindelig som Materien.“

Dersom nu en streng Empiriker, en Mand, der ikke veed af nogen Speculation, med Eet kom
til at høre denne Schellings Naturforklaring; vilde han udentvivl ansee det Hele for et
Hjernespind. Hermed er imidlertid ikke beviist, at det, der udsiges i hine speculative
Sætninger, er i sig meningsløst; hermed er kun givet, at Empirikeren ikke forstaaer
Sproget. Dersom En, der vilde forklare Verdens Opholdelse, pludselig greb Lagrange's
% [print: a small dot precedes „da“ in the numerator — possibly a differential dot]
Formel: $\int \dfrac{da}{dt}\,dt = \Sigma\,\dfrac{p}{in-i'n'}\,\sin(in-i'n'+q)$ med den
Forsikkring, at Solsystemet vedligeholdes alene af den Grund, at $i$ og $i'$ aldrig
samtidig kunne være lige med $0$, saa at Nævneren $in-i'n'$ ingensinde kan blive $0$,
fordi $n$ og $n'$ ere incommensurable; da vilde vel endogsaa de, der ellers tænke
fornuftig over virkelige Ting, men ikke have Begreb om Mathematik, bedømme denne
mathematiske Forklaring ligesaa strengt, som Empirikeren hiin idealistiske.

At Ukyndighed og Mangel paa philosophisk Sands forene sig om at nedrive Idealismen,
betyder Intet; for at forstaae det philosophiske Sprog, maa man ligesaa vist studere
Philosophie, som man, for at forstaae Mathematikeren, maa studere Mathematik.

% ---- printed p.52 (PDF 60) ----
Men selv, naar vi søge at forstaae Philosophien philosophisk, kan det ikke nægtes, at
den objective Idealisme har, med sin dybsindige Væsenserkjendelse, forsyndet sig mod
Virkeligheden paa dobbelt Maade, deels ved at abstrahere fra virkelige Udgangspunkter,
deels ved at ville tvinge et virkeligt Indhold ind i uvirkelige Begreber. „Den
newtoniske Theorie, hvorefter Lyset skal udbrede sig i Linier, eller Bølgetheorien,
hvorefter det skal udbrede sig bølgeformig, som den eulerske Æther eller som Lydens
Zittren, ere“, siger Hegel (Naturphil. S. 141), „materielle Forestillinger, „\textit{die
für die Erkenntniß des Lichts nichts nutzen.}““ Herved er da abstraheret fra alle
virkelige Udgangspunkter, og enhver controllerende Undersøgelse gjort umulig.

Istedetfor at sætte sig ind i de virkelige Undersøgelsers Gang, og ved Begrebets
dialektiske Overlegenhed føre Empirikeren i Kraft af hans egne vitterlige
Forudsætninger til Erkjendelse af Sagens Væsen, har denne Idealisme ikke alene med
vilkaarlig Opfatning af saadanne Phænomener som Lysbrydning og Farver, men fra Først
til Sidst, villet corrigere Empiriens Udtryk og Fremgangsmaade, uden dog ret at fatte,
end sige, gjennemskue Empiriens Problemer.

\begin{center}{\bfseries \S\,10.}\end{center}

Naar der handles om en Erkjendelse af det Virkelige, maa Erfaringslæren have den
afgjørende Stemme. Den hume'ske Skepticisme gjør i sin Erfaringsiver sand Erfaring
umulig. Ved at bevise de aprioriske Formers, om end kun subjective, Almeengyldighed
bringer Kant nu vistnok Philosophien i nærmere Forstaaelse med Empirien. Men,
endskjøndt Erfaringslæren til en vis Grad gjerne indrømmer Kant Umuligheden af at
erkjende Tingenes Indre, er den dog langt fra at opfatte Tid og Rum som blot værende i
vor Bevidsthed. Naar Astronomen taler om Maanens Afstand fra Jorden, anskuer han dog
visselig denne Afstand fra Jorden, ikke som en blot subjectiv Opfatningsmaade,

% ---- printed p.53 (PDF 61) ----
men som et virkeligt, objectivt Rumsforhold; naar det hedder, at Neptun perturberer
Uranus, søges Vexelvirkningen imellem Planeterne dog vistnok ikke i vor Forstand
alene, men i Tingenes Natur. Den Skilsmisse, Kant gjør imellem
Erfaringsphænomenerne og Tingene selv, er en Skoletheorie, der ikke har mindste
Betydning i Erfaringslæren. Ved at hævde de aprioriske Formers objective Gyldighed har
Hegel derimod, trods alle hans Anstød mod Virkeligheden, i denne Henseende paaviist
Philosophiens og Empiriens væsentlige Overeensstemmelse.

Skal denne Overeensstemmelse nu ogsaa i Begrebets Form gjennemføres fra Virkelighedens
Synspunkt, maae de rationelle og dog empiriske Bindeled imellem de aprioriske Former
og de umiddelbare Kjendsgjerninger lade sig paavise. Det er ikke nok, at et
almeengyldigt Forhold mellem Grund og Følge, Betingelser og det Betingede, Aarsag og
Virkning, o. s. fr. anerkjendes a priori; Opgaven er, at forvisse sig om de virkelige
Forhold.

Forholdet mellem Aarsag og Virkning er et Nødvendighedsforhold, der har sit blivende,
almene Udtryk i Loven. Aarsagen for sig er ikke en Lov, skjøndt det er en apriorisk
Lov, at Alt maa have en Aarsag; Virkningen for sig er naturligviis heller ikke Loven;
men Loven er, at Virkningen maa indtræde, naar Aarsagen finder Sted. Paa disse
aprioriske Forudsætninger beroe Iagttagelse og Experiment.

Naar en virkelig Aarsag, et virkeligt Grundforhold, en virkelig Lov, skal udfindes,
maa man begynde med en Antagelse. At Jordens Tiltrækning er Aarsag i Stenens Fald, at
en Vædskes Fordampning er Aarsag til Kulde, at ethvert Legeme har sit bestemte
Smeltepunkt, at Ebbe og Flod skyldes Solens og Maanens Indvirkning, alle slige
Antagelser vise oprindelig tilbage til et subjectivt Skjøn, til Erfaringsindtryk, som,
om de end gjentages nok saa ofte, dog umulig kunne bevise det Nødvendige, eller
godtgjøre det Almene.

% ---- printed p.54 (PDF 62) ----
For da at indsee, hvorvidt Erfaringslæren ved sine „Iagttagelser“ og „Forsøg“ er
istand til at sikkre sig det Almeengyldige, betragte vi særlig: Abstractionen,
Inductionen og den virkelige Lov.

\emph{a) Abstractionen.}

Det Abstraherede minder om det, hvoraf det er abstraheret; saaledes minder det rene
Rum om Materien, og fordrer, saa at sige, sit tabte Indhold tilbage. Til Rummets
Discretion svare Materiens Atomer, til Rummets Continuitet Materiens Masse. De
Almeenbestemmelser, der ligge i Geometriens Linier, Flader og Legemer, have tillige
Gyldighed, forsaavidt de udsiges om det Materielle. At den rette Linie er den korteste
Vei imellem to Punkter, at to rette Linier kun kunne skjære hinanden i eet Punkt o. s.
v., ere Sætninger, der kunne abstraheres af det Materielle, og anvendes paa det
Materielle. Forsaavidt Materien opfattes som det Abstrakte, træffer Erfaringsmomentet
her sammen med det Aprioriske.

Men selv, hvor Erfaringsmomentet træder stærkere frem, som, naar der tillægges den
rumopfyldende Materie en Modstandskraft, en Udvidekraft, Tyngde o. s. v., ere vi
ingenlunde, hvad Hume antager, blot henviste til Ideeassociationen og Vanen. „Alle
Legemer have Tyngde, alle Legemer have Sammenhængskraft, alle Legemer kunne deles“: en
simpel Benegtelse af disse Sætningers Almeengyldighed har Intet at betyde; de
understøttes af Millioner Erfaringsindtryk. Vel ere disse Sætninger i deres
Almindelighed abstracte, men dette Abstracte er en Reflex af det Concrete; Den, der
altsaa vil benegte Almeengyldigheden, maa udtrykkelig anføre Kjendsgjerninger, og
paatage sig \textit{onus probandi}.

Hvad den uophørlige Vexelbestemmelse af Erfaringselementet og det Aprioriske betyder,
viser Physiken. „Det er vigtigt, at de meest enkelte mathematiske Sandheders
Anvendelse i Naturlæren omhyggelig retfærdiggjøres, paa det

% ---- printed p.55 (PDF 63) ----
at ingen Tvivl levnes om Naturens Overeensstemmelse med Tænkningen“ (H. C. Ørsted); og
det er ligeledes vigtigt, at Naturlærens Sætninger opfattes mathematisk; thi det, der
især fremhjælper de strenge Undersøgelser, i hvilke Conseqvensen bliver kjendelig paa
Forudsætningen, og denne igjen paa hiin, er just den exacte Form, hvori Mathematiken
(Jfr. Philosophie og Mathematik S. 12—14) opfatter de paagjældende Functioner.

Forskjellen imellem den abstracte, philosophiske og den mathematisk-physiske Atomlære
bestaaer da ogsaa væsentlig deri, at denne strax griber Virkeligheden i dens
Punktualitet, medens hiin aldrig kommer ud over det eensidig Almene. Vistnok begynder
Mathematiken ogsaa med Abstractioner, men den viser tillige, hvorledes Abstractionerne
kunne integreres. Mathematikeren taler ikke blot om Atomer og Atomernes Forening i
Moleculer, men formulerer derhos Udtrykket saaledes, at det netop egner sig til en
Undersøgelse af det Virkelige. Moleculernes Antal være t. Ex. $m$, d. e. Massen, $v$
Massens Rumfang, der indeslutter et Antal Moleculer; man har da i Formlen:
$\dfrac{m}{v} = \varrho$, et exact Udtryk for Tætheden, og heri en Forudsætning for
Bestemmelsen af de forskjellige Legemers Vægtfylde.

Indeslutter nu $v$ et Antal af $n$ Moleculer, faaer man $\dfrac{v}{n} = e^3$, og $e^3$
altsaa som Beregningsudtryk for Moleculen. Bliver i et Legeme endvidere $m$, den
endelige Deel af Massen, uendelig lille, altsaa: $dm$, saa erholder man
$\dfrac{dm}{dv} = \varrho$, hvorved da ogsaa Delenes endelige Sum, $\Sigma$, gaaer over
til $\int$, den uendelige Sum; man har altsaa for det \emph{homogene} Legeme, idet
$\varrho$ er en Function af de forskjellige Steder i Massen, Formlen: $\varrho\int dv$,
og Formlen: $\int \varrho\, dv$ for det \emph{heterogene},

% ---- printed p.56 (PDF 64) ----
og heri et Grundbegreb, hvoraf utallige med Erfaringen stemmende Specificationer lade
sig udlede.

Den mathematiske Punktualitet muliggjør en Virkelighedsanalyse, som i Skarphed kan
maale sig med Dialektikens Væsensanalyse.

\emph{b) Inductionen.}

Ved Abstractionen sees der bort fra det enkelte Tilfælde; ved Inductionen faaer det
Enkelte igjen sin Betydning. Hvorledes de enkelte Led lede hen til det Almene,
forsaavidt dette har Tilværelse i dem; hvorledes Delene ved deres Eiendommelighed
bestemme det Hele, saa at Heelhedens Væsen finder sit fuldstændige Udtryk i Delenes
Virkelighed, viser Inductionen. „Det er det Mærkværdige ved denne Slutningsart, at
Slutningssætningen aldrig er en singulær eller særegen, men altid en \emph{almindelig}
Dom. Den kategoriske Slutning slutter fra en Regel til et Tilfælde; den disjunctive
derimod fra Tilfældene til Reglen. Hist er Reglen det Forudsatte, her det ved Slutning
Opnaaede. Kategorisk kan man kun deelviist, hypothetisk kun benegtende (\textit{modo
tollente}), slutte fra det Særlige til det Almene. Den disjunctive Slutningsmaade er
den eneste, der tilsteder en bekræftende og tillige almindelig Slutning fra det
Særlige til det Almene. Deri ligger Inductionens Tryllemagt: den lader os af
Iagttagelsernes og Kjendsgjerningernes Sammenstilling erkjende Loven.“ (\textit{Die
Theorie der Induction von E. F. Apelt. Leipzig 1854. S. 18}).

Med Hensyn paa Inductionens Uddannelse og Anvendelse kunne Sokrates, Aristoteles og
Keppler siges at betegne Vendepunkterne.

„To Ting,“ siger Aristoteles, „kan Ingen med Billighed gjøre Sokrates stridig,
Inductionen og Bestemmelsen af de almene Begreber.“ Hvorledes Sokrates gik frem, idet
han ved Betragtning af en Mængde enkelte Tilfælde og

% ---- printed p.57 (PDF 65) ----
Exempler søgte at tydeliggjøre de ethiske Begreber, er bekjendt. Men, uagtet denne
sokratiske ἐπαγωγή vistnok er en Fremgang fra det Enkelte og Særlige til det Almene,
er den dog, hvad Apelt ogsaa bemærker, ikke det, hvad vi nu forstaae ved Induction,
men snarere en analyserende Abstractionskunst.

Som Logiker og Naturphilosoph er Aristoteles opfordret til at give Inductionen en
bestemtere Form, og sætter da ἐπαγωγή deels imod συλλογισμός, deels imod ἐπιστήμη. I
første Tilfælde er ἐπαγωγή Slutningen fra det Særlige til det Almene, fra Delene til
det Hele, fra Tilfældene til Reglen; i det andet Tilfælde angiver ἐπαγωγή ogsaa
Begyndelsen til det Almene (ἀρχή ἐστι καὶ τοῦ καθόλου); medens ἐπιστήμη kun uddrager
Fornuftslutninger fra det Almene. Ἐπαγωγή fører til Principet, ἐπιστήμη udvikler sit
Indhold af Principet. Hertil bemærker nu Apelt: „Inductionen (ἐπαγωγή i Modsætning til
συλλογισμός) fører aldrig til Principet, man forudsætter det allerede; den fører meget
mere til en Lov, der knytter en bestemt Kreds af Phænomener til et Princip. Den
regressive, epagogiske Tankegang, der fører til Principet (ἐπαγωγή i Modsætning til
ἐπιστήμη), er ingen Slutningsforbindelse, men en Abstraction. Denne Forvexling af
Induction og Abstraction er i Philosophiens Historie bleven staaende indtil den nyere
Tid.“

Som den, der i den nyere Tid rigtig har opfattet og videnskabelig anvendt
Inductionen, nævnes Keppler; thi Baco af Verulam fremhæver kun dens Betydning, men
kjender ikke dens Anvendelse. Kepplers Opdagelse af Marsbanen skete ved Anvendelse af
Inductionen; denne Opgaves Løsning behandler Apelt med stor Skarpsindighed og
Sagkundskab. „Kunsten bestaaer i, af en Række numeriske Værdier for to sammenhørende
foranderlige Størrelser at udfinde Loven for deres gjensidige Afhængighed af hinanden
d. e. Functionsformen. Heri ligger Opgavens store Vanskelighed.“

% ---- printed p.58 (PDF 66) ----
Den logiske Form for den kepplerske Induction er denne: Marsbanen indeholder i sig
Marsstederne:
\[ O_1,\ O_2,\ O_3,\ \ldots O_n. \]
Paa Stedet $O_1$, d. e. til den Tid, da Mars's sande Anomalie er $v_1$, er Størrelsen
af Radius Vector $r_1$
\[
\begin{array}{ccccccccc}
\text{,,}&\text{,,}&\text{,,}& v_2 &\text{,,}&\text{,,}&\text{,,}&\text{,,}& r_2\\[2pt]
\text{,,}&\text{,,}&\text{,,}& v_3 &\text{,,}&\text{,,}&\text{,,}&\text{,,}& r_3\\[2pt]
\multicolumn{9}{c}{\vdots}\\[2pt]
\text{,,}&\text{,,}&\text{,,}& v_n &\text{,,}&\text{,,}&\text{,,}&\text{,,}& r_n\\
\end{array}
\]
Marsbanen er en Curve, hvis Polarligning er: $r = \dfrac{a(1-e^2)}{1-e\cos v}$. Men
$r = \dfrac{a(1-e^2)}{1-e\cos v}$ er Ellipsens Polarligning. Marsbanen er altsaa en
Ellipse, i hvis ene Brændpunkt Solen staaer.“

\emph{c) Den virkelige Lov.}

Hvorledes Inductionen, begyndende med Iagttagelsen af det Enkelte, fører fra det
Enkelte til det Totale, idet den fører gjennem Kjendsgjerningerne til det Almene, og
derved til Erkjendelsen af den virkelige Lov, kan sees af det anførte Exempel.

At Marsbanen virkelig er en Ellipse, er ikke beviist ved at iagttage alle enkelte
Punkter i Banen; en saadan Iagttagelse er umulig, da Banen indeholder uendelig mange
Punkter. Ikke desto mindre giver Inductionen tilstrækkelig Vished. Denne Vished er
ikke overfladisk støttet paa en Sandsynlighedsberegning, der kun vil nærme sig til at
eliminere Tilfældet: nei, her er en væsentlig Slutning, der begrunder en Læresætning.

For at fatte Inductionens Beviiskraft, maa man vel lægge Mærke til, at den jo ikke
skal afgjøre, hvorvidt det paagjældende Phænomen, her Planetens Bevægelse, er under-

% ---- printed p.59 (PDF 67) ----
kastet en Lov eller ikke (thi det er a priori afgjort); men derimod godtgjøre,
\emph{hvilken} Lov her i Virkeligheden har fundet sit Udtryk.

Imellem den aprioriske Forudsætning, at her er en Lov, og det empiriske Spørgsmaal:
hvilken Lov? er Afstanden saa stor, at Apelt med Grund anbefaler „heuristiske
Maximer,“ for at lede Iagttageren paa rette Spor. „Maximerne“ selv ere i denne
Sammenhæng uden særegen Interesse; hvorimod følgende instructive Anviisning (Anf. Skr.
S. 50) ikke bør forbigaaes.

„Ved den kategoriske Slutning ere de tre Termini hvergang givne, ved Inductionen
derimod maa Overbegrebet ofte først søges. Naar man skal søge Noget, maa man først
vide, \emph{hvad}, og \emph{hvorledes}, og \emph{hvor} man skal søge det. Inductionen
søger den Lov, af hvilken en bestemt Kreds af Phænomener afhænger. Man maa da aabenbart
i Forveien vide, at Lovens Nødvendighed, og ikke Tilfældet (Lovløshed), hersker i denne
Kreds. Saaledes er t. Ex. Stjernernes Løb bunden til en Lov, Antallet af de jordiske
Grundstoffer derimod ikke. Men den blotte Forudsætning af Lovmæssighed er endnu ikke
tilstrækkelig; man maa ogsaa have en vis Forestilling om den Lovs Beskaffenhed, man vil
søge. Da Keppler søgte Marsbanens Figur, da vidste han allerede forud for sin
Induction d. e. a priori, at Planeterne ikke gaae deres Vei gjennem Himmelrummet efter
Lune og Tilfælde, men at de omkring Solen beskrive en Bane, der lader sig betegne efter
Geometriens Love. Man maa, naar man ikke vil forfeile Veien til Sandhed, vide, om den
postulerede Lovmæssighed blot er geometrisk, eller mechanisk, eller teleologisk.
Saaledes søge endnu den Dag idag Nogle teleologiske Forklaringsgrunde for Livets
Phænomener i Physiologien, Andre physikalske (mechaniske og chemiske), fordi man endnu
ikke er enig angaaende de paa dette Omraade herskende Love.“

% ---- printed p.60 (PDF 68) ----
At Galilei, der fremhæver Grundsætningen om Materiens Inertie, repræsenterer
Abstractionen, ligesom Keppler Inductionen, og at begge Aandsretninger ere forenede i
Newton (\textit{Philosophiæ naturalis principia mathematica}), er en sindrig
Bemærkning af Apelt; hvorimod den speculative Udtydning, Hegel (Naturphil. S. 94—124)
giver af Faldbevægelsen og de kepplerske Love, er ligesaa holdningsløs, som hans Dom om
Newtons Gravitationslære er umotiveret.

Hverken Sensualisten eller Idealisten er istand til at erkjende den virkelige Lov;
hiin ikke, fordi han negter det Aprioriske; denne ikke, fordi han overseer
Kjendsgjerningen.

\begin{center}{\bfseries \S\,11.}\end{center}

At en Erkjendelse af de virkelige Love er en ufravigelig Betingelse for at erkjende de
tilværende Ting, vil neppe Nogen drage i Tvivl. Men det Spørgsmaal, hvorledes en
Erkjendelse af det Virkelige forholder sig til en Erkjendelse af Væsenet, angiver dog i
denne Sammenhæng et nyt Synspunkt, hvorfra Modsætningen mellem Kant og Hegel nærmere
kan bedømmes. Ifølge Kant vilde det da altsaa være muligt at drive det endogsaa meget
vidt i Erkjendelse af det Virkelige, uden derved i mindste Maade at komme Tingenes
Væsen nærmere; ifølge Hegel vilde det omvendt være muligt, at erkjende det evige
Væsen, uden at bekymre sig meget om de virkelige Tings Forskjellighed. Hvorledes dette
er at forstaae, vil først blive os klart, efterhaanden som vi nøiere undersøge
Forholdet imellem Tingene og Lovene, Tingene og deres Egenskaber, den væsentlige Viden
og de virkelige Objecter.

\emph{a) Tingene og Lovene.}

Selv naar man indrømmer det i Lovene Almeengyldige, gjør det dog en væsentlig Forskjel,
om man opfatter disse Love som væsenløse Regler uden Undtagelser, eller seer i deres
Opfyldelse en Yttring af Tingenes Natur.

% ---- printed p.61 (PDF 69) ----
I Erfaringslæren tales et dobbelt Sprog. Paa den ene Side hedder det: „Vi vide ikke,
hvad Materien er i sig selv; Kraften er os et Ubekjendt, vi kjende kun dens Yttringer;
Lysets Væsen er os ubekjendt, vi kjende kun Phænomenet; vi kjende ikke Varmens Aarsag,
vi kjende kun dens Virkninger.“ Paa den anden Side hedder det igjen: „det er ikke ved
Speculationer og Theorier, men ved Iagttagelser og Forsøg, at vi trænge ind i Tingenes
Væsen, og lære de skjulte Kræfter at kjende. Vi erkjende Naturens Love; disse Love
udgjøre det Almindelige i Mangfoldigheden, det Evige i Omskiftningerne; disse Naturlove
ere sande Naturtanker, saa at det, vi erkjende i Naturen, ikke er noget Fremmed, men
det Samme som det, der lever i os selv.“

Disse to Synsmaader staae i aabenbar Strid med hinanden. Er det sandt, at vi kun
erkjende Phænomenet, ikke Væsenet, kun Yttringen, ikke Kraften, kun Virkningen, ikke
Aarsagen; da er det usandt, at Tingenes Love udtrykke Tingenes Væsen. Er det sandt, at
Naturlovene, som Aabenbarelser af en evig virksom Fornuft, udtrykke Tingenes Væsen, da
er Umuligheden af at erkjende Væsenet, Kraften, Aarsagen, kun en Fiction.

Saadanne Spørgsmaal som: hvad er Rummet i og for sig? hvad er Materie i og for sig?
hvad er Kraft i og for sig? kunne forraade megen Uklarhed i det Aprioriske. Materien i
og for sig, den rene Materie, er ligesaa vist en Abstraction, som det rene Rum. Man kan
ikke spørge, \emph{hvorfra} eller \emph{hvorledes} Materien er kommen ind i Rummet, thi
Materiens og Rummets Væren er i Principet den samme. Her er en væsentlig Identitet i en
virkelig Sondring; som det sandselige Indhold danner Materien en virkelig Modsætning
til den udvortes Form. Det er saa langt fra, at Anskuelsen af et tomt Rum skulde være
blot subjectiv, at den tvertimod viser Rummets objective Gyldighed som virkelig Afstand
imellem to fra hinanden fjernede Legemer.

% ---- printed p.62 (PDF 70) ----
Saalidt som Materien kan være i sig, uden at være i Rummet, ligesaalidt kan Kraften
være i sig, uden at være i Materien. Man nævne hvilken Kraft, man vil, den vil altid
være bestemt ved et eiendommeligt Forhold, hvori den ene Materie kan befinde sig til
den anden. Hvad er Sammenhængskraft? Et Udtryk for den Modstand, et Legeme gjør imod
Delenes Adskillelse. Her er da i Legemet ikke først en Mængde Dele i og for sig, og
derhos en Kraft i og for sig; nei, Sammenhængskraften er en Væsensyttring af de
materielle Dele; den kraftige Sammenhæng udtrykker, at Delene kun ere Noget i sig
derved, at de ere Noget for hverandre indbyrdes. Som det gaaer med
Sammenhængskraften, gaaer det ogsaa med Tyngdekraften, med de magnetiske, elektriske,
chemiske Kræfter o. s. v.

Tænkes al Kraft borte, er Materien kun en tom Abstraction; tænkes Materien borte, er
Kraften en Abstraction. Kan man ikke sige, at det er Materien, der oprindelig giver
Kraften Tilværelse, efterdi Materien, for selv at være til, maa forudsætte Kraften;
saa kan man da paa den anden Side heller ikke antage, at Kraften giver Materien
Tilværelse, da enhver Kraftyttring jo maa forudsætte en værende Materie, hvorfra den
kan udgaae. Man behøver ikke at gruble over det Spørgsmaal, hvorledes Kraften
oprindelig er kommen ind i Materien, da Kraftens og Materiens Væren i Principet er den
samme.

Hvad Chemiens saakaldte Grundstoffer angaaer, er det ligesaa modsigende at ville
bestemme deres Natur i og for sig, som at ville udspeculere, hvad den chemiske
Tiltrækning er allene i og for sig. Ilten, Brinten, Svovlet, slige Grundstoffer ere
kun i og for sig, hvad de udtrykke i deres gjensidige Forhold og Forbindelser. Det er
ikke en Mangel ved vor Erkjendelse, at vi ikke kunne indsee, hvad den chemiske Kraft
er ubetinget i og for sig; thi den er, som et Ubetinget i og for sig, et Intet, en
Abstraction, en Indbildning.

% ---- printed p.63 (PDF 71) ----
At Kant i sin Fornuftkritik maatte støde paa en uforklarlig „Ting i sig“, er
begribeligt; men at ville, efter Alt, hvad den hegelske Dialektik har oplyst, endnu
fastholde en saadan „Ting i sig“, og derved mene at undgaae de dialektiske
Modsigelser, er kun at skuffe sig selv: en Ting i sig, en Ting bag ved alle
Betingelser, er en lovløs Ting; thi Tingene ere kun under Loven, forsaavidt de med
indre Nødvendighed ere for hverandre indbyrdes, og betinge hverandre gjensidig; en
Ting i sig, en ubetinget Ting, er ikke blot en Modsigelse, men, naar den fastholdes,
en Abstractionsforhærdelse, en Uting.

Er det nu indlysende, at Tingene kun i deres indbyrdes Vexelvirkning lægge deres Væsen
for Dagen, og er Loven det formelle Nødvendighedsudtryk, som for denne Vexelvirkning
angiver de bestemte Forudsætningers bestemte Følger; da er det vel ogsaa uimodsigeligt,
at Loven er mere end en almindelig Regel uden Væsenhed, at en virkelig Loverkjendelse
tillige er en Væsenserkjendelse, at de forskjellige Love angive forskjellige Sider af
Tingenes Væsen.

\emph{b) Tingene og deres Egenskaber.}

Som Væsenet kjendes paa dets Phænomen, maa Tingen kjendes paa dens Egenskaber. Tingen
er Egenskabernes existerende Eenhed, og Egenskaberne ere igjen Tingenes indre
Bestemmelser i tilværende Fleerhed. Men, uagtet Egenskaben er det for Tingen Egne, kan
den ene Tings Egenskab dog i Virkeligheden kun blive kjendelig ved dens endelige
Forhold til den anden Tings Egenskab; til Sagens Belysning kræves, at Egenskabernes
Vexelvirkning sees fra den virkelige Side.

Den Forskjel, der i Virkeligheden er imellem det Umiddelbare og det gjennem mange
Mellemled Formidlede, imellem det Partielle, der opfattes enkeltviis, og det Totale,
hvis indre Sammenhæng kræver alsidig, nøiagtig, udholdende Forskning, bevirker, at
Afstanden mellem en forberedende

% ---- printed p.64 (PDF 72) ----
Erkjendelse af Phænomenerne og en afsluttende Væsenserkjendelse træder tydelig frem.

Ved blot at holde sig til de optiske Phænomener, faaer man endnu ingen Forestilling om
Lysets chemiske Virkninger; af den blotte Varmefornemmelse kan man ikke forklare sig,
at Varmen udvider Legemerne. Hvorledes gaaer det nu til, at Lyset, den samme Aarsag,
kan have saa mange og høist forskjellige Virkninger? Kan det Samme, der er lysende,
ogsaa være varmende, eller gives der maaskee et eget Varmestof? Kan Lys forvandles til
Varme, og Varme til Lys, hvori bestaaer da Forvandlingen? Naar slige Spørgsmaal sættes
i Bevægelse, indsees det, at Den, der virkelig vil en sand Erkjendelse af Lysets og
Varmens Væsen, neppe, som Hegel, kan oversee „den newtoniske Theorie“ og „den eulerske
Æther.“

Hvilken Betydning de af Hegel forsmaaede Hypotheser maae have endog for
Væsenserkjendelsen, sees allerede deraf, at det gjælder en Undersøgelse af Materiens
egentlige Discretion (\textit{Fechner: Atomenlehre. Leipzig 1855. S. 18—32}), naar det
gjælder en Undersøgelse af Varmens Identitet med Lyset.

At Emanationslæren forudsætter Atomer af endelig Størrelse, forstaaer sig; men heller
ikke Undulationstheorien kan gjennemføres, uden under en lignende Forudsætning. Under
Straalebrydningen spreder den brudte hvide Straale sig i en smal Farvevifte; dette
Phænomen kan kun bringes i mathematisk Overeensstemmelse med Svingningshypothesen, naar
Ætherdelene antages at være sondrede ved endelig Discretion. Den smalle Farvevifte
beroer nemlig paa, at de forskjellige Straaler have forskjellig Brydbarhed, og at
Ætherdelenes Afstand ikke er forsvindende i Sammenligning med en Lysbølges Brede. Den
samme Synsmaade, at nemlig Ætherdelene ere i endelig Sondring, stadfæstes endvidere,
naar det sædvanlige Lysphænomen sammenlignes med Polarisationsphænomenet. I en
sædvanlig Lysstraale have

% ---- printed p.65 (PDF 73) ----
Svingningerne alle mulige Retninger, i en polariseret Lysstraale ere de parallele, i
begge ere de transversale. Men transversale Svingningers Fortsættelse er kun mulig
under Forudsætning af Ætherens Discretion.

Som nu Lysphænomenet i alle Henseender kan tydes efter Svingningshypothesen, kan
Varmevirksomheden det ogsaa. Vel er Varmens Udbredelse ved Berøring i den Grad
forskjellig fra Straalingen, at det ved første Øiekast skulde synes umuligt at forklare
saa forskjelligartede Phænomener af samme almindelige Forudsætning; men Fourier har
viist, at Phænomenerne lade sig forene, naar man nemlig antager, at Legemerne ikke ere
Continua, men bestaae af sondrede Smaadele. Legemernes Ledningsevne bliver da afhængig
af den ledede Varmes Svingningshastighed, altsaa af Temperaturen: Varmen og det
Opvarmede charakterisere hinanden gjensidig.

At Varmestraalerne virke stærkest, naar de falde lodret paa Legemets Overflade, og
derimod svagere, naar de falde i skraa Retning, er uforklarligt, saalænge man vedbliver
at gaae ud fra Materiens Continuitet; forudsætter man derimod en endelig Discretion, er
Phænomenet ikke blot forklarligt, men endogsaa beregneligt, idet Straalevirkningen
svækkes efter Loven for Sinus.

I Chemien og Krystallæren finder ikke alene den herved paapegede Atomlære særegen
Bekræftelse, men Lys- og Varmelæren gjennemgaae med det Samme en ny Prøvelse. See vi da
paa Sagen i dens Heelhed, kan Lysets Væsen kun erkjendes af Phænomenernes Totalitet;
men Totaliteten har et uendeligt Omfang: den ene Totalitet griber ind i den anden. At
undersøge Lyset som det rene Lys i og for sig er umuligt; Lysets Egenskaber kjendes kun
paa de Tings Egenskaber, der paavirkes af Lyset; det Samme gjælder om Varmen, om
Magnetismen, om Elektriciteten o. s. fr. Disse „Potenser“ ere ikke blot, hvad de ere,
hver for sig i deres Egenskabers Vexelvirkning med de paagjældende Tings Egen-

% ---- printed p.66 (PDF 74) ----
skaber; men de ere tillige for sig, hvad de ere i deres Vexelvirkning med hinanden
indbyrdes.

\emph{c) Den væsentlige Viden og de virkelige Objecter.}

I Forhold til den almene Væsenhed har Hegel Ret imod Kant; men Virkeligheden har Ret
imod Hegel. Istedetfor at erkjende det virkelige Forhold imellem Ting og Egenskaber,
har Hegel allerede i „\textit{Phänomenologie des Geistes}“ anlagt Begrebsudviklingen
saaledes, at der i Reflexionen skrides frem fra Tingens lavere til Objectivitetens
høiere Kategorie. Men ved denne logiske Fremadskriden er det saa langt fra, at
Virkelighedsproblemet bliver kjendeligt, at det meget mere tilhylles. Det Almene er det
Objective, men det Almene er identisk med Tanken; det Objective er i dets eiendommelig
bestemte Objectivitet et Object, men Objectiviteten er identisk med sit Begreb: Objectet
maa objectiveres af Subjectet. Ad denne Vei kommer man aldrig ud over den almindelige
Identitet af Tænken og Væren.

Lade vi derimod det Existerende, Tingen med dens Egenskaber, stille sig i sandselig
bestemt Modsætning til den subjective Reflexion, træder Virkelighedsproblemet i Kraft.
I dette Punkt seer Locke klarere end Hegel.

Strengt taget, kan „Substratets“ Identitet med sig i dets mange Egenskaber og
Accidenser ikke for noget virkeligt Object hævdes med apriorisk Nødvendighed. Hvad er
dette? „Sølv! man seer det strax af den eiendommelige Glands.“ — Men Glandsen kan jo
blænde; et andet Legeme kunde jo muligviis have en lignende Hvidhed. Af en eller to
Egenskaber kan et Object ikke kjendes; vi maae da for Sølvets Vedkommende prøve en
Mængde Egenskaber: Bøielighed, Strækbarhed, Vægtfylde, Krystalform, Smeltepunkt,
Tiltrækning til Chlor, til Svovl, til Ilt o. s. v. Hver Egenskab maa prøves
tilstrækkelig, og et tilstrækkeligt Antal Egenskaber maa prøves: her er en
Virkelighedsopgave i dens Uendelighed.

% ---- printed p.67 (PDF 75) ----
Det kan indrømmes den hegelske Begrebsanalyse, at den imod Hume saavelsom imod Kant og
Herbart tydeliggjør den dialektiske Identitet af Tingen og dens Egenskaber, af
Substansen og dens Accidenser, af Begrebet og Objectet i Almindelighed, og derved a
priori sikkrer Erfaringslæren et objectivt Grundlag: at Sølvsubstansen er i væsentlig
Identitet med Sølvets Egenskaber, kan ikke betvivles uden at Virkelighedsloven maa
krænkes, og Empirien modsige sig selv; men \emph{hvilke} Egenskaber her ere forenede,
og \emph{hvorledes} de ere forenede i denne virkelige Substans, maa prøves ved
Iagttagelse og Forsøg; Realbegrebet henter ethvert Moment fra Empirien.

I Kraft af den empiriske Forsknings indre oprindelige Sammenhæng med det Aprioriske,
gives der altsaa en væsentlig Viden; men ifølge Tænkningens og Sandsningens
Ueensartethed er Realerkjendelsen underlagt endelige Betingelser, og al vor Viden
begrændset. Videns Grændse er en dobbelt: en \emph{relativ}, forsaavidt her altid er
Noget, vi endnu ikke vide, men som vi dog med Tiden kunne faae at vide; en
\emph{absolut}, forsaavidt en endelig Adskillelse af Bevidsthedens og Virkelighedens
Indhold er begrundet i vor Organisation: her er altsaa Noget, vi i denne Tilværelse
aldrig faae at vide.

\begin{center}\rule{0.35\linewidth}{0.4pt}\end{center}

% ---- printed p.68 (PDF 76): II. Videnskabslære ----
\begin{center}
  {\large\bfseries II.}\par\vspace{4pt}
  {\Large\bfseries Videnskabslære.}\par\vspace{4pt}
  \rule{2.5cm}{0.4pt}
\end{center}
\phantomsection
\addcontentsline{toc}{section}{\texorpdfstring{II.\quad Videnskabslære}{II. Videnskabslære}}
\markboth{Videnskabslære}{Videnskabslære}

\begin{center}{\bfseries \S\,12.}\end{center}

Hvor Erkjendelseslæren ender i tankeforvirrende eller blot negative Resultater
(Sophistik, Skepticisme), kan der naturligviis ikke være Tale om sand Videnskab. Som
Erkjendelseslæren er, maa Videnskabslæren blive. Drager Erkjendelseslæren en kritisk
Grændse mellem Phænomen og Væsen (\textit{Kritik der reinen Vernunft}), vil
Videnskabslæren paa alle dens Enemærker minde om Grændsen, og som i en \textit{reservatio
mentalis} altid underforstaae Umuligheden af at erkjende „Tingen i sig.“ Culminerer
Erkjendelseslæren i subjectiv Idealisme (\textit{Wissenschaftslehre}), ville de
videnskabelige Lærebygninger, forsaavidt de ere Principets Værk, forvandle sig til
Reflexionsspeile for „det rene Jeg“. Hæver Erkjendelseslæren sig til den objective
Idealismes absolute Viden (\textit{Phänomenologie des Geistes}), ville Sandsetingene
falde, men Væsensbegreberne staae som dialektiske Realiteter.

Paa sand Indsigt i Erkjendelsens Betingelser beroer den sande Indsigt i Videnskabernes
Væsen. Da nu Erkjendelsens Betingelser ikke blot ere indvortes subjective, men tillige
udvortes objective, er Videnskabens Væsen i dets forskjellige Gjennembrud og
eiendommelige Tilsyneladelser derhos tillige saa nøie sammenvoxet med givne Tiders
Vilkaar, med Livets Formaal, med Historiens Begivenheder, at man for at forstaae
Theorien maa see paa Tidsalderen.

Hvad Philosophien var for den aristokratiske Oldtid, hvad Theologien var for den
hierarchiske Middelalder, det er Physiken for den borgerlige Nutid: Culturhistorien og
Videnskabernes Historie belyse hinanden.

Hvorledes Bevægelser i det Udvortes forholde sig til Bevægelser i det Indvortes, vise
Opdagelserne i det 15de, og Reformationen i det 16de Aarhundrede. For at bedømme de
intellectuelle Fremskridt maa man imidlertid skjelne imel-

% ---- printed p.69 (PDF 77) ----
lem de Impulser og Vækkelser, hvorved store Hovedtanker fremkaldes, og den stille,
udholdende Aandsvirksomhed, der prøver Synspunkterne og giver Indholdet videnskabelig
Udarbeidelse.

Kan end Videnskaben stundom indslumre, fordi Aanderne sløves i det Overleverede, og
Tidsalderen mangler inspirerende Elementer; saa kan dog ogsaa paa den anden Side
Tænkningen overspændes og Kræfterne forvilde sig, naar Gjæringen er for stærk og
Indtrykkene for overvættes (Cardanus, Telesius, Campanella, Bruno.)

Hvad Videnskabsinteressen for sig angaaer, kan Begeistringen enten udelukkende
fastholde theoretiske Idealer (Copernikus, Le Verrier), eller holde sig nærmere til
praktiske Resultater (J. F. Böttiger, Daguerre); begge Retninger fremgaae af
Videnskabens Aand, thi det Theoretiske vil bevise sin Realitet i det Praktiske.

Men uagtet Videnskaben kun lever et kraftigt Liv, naar den finder sig i sund og sand
Forstaaelse med Livets Magter, kan den dog ikke drages ned i Midlernes Række, og
underordnes Endelighedens egennyttige Hensyn. Kun den Videnskab dyrkes, der dyrkes for
dens egen Skyld. At ville det Videnskabelige blot for Nyttens Skyld er netop ikke at
ville det Videnskabelige. Egennytten bedrager sig selv; den vil kun Midlet; men Midlet
udebliver, hvor Aanden mangler: det blotte Haandværksblik seer ikke saa dybt, at det
seer de nye Muligheder.

Hvor videnskabelig Indsigt derimod i enhver Retning erkjendes som aandig Magt, og uden
Hensyn til Egennyttens smaalige Hvorfor? omfattes med Uendelighedens, med Idealitetens
Interesse, vil Videnskaben ikke blive staaende som bleg theoretisk Skygge, men liig et
frodigt Træ bære gode, for Livet gavnlige Frugter. Naturvidenskabens Opdagelser i
nærværende Aarhundrede (Galvani og Volta; Ørsted og Humboldt) afgive en Række
storartede Exempler paa sand Vexelvirkning af theoretiske og praktiske Interesser.

\begin{center}\rule{0.35\linewidth}{0.4pt}\end{center}

% ---- printed p.70 (PDF 78): A. Videnskabernes Videnskab ----
\begin{center}
  {\large\bfseries A.}\par\vspace{2pt}
  {\bfseries Videnskabernes Videnskab.}\par\vspace{4pt}
  \rule{2cm}{0.4pt}
\end{center}
\phantomsection
\addcontentsline{toc}{subsection}{\texorpdfstring{A.\quad Videnskabernes Videnskab}{A. Videnskabernes Videnskab}}

\begin{center}{\bfseries \S\,13.}\end{center}

Da Videnskaben har sin Rod i Menneskelivet, hvor Phantasiens Idealer gaae forud for
Tænkningens, er det naturligt, at de forskjellige Videnskaber fra først af komme til
Verden i Indbildningens Svøb, og ere behæftede med uvidenskabelige Tilsætninger.
Saaledes var Astronomien længe sammenvoxet med Astrologie, Chemien med Alchemie,
Philosophien med Mythologie.

Vistnok ligger det i Videnskabens Drift at hæve Forpupningen, bryde sit Hylster og
udrense de fremmedartede Bestanddele. Men Dannelsen koster Tid, Udrensningen
Anstrengelse; de Hindringer, Menneskeaanden her har at overvinde, ere af forskjellig
Art.

a) En \emph{udvortes Hindring} fremgaaer af den naturlige Strid imellem Egoismen og
Sandheden. Den intellectuelle Sandhedskjærlighed har til alle Tider ført Kamp, ikke
blot mod Tidsaldernes Fordomme og Mængdens Tankeløshed, men fremfor Alt mod deres
Egoisme, hvis Magt og Indflydelse saae sig truet af det opgaaende Lys. Saalænge en
given Tilstand er af den Beskaffenhed, at Befolkningens aandelige Formyndere have
Interesse af at holde paa det Usande, ville de ikke let overbevise sig om det Sande.
Vistnok er Videnskaben objectiv; men det Subjective kommer jo med i Overbeviisningen:
Indsigten betinges ved Udsigten. Gaves der endnu i Landene anseelige Stillinger for
Astrologer; mon der da ikke ogsaa endnu skulde gives dem, som følte sig kaldede til at
studere Astrologie, naturligviis ikke for Anseelsens Skyld, men fordi de virkelig
havde overbeviist sig om det dybere Speculative i Astrologien? De, der ville Sandheden,
maae ivre for Sandheden; ogsaa Videnskaben har sine Marthyrer (Vanini, Galilei, J. G.
Fichte).

b) Af ædlere, skjøndt ikke mindre farlig Natur, ere de \emph{indvortes Hindringer}, den
rigtbegavede, overgribende

% ---- printed p.71 (PDF 79) ----
Subjectivitet kan berede sig selv. I begeistret Stræben efter dybere Indsigt hændes
det let, at Tænkeren forvexler den subjective Forestilling med den objective Gjenstand,
og saa mener at finde Det i Virkeligheden, som dog alligevel kun findes i hans egen
Phantasie. Hos de største Aander, de ædleste Naturer, en Plato (\textit{Timæus}), en
Keppler (\textit{Mysterium cosmographicum}), en Schelling, en Hegel, ville Blandinger
kunne paavises, hvorved deels i sig uvidenskabelige Elementer ere indvirkede i
videnskabelig Form, deels i sig videnskabelige Bestanddele findes anbragte i
uvidenskabelig, phantastisk Sammenhæng.

c) Hertil kommer endnu den Spænding mellem Videnskab og Livsanskuelse, som forøges ved
en stedse voxende \emph{Conflict mellem det Virkelige og det Ideale}.

Jo mere Videnskaben i endelig Forstand fastholder det Virkelige, desto lettere vil den
unegtelig undgaae det Phantastiske, men desto vanskeligere vil den paa den anden Side
tilfredsstille Personlighedens ideale Fordringer. Jo stærkere det videnskabelige
Opsving er mod det Ideale, desto lettere vil vistnok Videnskaben tilfredsstille
Aandsfordringen, men desto lettere vil den da ogsaa udsætte sig for at overfare
Virkelighedens Betingelser, og gribes i Illusioner.

Hvad er da Videnskab? Hvorpaa kjender man overalt og under alle Skikkelser det
Videnskabelige fra det Uvidenskabelige? Hvorledes bør den personlige Livsanskuelse
forholde sig til Videnskabens objective, upersonlige Resultater? Til Besvarelse af
disse Spørgsmaal kræves en omfattende Videnskabslære, en Videnskabernes Videnskab.

\begin{center}{\bfseries \S\,14.}\end{center}

Af Forholdet mellem Erkjendelseslæren og Videnskabslæren fremgaaer, at den sande
Videnskabernes Videnskab ikke kan være nogen anden end Philosophien. Hermed stemmer det
nøiagtig, at den sande omfattende Videnskabslære kun kan komme istand derved, at de
forskjelligartede Videnskabers Grundinteresser repræsenteres i de Forhandlinger, hvoraf

% ---- printed p.72 (PDF 80) ----
den fremgaaer. Under „den omfattende Discussion og Debat af Alt imod Alt“ (Sibbern) har
altsaa hver Videnskab sit Votum, og hver sit Bidrag. At Philosophien sammenarbeider,
klarer og organiserer disse Bidrag til et intellectuelt Hele, beviser, at den i saa
Henseende ikke afgjør Sagen paa egen Haand, men virker som Videnskabernes
Fælledsorgan, som deres egen Almeenbevidsthed, det Væsen, der er deres Slægtskab, det
Lys, hvori de kjende og erkjende hverandre indbyrdes.

Det er altsaa ikke anmassende Vilkaarlighed, naar Philosophien stiller de forskjellige
Lærebygninger frem for Videnskabslærens Domstol, og fordrer, at der med Hensyn paa
enhver især af dem skal gjøres Rede saavel for Bygningens Grundlag, som for
Bygningsmaaden; det er ikke selvtagen Myndighed, naar den efter omhyggelig afholdt
Prøvelse tillader sig at anvise de Theorier, som i Grundlag og Bygningsmaade have mere
Lighed med æsthetiske Luftcasteller end videnskabelige Lærebygninger „en egen Plads“
udenfor Videnskaben.

Philosophien gjør med alt dette kun sin Pligt, idet den, bunden ved Erkjendelsens Love,
ikke handler efter Lune, men som Videnskabernes lovlige Organ, Videnskabelighedens
underdanige Tjenerinde.

Sammenstille vi Aandsudviklingens store Repræsentanter: Philosophie, Theologie og
Physik med hverandre indbyrdes; da kan ikke alene den første, forsaavidt den stilles i
Modsætning til de tvende andre, ingenlunde ville fritage sig for videnskabelig
Retfærdiggjørelse, men fornuftigviis heller ikke de sidste.

Vilde Physiken, blændet af dens nuværende Popularitet, kaste Vrag paa den undersøgende
Videnskabslære, og istedetfor at agte paa det Tankens Lys, der vil gjennemskinne dens
Grund og Væsen, kun stole paa det umiddelbart Givne, og paaberaabe sig Sandsernes
uimodsigelige Autoritet, vilde den kun derved stille sig selv lavere end den
middelalderlige

% ---- printed p.73 (PDF 81) ----
Theologie. Ogsaa Theologien kan tale med om det Givne, benytte det Populære, og
paaberaabe sig en uimodsigelig Autoritet; og saalænge der i Menneskenaturen er Gnist af
et Høiere, vil Troens Magtsprog altid kunne opveie Sandselighedens Magtsprog.

Hvor da saaledes Supranaturalismens og Naturalismens Magtsprog, snart vexelviis, snart
deelviis, snart i aabenbart Fjendskab, snart i mystisk Overeenskomst, træder i
Videnskabslærens Sted og derved giver sig videnskabelig Vigtighed, er det jo vistnok
naturligt, at Philosophien enten ligefrem foragtes, eller dog betragtes som et
intellectuelt Legetøi; men saa er ogsaa Videnskabeligheden i Fare.

Men vil Theologien da maaskee ingen Videnskabslære? Dersom man paa een Gang kunde holde
sig ved Anselmus og dog profitere af Cartesius, kunde faae \textit{credo ut intelligam}
til at sige: \textit{de omnibus dubitandum est}, og omvendt, \textit{de omnibus
dubitandum est} til at forkynde sig som det rette \textit{credo, ut intelligam}, skulde
Theologien vist gjerne forene sig med Philosophien om at anprise en grundig og sund
Videnskabslære; men paa disse Betingelser er det Philosophien en Umulighed at forene
sig med Theologien, da Sagen selv er en Modsigelse.

Uagtet Philologie, Kritik og Historie, som i og for sig høre under Videnskaben,
udgjøre Bestanddele af Theologien, udgjøre de dog ingenlunde det i Theologien
Theologiske. Forsaavidt de nævnte Discipliner behandles videnskabelig, behandles de
utheologisk; forsaavidt de derimod behandles theologisk, behandles de uvidenskabelig.
Det i Theologien Theologiske kjendes nemlig derpaa, at den Aarsag, hvoraf alle
Virkninger i Naturens som i Aandens Rige deels umiddelbart, deels middelbart afledes og
forklares, ikke søges i Tingenes Natur eller Fornuftens Væsen, men i den absolute
Villie, i den Almægtige, der skaber af Intet. Men selv den reneste, skarpsindigste og
conseqventeste Anvendelse af et

% ---- printed p.74 (PDF 82) ----
saadant Princip (Occasionalismen; Malebranche), gjør Videnskaben uvidenskabelig.

Naar der i Videnskabslæren handles om et Princip, kan Sagen ikke afgjøres ved et Mere
eller Mindre, som om En ved at følge „Middelveien“ vilde mene, at man endda nok kunde
komme videnskabelig igjennem ved at antage Lidt af det Overnaturlige, naar man kun ikke
gjorde formeget af det.

Spørgsmaalet om et Princip er et Spørgsmaal om en ubetinget, Alt betingende
Forudsætning; et saadant Spørgsmaal maa derfor altid besvares med Ja eller Nei. Enten
maa da Theologien een Gang for alle vedkjende sig den overnaturlige Causalitet, og har
saa med det Samme bortstødt det objectiv rationelle Grundlag, eller ogsaa een Gang for
alle vedkjende sig dette Grundlag, og har saa med det Samme fraskrevet sig ethvert
Tilhold i det Overnaturlige. I første Tilfælde vil Theologien altsaa ophøre at være
Videnskab, i sidste Tilfælde at være Theologie.

\begin{center}{\bfseries \S\,15.}\end{center}

At Liv og Videnskab ikke kunne forholde sig udvortes og ligegyldig til hinanden, følger
umiddelbart af begges inderlige Sammenhæng med Cultur og Humanitet; at Liv og Videnskab
desuagtet ere selvstændige, beroer paa en principiel Adskillelse af det Subjective og
det Objective.

I Videnskaben have subjective Stemninger: Tilbøielighed, Utilbøielighed, Frygt, Haab o.
desl. Intet at betyde. Forsaavidt her spørges om det Virkelige, ville vi ikke vide,
hvorledes dette Virkelige efter vor Mening helst burde være, men hvorledes det i sig
er, og nødvendigviis maa være; forsaavidt her spørges om det Mulige, spørges her ikke
om, hvorvidt dette Mulige bedst tilfredsstiller vore Forhaabninger, men hvorvidt det
stemmer med Tilværelsens Betingelser.

Selv naar Videnskaben arbeider i Livsinteressernes Tjeneste, virker den i Kraft af det
Objective. Det er ikke subjective Følelser, ikke fromme Ønsker, men Indsigt i Tingenes

% ---- printed p.75 (PDF 83) ----
Væsen, der afleder Lynstraalen, og sætter Dampvognen i Bevægelse.

Vistnok er Videnskab umulig uden Subjectivitet; men Subjectiviteten er tjenende. Her
sandses, her føles, her tænkes; men det i Sandsning og Tænkning blot Subjective maa
overalt udrenses som det Tilfældige, det Forstyrrende, det Usande; hver Forudsætning er
objectiv. Selv der, hvor Indholdet berører Subjectiviteten saa umiddelbart som i det
Æsthetiske, seer Videnskaben kun det Sande, forsaavidt den seer det i Sagen
Uimodsigelige, det af vilkaarlig Forkjærlighed Uafhængige, det Almeengyldige.

Fra Livets Standpunkt hersker derimod Subjectiviteten; her er alt det Objective
tjenende. I uendelig Vexelvirkning med Omverdenen virker Personligheden, selv i
Hengivenhed og Opoffrelse, for subjective Interesser, subjective Formaal, subjective
Idealer. Overalt hvor Subjectiviteten her er i Forhold til Andet, er den igjennem dette
Andet i et eiendommeligt Forhold til sig selv. Paa Selvhed, Individualitet, Charakter,
Personlighed lægger Livet uendeligt Eftertryk. At Subjectiviteten blandt flere mulige
Opgaver nu ogsaa vælger Videnskaben, ophæver ingenlunde Modsætningen af Liv og
Videnskab. Eet er det at leve for en Videnskab, et Andet at lægge Videnskaben til Grund
for Livet.

Forsaavidt Subjectiviteten i stigende Conflict med den objective Verden udvikler de
personlige Kræfter, kunne vi her adskille følgende Hovedtrin.

a) \emph{Den relative Subjectivitet.} Paa dette Trin bliver det Objective ligelig sat
som Middel og som Maal. Ved at sætte det som Maal erkjender Subjectiviteten sin Grændse
og forholder sig tjenende; ved at sætte det som Middel tilfredsstiller den Selvet, og
hævder Friheden. Af Livsinteresse for Selvet og Værket, for Arbeideren og Arbeidet vil
Subjectiviteten da her træffe mangehaande Overeenskomster med det Objective, vove sig i
mangen Strid, lide mange Nederlag og vinde mange Seire, idet den gjennem-

% ---- printed p.76 (PDF 84) ----
løber Udviklingens Skala fra det Nyttige til det Ophøiede, fra den borgerlige Industrie
til den heroiske Bedrift.

b) \emph{Den uendelige Subjectivitet.} Ved at gjennembryde den objective Skranke og
overstige det Endeliges Grændse stræber Menneskets Bevidsthed at svinge sig op til
uendelig Frihed. Da Virkeligheden ikke kan negtes, og en materiel Omverden ikke
tilintetgjøres, maa Selvbefrielsen være ideel. I kraftig Eensidighed vælges mellem
tvende hinanden modsatte Udveie. Deels bliver det Objective under alle Skikkelser
saaledes nedsat til et blot Middel, at den nydende Individualitet i hvert Øieblik selv
kan gjøre sig til Tingenes Herre (Hedonisme; Aristipp); deels bliver det Objective som
det Almeenfornuftige saaledes indavnet i Bevidstheden, at alt Endeligt forvandles til
et Ligegyldigt; luttret og hærdet i uendelig Resignation, vil den handlende
Subjectivitet bestemme sig selv, idet den bestemmes af Skjæbnen (Stoicisme; Zeno —
Epictet).

En indre negativ Uendelighed, en dialektisk Eenhed af sædelig Alvor og gjennemført
Ironi er disse Eensidigheders dybere Forudsætning og uopnaaede Ideal. (Sokrates).

c) \emph{Den absolute Subjectivitet.} Den formelle Selvbefrielse ender i det Negative;
Subjectiviteten søger Livet, men finder Døden. Kun det Liv, der overvinder Døden, er
det sande, det evige Liv. At Individets medfødte Egoisme er dets egen oprindelige
Skyld, en Synd, som fortjener Døden; at den rene, skyldfrie Subjectivitet har frivillig
baaret Synden, og frelst de Skyldige fra Døden; at det naturlige Væsen er idel Afmagt,
kun Aandens Villie ubetinget Magt, at det Nærværende kun er Midlet, det Tilkommende
Maalet: slige Forudsætninger fordre Tro, men udelukke Viden.

I aabenbar Strid med Alt, hvad videnskabelig Indsigt paa nogen Maade kan anerkjende som
objectivt, er Troen saa langt fra at opnaae nogen Understøttelse hos Videnskaben, at
den just ved Sammenstød med den objective

% ---- printed p.77 (PDF 85) ----
Viden maa stødes tilbage i sig selv, og tilstaae, at dens eneste Grund er en
Subjectivitetens Trang, en indre uafviselig Fordring. Paa dette Trin er alt det
Objective forvandlet til Gjennemgangspunkt for det Subjective. Her handles ikke om
Livets Forhold til det Livløse, men om de Levendes Forhold til den absolut Levende; her
spørges ikke om Villiens Forhold til det Villieløse, men om de Villendes Forhold til
den absolut Villende.

Livet er subjectivt, det evige, sande Liv absolut subjectivt; Villien er subjectiv, den
evige, ubetingede Villie absolut subjectiv; det Objective er altsaa som objectivt ikke
det Sande: „Subjectiviteten er Sandheden“ (S. Kierkegaard).

\begin{center}{\bfseries \S\,16.}\end{center}

Idet Striden mellem Tro og Viden saaledes er dreven til sit Yderste, er den
intellectuelle Forsoning med det Samme indledet: de absolut ueensartede Principer
kunne, i væsentlig Forstand, ikke gjøre hinanden Afbræk. Den Tro, der hader Viden,
fordi den maa styrke sig ved Uvidenhed, er en falsk Tro. Virkeligheden vil med; det
nytter ikke at lukke Øinene: gaaer Jorden virkelig rundt, da gaaer den jo ligesaavel
rundt med de Troende som med de Ikke-Troende. Idet Troen erklærer, at Subjectiviteten
er Sandheden, taber den vistnok al Betydning for Videnskaben; men, idet Videnskaben
erklærer, at Objectiviteten er Sandheden, fraskriver den sig med det Samme al
Indflydelse paa Troen.

Tro og Viden kunne da uden Modsigelse forenes i een Bevidsthed; deres indbyrdes Strid
hidrører ikke fra det Principielle. Striden er psychologisk; det er den menneskelige
Skrøbelighed, der hverken har Resignation nok til at forsone sig med Videnskabens
Objective, eller Inderlighed nok til at fastholde Troen.

At Troen hader Videnskaben, er usandt; at den gjør Subjectiviteten uafhængig af
Videnskaben, er derimod sandt; At sætte Troen øverst og dog beholde Videnskaben, er
muligt; at sætte Videnskaben øverst, og dog bevare Troen, er umuligt.

% ---- printed p.78 (PDF 86) ----
Men uagtet Videnskaben selv maa indsee, at den hverken kan anerkjende den absolute
Subjectivitet eller foreskrive den Love, er det dog alligevel i dens Interesse at faae
denne Subjectivitet belyst saa alsidig som muligt.

Adskillelsen af Videnskabeligt og Uvidenskabeligt vanskeliggjøres derved, at det, som
ifølge sin Natur maa forkastes af Videnskaben, ikke desto mindre bliver Gjenstand for
videnskabelig Behandling, og saaledes dog paa en Maade gaaer med i Videnskaben. Galskab
er t. Ex. uforenelig med Videnskab; ikke desto mindre blive Galskabssymptomerne og de
Gales Behandling en vigtig Gjenstand for Psychologie og Lægevidenskab. Som et Foster af
overtroisk Phantasie er Mythen i Strid med det Videnskabelige; ikke desto mindre gives
der en Mythevidenskab, en videnskabelig Mythologie.

Kan det Allerforskjelligste stilles under videnskabelig Belysning; kan endog det meest
Subjective blive objectivt, forsaavidt det objectiveres i Reflexionen, da kan der fra
dette Synspunkt heller ikke indvendes Noget imod, at der danner sig en
Subjectivitetsvidenskab, en Videnskab om det Subjective.

At en saadan Subjectivitetsvidenskab (Psychologie, Æsthetik, Ethik,
Religionsphilosophie) imidlertid ikke maa forvexles med, hvad man ellers pleier at
kalde Troesvidenskab, lader sig indsee, naar vi tydeliggjøre os Forskjellen imellem det
i videre og det i snævrere Forstand Videnskabelige.

Videnskabeligt i videre Forstand er Alt, hvad der bliver Gjenstand for Videnskabens
Behandling: det Videnskabelige ligger i Behandlingen; videnskabeligt i snævrere Forstand
er derimod kun det, hvis Gyldighed Videnskaben indseer, hvis Sandhed den kan bedømme,
prøve og anerkjende.

Ved at benytte den Tvetydighed, der ligger i Kategorien: „Videnskabeligt“ lykkes det
undertiden Troeslæren at give sig Skin af Videnskab; det er da Videnskabslærens Opgave
at paavise de vigtigste herhen hørende Sophismer.

\emph{Første Sophisme:} „Skal Troen ikke beroe paa subjective Indfald, og priisgives
Individernes Luner, maa

% ---- printed p.79 (PDF 87) ----
det udtrykkelig kunne paavises, hvad man skal troe, og paa hvem man skal troe: Troen
maa have sin Gjenstand. At undersøge Troesgjenstandene grundig er at undersøge dem
videnskabelig. Troeslæren beroer paa en saadan Undersøgelse; altsaa er Troeslæren en
Videnskab.“

Denne Sophisme giver Sagen det Udseende, at Troesgjenstandene sikkres ved Videnskaben,
fordi de prøves i Videnskaben. Sagen forholder sig omvendt. Troesgjenstandene opløse
sig i Dunst, naar Videnskaben først faaer Lov at anbringe sin Kritik. Det er den
videnskabelige Kritik ligesaa umuligt at anerkjende Evangeliet om den Opstandne, som
det er Videnskaben umuligt at anerkjende Mythen om Herkules. At anerkjende et Mirakel i
Kraft af videnskabelig Kritik er en Modsigelse, da Muligheden af et Mirakel gjør
videnskabelig Kritik umulig: hvor Underet taler, maa Forstanden tie, hvor Miraklet
begynder, maa Videnskaben ende.

\emph{Anden Sophisme:} „Enhver eiendommelig Videnskab maa nødvendigviis have sine
Forudsætninger. Ligesom Logiken forudsætter og retfærdiggjør Tænkningens Gyldighed,
Physiken Erfaringens Gyldighed og Æsthetiken det Skjønnes Gyldighed, saaledes maa
Troeslæren for sit Vedkommende forudsætte og retfærdiggjøre Gyldigheden af Troen og
Troens Indhold. Den, der beviser, at Troeslæren ikke er en Videnskab, fordi den har
sine særegne Forudsætninger, beviser med det Samme, at der slet ingen Videnskab gives;
men her gjælder det: \textit{qui nimium probat, nihil probat}. Altsaa er Troeslæren en
Videnskab.“

Denne Sophisme blænder ved den Tvetydighed, der ligger i at medbringe en Forudsætning.
Hvad der gjør Troeslæren uvidenskabelig, er ingenlunde den Omstændighed, at den
medbringer Forudsætninger, men netop den Omstændighed, at de Forudsætninger, den
medbringer, stride mod Videnskaben. Den forudsætter Miraklet, den forudsætter
Subjectivitetens og den subjective Forvisnings absolute Gyldighed; med disse
Forudsætninger har den saa at sige væltet

% ---- printed p.80 (PDF 88) ----
sig ind paa Videnskaben, og brudt den videnskabelige Maalestok. Skulle Conseqvenserne
være videnskabelige, maae Forudsætningerne ogsaa være det; paa uvidenskabelig Grund kan
ingen videnskabelig Bygning opføres: Naturvidenskab er mulig, men Mirakelvidenskab er
umulig.

\emph{Tredie Sophisme:} „Troeslærens overnaturlige Forudsætninger have, hver for sig,
vistnok kun Gyldighed for den Troende; men imellem disse Sætninger indbyrdes gives der
en logisk og forsaavidt objectiv Sammenhæng. Da Sammenhængen er objectiv, og
Troessætningerne netop finde deres sande Udtydning og dybere Begrundelse i denne
Sammenhæng, er Begrundelsen selv objectiv, og Troeslæren altsaa i Kraft af denne
Begrundelse en Videnskab.“

Det Blændende i denne Sophisme beroer derpaa, at en logisk Sammenhæng uden videre
bliver til en objectiv Sammenhæng, og denne igjen til en objectiv Begrundelse. „Alle
Mennesker ere Triangler; Cajus er et Menneske, altsaa er Cajus et Triangel. I ethvert
Triangel er Vinklernes Sum liig to Rette; Cajus er et Triangel, altsaa er Vinklernes
Sum i Cajus liig to Rette.“ Imellem disse Sætninger er her unegtelig en logisk
Sammenhæng, men er her derfor en objectiv videnskabelig Begrundelse? Fra Livets
Standpunkt kan der være en Evighedsforskjel imellem det Hellige og det Profane, det
Evangeliske og det Mythiske; men, naar Videnskaben nu eengang ikke kan bedømme dette
Hellige, ikke magte dette Evangeliske, er Forskjellen i Videnskabens Øine, som om den
slet ikke var.

Da Troessætningerne hver for sig ere ubetinget subjective, ere de i Videnskaben hver
for sig ligesaa vilkaarlige, som de meest eventyrlige Phantasier. Gives der altsaa en
videnskabelig Begrundelse af Troeslærdomme, saa gives der ogsaa en videnskabelig
Begrundelse af Mythelærdomme, en videnskabelig Begrundelse af Nonsens.

\begin{center}\rule{0.35\linewidth}{0.4pt}\end{center}

% ---- printed p.81 (PDF 89): B. De særskilte Videnskaber ----
\begin{center}
  {\large\bfseries B.}\par\vspace{2pt}
  {\bfseries De særskilte Videnskaber.}\par\vspace{4pt}
  \rule{2cm}{0.4pt}
\end{center}
\phantomsection
\addcontentsline{toc}{subsection}{\texorpdfstring{B.\quad De særskilte Videnskaber}{B. De særskilte Videnskaber}}

\begin{center}{\bfseries \S\,17.}\end{center}

At Philosophien kalder sig selv Videnskabernes Videnskab, maa ikke misforstaaes, som om
den hermed vilde udgive sig for et Videnskabernes Orakel, den eneste, egentlige og
sande Videnskab. Hvorledes en saadan Misforstaaelse opkommer, hvilke Forvirringer den
medfører, bør paaagtes.

Enhver Viden bestemmes ved sin Gjenstand, det særegne Stof, Erkjendelsen bearbeider. Et
omfattende Vidensapparat, en planmæssig erhvervet Kundskabsrigdom giver Lærdom, men
Lærdom er endnu ikke Videnskab; først med Principernes Udvikling og Stoffets rationelle
Gjennemarbeidelse kommer det Videnskabelige. Men, endskjøndt der til en vis Grændse
ligesaavel kan gives Videnskab uden Lærdom, som Lærdom uden Videnskab, vil enhver
indtrængende Erkjendelse af virkelige Kjendsgjerningers Sammenhæng dog altid udkræve en
egen Forening af Lærdom og Videnskab. Alle saakaldte positive Videnskaber ere lærde
Videnskaber.

Det Specielle er altid underordnet det Almene; men for at berige vor Indsigt i det
Almene, maae vi begynde med det Specielle. Jo mere en Hovedvidenskab udvider sig, desto
flere Bifag vil den udsondre. Det gaaer med Videnskaben som med Alexanders Erobringer;
jo større Riget bliver, desto nærmere er det Vendepunkt, da Ringen skiftes, og Delingen
foregaaer. Skal nu en saadan Deling for Videnskabens Vedkommende ikke betegne Rigets
Undergang, men en Forfatningens Fremgang, maae de mindre Afdelinger strengt underordnes
de større, og disse det organiserede Hele: Totaliteten er Sandheden.

Bliver nu Sandhedstotaliteten opfattet som apriorisk-metaphysisk Væsenhed, saa forstaaes
heraf, hvorledes den

% ---- printed p.82 (PDF 90) ----
absolute Videns Philosophie har kunnet sætte sig som den ene Vidende paa Thronen og see
ned paa de lærde Videnskaber, som paa sine Undergivne. Lærdom er Middel, Videnskab er
Philosophie. De særskilte Videnskaber maae nu enten arbeide som Haandlangere, eller see
at lære Philosophie, og forklare sig Speculationen.

„Speculation er Alt, er i Gud at skue, at betragte det, som er. Videnskaben selv har
kun forsaavidt Værd, som den er speculativ, det er Contemplation af Gud, som han er.
Men den Tid vil komme, da Videnskaberne mere og mere ville ophøre, og den umiddelbare
Erkjendelse indtræde. Kun i den høieste Videnskab lukker det dødelige Øie sig, idet ikke
mere Mennesket seer, men den evige Seen selv er seende i Mennesket.“ (Schelling).

Men saameget disse og lignende Forkyndelser end vidne om Sammensmeltning af religiøs og
videnskabelig Begeistring, saa kan dog en saadan apokalyptisk Henviisning til et
Vendepunkt, da „den evige Seen selv“ vil blive „seende i Mennesket,“ ikke erstatte
Videnskaben den Mangel paa Besindighed og methodisk Forstandighed, der er en
ufravigelig Betingelse for Fremskridt i Viden.

Eensidigheden har, som Erkjendelseslæren viser, sin dybere Grund i den Miskjendelse af
Forestillingselementet, der charakteriserer den objective Idealisme.

Ligesom Lyset i givne Tilfælde maa forudsætte en Flamme, Flammen en Forbrænding, og
Forbrændingen brændbare Stoffer; saaledes maa den absolute Viden vistnok forudsætte
forskjelligartede Forestillingsstoffer, men kun for at lade dem forbrænde i den
Dialektik, hvoraf Begrebet lyser. Bliver nu Forestillingsstoffet paa denne Maade
behandlet, ikke som et intellectuelt Bygningsmateriale, endnu mindre som Tankernes Kjød
og Blod, men kun som Brændsel; da følger det af sig selv, at man ikke tager det saa nøie
med Kjendsgjerningens enkelte Stykker.

% ---- printed p.83 (PDF 91) ----
Følgen heraf er naturligviis Feiltagelser i det Enkelte og Vilkaarligheder i
Opfatningen af det Totale; men, selv naar disse Misgreb rettes, vil Standpunktet
ligefuldt være falsk.

For at vurdere Forestillingselementet, erindre man, at det intellectuelle Lys ligesaa
lidt som det physiske kan afgive Farver og Synsbilleder, med mindre det brydes i de
forskjellige Media, tilbagekastes af de forskjellige Gjenstande, opfanges og
tydeliggjøres i bestemte Indtryk. Selv, naar Fremstillingen af de til et positivt
Indhold svarende Kategorier nærmest holder sig i det Almindelige, mærker man en
væsentlig Forskjel imellem de golde, for Sagerkjendelse blottede Abstractioners Tomhed
og det Farvespil, den naturlige Friskhed, hvormed Begreber udspringe, der ere
gjennemvirkede af Anskuelse, og vidne om Autopsie.

\begin{center}{\bfseries \S\,18.}\end{center}

I sin Afhandling „\textit{Ueber den Begriff der Wissenschaftslehre}“ (\textit{Sämmtliche
Werke, Berlin 1845}) har Fichte opkastet det Spørgsmaal, hvorvidt Videnskabslæren kan
være forvisset om, at den i det Hele har udtømt den menneskelige Viden. Besvarelsen af
dette Spørgsmaal reduceres til følgende Cirkel: „Naar Sætningen $X$ er den første,
høieste og absolute Grundsætning for den menneskelige Viden, saa er den menneskelige
Viden et eneste System; thi det Sidste følger af Sætningen $X$. Da der nu skal være et
eneste System i den menneskelige Viden, saa er Sætningen $X$, der virkelig (ifølge den
opstillede Videnskab) begrunder et System, den menneskelige Videns Grundsætning
overhovedet; det paa den begrundede System er hiint den menneskelige Videns eneste
System. Ud over denne Cirkel kan man ikke komme.“ Endvidere hedder det (S. 65) om
Forholdet mellem Videnskabslæren og de særskilte Videnskaber: „Man indseer, hvorfor kun
Videnskabslæren har absolut Totalitet, medens alle særskilte Videnskaber ville være
uendelige.

% ---- printed p.84 (PDF 92) ----
Videnskabslæren indeholder blot det Nødvendige; alle øvrige Videnskaber gaae paa
Friheden saavel med Hensyn til vor egen Aand, som til den af os uafhængige Natur.“

Fra den absolute Videns Standpunkt har Hegel givet en Fremstilling af de forskjellige
Videnskabers indbyrdes Sammenhæng.

„I en sædvanlig Encyklopædie tages Videnskaberne empirisk, som de forefindes. De skulle
deri fuldstændig opregnes, og ordnes derved, at det Lignende, det under fælleds
Bestemmelser Sammentræffende stilles sammen. Den philosophiske Encyklopædie er derimod
Videnskaben om den nødvendige, ved Begrebet bestemte Sammenhæng, og om den philosophiske
Tilbliven af Videnskabernes Grundbegreber og Grundsætninger.“

„Videnskaberne ere efter deres Erkjendelsesmaade enten empiriske eller reent rationelle.
Absolut betragtede skulle begge have samme Indhold. Det er den videnskabelige Stræbens
Maal at hæve det, der vides empirisk, til det altid Sande, til Begrebet, at gjøre det
rationelt, og derved indlemme det i den rationelle Videnskab.“

„Videnskabernes rationelle Udvidelse er med det Samme en Udvidelse af Philosophien selv.
Videnskabens Hele deler sig da i 3 Hoveddele: 1) \emph{Logik}. 2) \emph{Naturvidenskab}.
3) \emph{Aandsvidenskab}. Naturens og Aandens Videnskaber kunne da betragtes som de
\emph{reale} eller særskilte Videnskabers System til Forskjel fra den rene Videnskab
eller Logiken, efterdi de ere den rene Videnskabs System i Skikkelse af Naturen og
Aanden.“

Som Overgang fra Logik til Naturvidenskab opføres Mathematiken med Hensyn paa Rummet og
Tiden. Naturvidenskaben selv falder i: a) \emph{Mechanik} med Hensyn til Materie og
Bevægelse. b) \emph{Det Uorganiskes Physik} med Hensyn paa Lys, Magnetisme,
Electricitet, de physiske Legemer og den chemiske Proces. c) \emph{Det Organiskes
Physik}, hvorunder da igjen hører Geologien, der be-
% ---- printed p.85 (PDF 93) ----
handler Jorddannelsen; Botanik og Plantephysiologie, hvis Gjenstand er den vegetabile
Natur; Physiologie, \emph{comparativ} Anatomie, Zoologie, der behandle den sunde, animale
Natur, ligesom Medicinen den syge. Ogsaa Aandsvidenskaben falder i tre Afsnit: a) Aanden
i sit Begreb, \emph{Anthropologie}, \emph{Psychologie} (Sproget — \emph{Philologie}.) b) Den
praktiske Aand, der yttrer sig i Daad og Handling; herunder \emph{Retslæren},
\emph{Statsvidenskaben}, \emph{Historien}. c) Aandens Fuldendelse i Kunst, Religion og
Videnskab, hvortil svare \emph{Æsthetik}, \emph{Religionslære}, \emph{Philosophie}.

Men uagtet de særskilte Videnskaber saaledes opstilles mod hverandre indbyrdes og imod
Philosophien, er deres ideelle Selvstændighed dog alligevel kun et Skin. Encyklopædien
„er egentlig en Fremstilling af Philosophiens almindelige Indhold; thi hvad der i
Videnskaberne er grundet paa Fornuft, afhænger af Philosophien, hvad der derimod er i dem
af det, der beroer paa vilkaarlige og udvortes Bestemmelser, eller, som det hedder, er
positivt og statutarisk, ligger udenfor den.“

Fra sit kantiske Standpunkt har Apelt (\textit{die Theorie der Induction}, S. 65—70)
angivet det Principielle i Adskillelsen af Videnskabernes forskjelligartede
Grundcharakteer.

„Der gives,“ hedder det, „to forskjellige Kilder, hvoraf vore Erkjendelser udspringe,
Sandsen og den rene Fornuft. Fra hiin kommer Sandseerkjendelsen, fra denne de rene
Fornufterkjendelser. De sidste falde i Mathematik og Philosophie. Der gives som Følge
heraf tre forskjellige Slags Erkjendelser: \emph{empiriske}, \emph{mathematiske},
\emph{philosophiske}. Enhver af disse tre Erkjendelsesmaader har sin særegne Charakteer:
Tilfældighed og Anskuelighed tilkommer den \emph{empiriske}, Anskuelighed og Nødvendighed
den \emph{mathematiske}, Nødvendighed uden Anskuelighed den \emph{philosophiske}.“

„I de mathematiske Erkjendelser blive vi os Tingenes Sammensætningsformer, i de
philosophiske deres Sammen-
% ---- printed p.86 (PDF 94) ----
knytningsformer \textit{in abstracto} bevidste. Kant kaldte hine: Former for den figurlige
Synthesis, d. e. anskuelige Forbindelsesformer, disse: Former for den intellectuelle
Synthesis, d. e. kun tænkelige Forbindelsesformer, Love.“

„Betragtningen af de forskjellige Systemformers Forhold til hinanden og til de forskjellige
Erkjendelsesarter fører til en videnskabelig Architektonik, d. e. til en Lære om et System
af alle menneskelige Videnskaber.“

„De tre forskjellige Erkjendelsesmaader svare til de tre forskjellige Systemformer. Den
rene philosophiske Erkjendelsesform hører under det kategoriske System; den rene
mathematiske Form under det hypothetiske System, og den rene empiriske Gehalt under det
conjunctive System. Philosophien er en Erkjendelse af blotte Begreber. Men med en saadan
Erkjendelsesmaade er ingen anden systematisk Afleden mulig end ved en kategorisk Indordnen
af særegne Tilfælde i Reglens Almindelighed. Mathematiken er en Erkjendelse ved Begreber af
reen Anskuelse. I denne Anskuelse er det mindst Sammensatte det Oprindeligste og den Grund,
hvoraf nye Følger afledes ved bestandig nye Sammensætninger. Det Sammensattes Afledning af
det Usammensatte skeer her ved hypothetisk Tankeforbindelse. Endelig har den rene Empirie
kun med Opfatning af det enkelte ved Sandseanskuelsen Givne at gjøre. Hver Kjendsgjerning
beroer paa sig selv, og den ene staaer ved Siden af den anden i et conjunctivt System, der
sammenordner Dele i et Hele.“

Det i denne Deduction kritisk Træffende er at anerkjende; men det viser sig her, som
overalt, at det blot Kritiske uden sand Egenhed med det Dialektiske, er i Philosophien
ufrugtbart, og fører til Schematisme. Den Adskillelse af Sandsen og den rene Fornuft,
hvorpaa denne hele Architektonik er begrundet, har i Virkeligheden ingen Betydning.
Adskillelsen beroer paa en Abstraction, men det Abstraherede vil igjen integreres ved det,
hvoraf det er abstraheret.
% ---- printed p.87 (PDF 95) ----
At Sandseanskuelsen er ueensartet med Tænkningen, er uimodsigeligt; men disse Ueensartede
ere for Reflexionen dialektiske Momenter i den samme Realerkjendelse. Det Sandselige og det
Tænkelige, det Individuelle og det Almene ere i Tingen ikke det Samme og dog det Samme. I
Virkeligheden er Alt concret. Virkelighedserkjendelsen betinges af Concretioner; men de
intellectuelle Concretioner betinges igjen af de abstracte Formers Integration.

Ingen, der har endog blot et almindeligt Begreb om Differential- og Integralregning, vil,
naar han reflecterer philosophisk, kunne negte det dialektiske Forhold mellem den tænkte
Linie, Fladens negative Grændse, den bestemte Negation, der ikke kan integreres, og den
anskuede Linie, Rectanglets Differential, den positive Grændse, der kan integreres til et
endeligt Rectangel (Jfr. Philosophie og Mathem. S. 65—67; 75—76).

Den immanente Dialektik vil opløse de særskilte Videnskabers Selvstændighed, og forvandle
dem til Momenter i Philosophiens System; den afgrændsende Kritik vil indskrænke Philosophien
til en egen Rubrik, og mener i det Hele at hjælpe ved Schemata og Rubriker. Begge
Bestræbelser ere eensidige. Skjøndt en Videnskab for sig, kan Philosophien ikke ved nogen
kritisk Demarkationslinie afsondres fra de øvrige Videnskaber; og endskjøndt de øvrige
Videnskaber have hver i sig det Almeenvidenskabelige, have de det dog ikke saaledes, at
„hvad der i Videnskaberne er grundet paa Fornuft“ derfor skulde kunne siges at være
„afhængigt af Philosophien.“

Mathematikens Forskjel fra Philosophien bestemmes ikke ligefrem ved Qvantitetens Anskuelse
og Ikke-Anskuelse, men ved den eiendommelige Analyse, Mathematiken underkaster det
Qvantitative. (Phil. og Mathem. S. 9—14).

Hvad de positive Videnskaber angaaer, er deres Selvstændighed umiddelbart givet med det for
enhver især eiendommelige Stof, uden at det dermed er givet, at Philoso-
% ---- printed p.88 (PDF 96) ----
phien ikke ogsaa paa sin Maade skulde kunne tilegne sig det samme Stof. Jo haardere
Modsætningen er imellem Objectet og Reflexionen, desto fjernere vil den positive Videnskab
ligge for „den absolute Videns“ Philosophie. Det gjør da her en væsentlig Forskjel, om
Tænkningens givne Object er æsthetisk, historisk eller physisk.

Det æsthetiske Object, det givne Kunstværk, er jo vistnok altid fra en eller anden Side
beslægtet med et Sandseligt, og maa forsaavidt opfattes empirisk. Men det Empiriske er her
saa gjennemskinnet af Bevidsthed, at det Aandige hersker. Den objective Idealisme har da
ogsaa uden videre vindiceret sig Æsthetiken som Deel af Philosophien. Imidlertid har
Æsthetiken med sine Specialanalyser, sin Autopsie, sin Kunsthistorie o. s. v. en anden
Side, hvorfra den, uden at opgive sin videnskabelige Charakteer, vel kan siges at indtage
en selvstændig Stilling ligeoverfor Philosophien.

At Historien med sine Kjendsgjerninger og sit Lærdomsapparat staaer endnu selvstændigere,
er indlysende; ikke destomindre har det historiske Object, som tilhørende
Forbigangenheden, Aandslivet og Erindringen, endnu et saa ideelt Anstrøg, at Idealismen
selv inden disse Enemærker kan bevæge sig, som om den kun var paa sit Eget. Kun Physiken
er haard, staaer ubetinget imod, og tillader ingen idealistiske Overgreb. Mod det physiske
Object, mod Virkelighedens Polygon støder „den absolute Viden“ an, og forraader sin
Svaghed.

\begin{center}{\bfseries \S\,19.}\end{center}

\emph{Historien fortæller}, \emph{Philosophien reflecterer}: her er en Grændse, som ikke kan
jævnes ved et Mere eller Mindre. Vistnok kan Fortællingen optage Reflexionen, og omvendt;
men derfor er alligevel hiin et Historiens, denne et Philosophiens Mærke.

At fortælle er at bestemme det Almene ved det Individuelle saaledes, at Tanken dramatiseres
i Begivenheder, og det Indre gaaer op i det Ydre: her er Forestillings-
% ---- printed p.89 (PDF 97) ----
elementet overveiende. At reflectere er at bestemme det Almene saaledes, at Begivenheder
afgive Begreber, at Anskuelsen fordybes i det Indre, og Forestillingen løser sig op i
Tanken.

Af denne Forskjel i Fremstillingsmaaden er en Forskjel i Behandlingen af Heelheden og dens
Dele en naturlig Følge. Den fortællende Historie bevæger sig gjennem Delenes Enkeltheder
Skridt for Skridt fremad til en middelbar Anskueliggjørelse af det Hele; den reflecterende
Philosophie begynder strax med Indblik i det Totale, opfatter de særskilte Dele som
Væsensmomenter i det Heles Indhold, og kan ikke give sig hen i Detaillen. Selv naar
Philosophien allermeest interesserer sig for det Enkelte, er dette Enkelte altid
gjennemskinnet af det Totale. Philosophien kan fra forskjellige Sider belyse det Totales
Selvformidling gjennem det Specielle; men den taaler ikke det Specielles Umiddelbarhed, og
kan ikke følge den fortællende Historie ad Middelbarhedens lange Omvei. Philosophien har,
som Historiens Philosophie, et andet Formaal, end den lærde Historie; men af samme Grund
har den lærde Historie da ogsaa et andet Formaal, end Philosophien.

Er her nu i Henseende til Formaalet og Midlerne, Opgaven og dens Behandling en principiel
Forskjel imellem det Historiske og det Philosophiske, da kunne Historien og Philosophien
heller ikke skride frem i convergerende Linier, saa at de nærme sig til en Sammensmeltning,
efterhaanden som de nærme sig til en Fuldendelse. Fremstillinger, hvori det Historiske og
det Philosophiske ligelig opveie, og da ogsaa paa en Maade halvere hinanden, ere kun
Mellemformer, neutrale Gebeter, som med deres midlertidige Frugtbarhed ikke vedkomme
Principerne.

Den lærde Historie og Historiens Philosophie kunne middelbart fremhjælpe hinanden,
corrigere og berige hinanden, og derved komme i en stedse inderligere Forstaaelse; men just
denne Forstaaelse vil stadfæste Forskjelligheden.
% ---- printed p.90 (PDF 98) ----
For at skille Sandt fra Falsk, maa Historien være kritisk; for at paavise Begivenhedernes
Aarsagsforbindelse, pragmatisk; for at finde og fremstille den Aand, hvis Interesser og
Plan Begivenhederne tjene, philosophisk. I alle tre Henseender kan den lærde Historiker vel
behøve philosophisk Dannelse; men hans Aandsvirksomhed gaaer i ingen Henseende ind under
Philosophien. Overalt henviist til Kjendsgjerninger og sammenvoxet med det Givne vil
Historikerens Tænkning, langt fra at tabe sig i det philosophisk Almene, tvertimod
articuleres ved det Anskuelige, og fuldendes i kunstnerisk Individuation: Historieskrivning
er en Kunst.

Har en Historiker gjort et eller andet Hovedparti til Gjenstand for et alsidigt
indtrængende Kildestudium, er han derved kommen en forgangen Virkelighed saa nær, at han
ligesom seer den Ansigt til Ansigt; mon da ikke hans Phantasie vil være saa opfyldt af
Billeder, og hans Sind saa gjennemkrydset af individuelle Indtryk, at han neppe uden
forudgaaende Gjæring kan ordne sit Stof og paabegynde sin Fremstilling af det Anskuede?

Væsentlig seet beroer Spændingen paa et eget Vexelforhold mellem det Virkelige og dets
Afbildede. I Afbilledet skal det Virkelige gjengives; med al Adskillelse af Væsentligt og
Uvæsentligt kan her ikke, som ellers i en Kunst, tænkes paa æsthetisk Idealisering. Det er
en virkelig Personlighed, en virkelig Stat, en virkelig Tidsalder, hvis Physiognomie skal
sees i givne Træk, usminket som det er, med Rynker og Fregner. Og dog skal Aanden være med,
og præge sig i Afbilledet, at det ikke skal overlæsses med det Trivielle, og løbe ud i det
Virkeliges endeløse Vidtløftighed.

Med en saadan Opgave for Øie vil Historikeren indsee, at der ikke gives nogen særegen Form,
der fuldkommen udtømmer Virkelighedens Indhold, men at der i Fremstillingen uophørlig
strides mellem Tab og Vinding. Hvad her i een Henseende er Vinding, er i en anden Henseende
% ---- printed p.91 (PDF 99) ----
Tab; dette gjælder ikke blot om Anlæget i det Hele, men gjentager sig endog i Udførelsen af
det Enkelte. At Partiaanden ikke altid behøver at gjøre Vold paa Kjendsgjerningerne, for at
faae Historien paa sin Side, skylder den vistnok nærmest sin egen sophistiske Smidighed, men
dog ogsaa væsentlig den Omstændighed, at det i Historien Objective, om det end fastholdes
nok saa strengt, nødvendigviis erholder et Farveskjær af den objectiverende Bevidsthed.
Forskjellige Fremstillinger af det samme Emne kunne ifølge Individuationen være lige
berettigede.

Men er den historiske Form saa kunstnerisk individuel, at een og samme almindelige Opgave
(t. Ex. den franske Revolutionshistorie) først finder sin fuldstændige Løsning, naar den
løses paa flere eiendommelige Maader; saa indsees heraf, i hvilken Grad
Forestillingselementet maa protestere imod en Sammensmeltning af den lærde Historie og
Historiens Philosophie.

\begin{center}{\bfseries \S\,20.}\end{center}

Det er ikke Tilfældets Lune, og ikke menneskelig Vilkaarlighed, der tør drage Grændsen
imellem de særskilte Videnskaber. Hver Videnskab er i sig en væsentlig Totalitet; dens
Grændser, Afdelinger og Bearbeidelse betinges af Stoffet, beherskes af Tanken, og bestemmes
ved den omfattende Opgave.

I enhver Videnskabs Methode sees det Almeenvidenskabelige: overalt en Vexelbestemmelse af
det i Erkjendelsen Umiddelbare og Middelbare, af Omvei og Gjenvei, af nødvendig
Sagbevægelse og fri Tankebevægelse (Philosophie og Mathematik. S. 44—46); ligeledes vil
enhver Videnskabs systematiserende Stræben udkræve Definitioner, Inddelinger, lavere og
høiere Eenheders Underordnen under hinanden indbyrdes og under en høieste Eenhed o.s.v. Men
det særegne Stof og det særegne Øiemed, hvori det bearbeides, vil dog ikke destomindre i
Methoden som i Systemdannelsen
% ---- printed p.92 (PDF 100) ----
medføre saa afgjørende Eiendommeligheder, at hver Videnskab derved ikke alene bliver en
Lærebygning i særegen Stiil, men endog paa indre uendelig Maade organiseres til en relativ
Monade, en intellectuel Individualitet.

I hvilken Grad det eiendommelige Stof indvirker paa Systemdannelsen, vil især fremlyse af de
Videnskaber, hvis Indhold er udtalt i typiske Naturformer. Her er i Mineralriget den døde
Form: Krystallen; i Planteriget den levende Form: Organismen; i Dyreriget den besjælede
Form: Individet.

a) \emph{Den døde Form}. I sin typiske Reenhed er den substantielle Begrændsningsform
mathematisk. Ved at lægge et mathematisk Legeme (f. Ex. Tærningen) til Grund for
Bestemmelsen af Axerne og det tilsvarende Snit, kan den rene Krystallographie med en exact
Videnskabs Sikkerhed afrunde sig til et i særskilte Systemer, — det sphæroedriske, det
rhomboedriske, det pyramidale, det prismatiske — articuleret Hele. Herved bliver det da ikke
alene muligt at tilbageføre en uendelig Mangfoldighed af afledede Specialformer til deres i
hvert System gjældende Grundform; men selv de Combinationer, Ændringer og Forqvaklinger, som
hidrøre fra Dannelser i den frie Natur, ere saa langt fra at bringe det Hele i Forvirring, at
de tvertimod blive forstaaelige ved de rene Typers Systematik (Anvendt Krystallographie).

Anderledes, naar der ikke blot spørges om de substantielle Begrændsningsformer, men tillige
om de tilsvarende Substansers chemiske og physiske Egenskaber, om Mineralets Charakteristik.
Her kommer det overveiende Hensyn til den chemiske Substantialitet (Berzelius) i aabenbar
Strid med et ligesaa overveiende Hensyn til den i de physiske Egenskaber udtalte
Individualitet (Mohs), den strenge Fordring paa systematisk Begrundelse i Strid med Schemaet
og dets mange Rubriker.

Et ligeligere Hensyn til chemiske og physiske Charakteermærker kan nu vistnok føre til
sindrige Combinationer af
% ---- printed p.93 (PDF 101) ----
System og Schema; men først efterhaanden som Videnskaben trænger ind i den væsentlige
Sammenhæng af Stoffernes chemiske og physiske Egenskaber, vil det Schematiske gaae op i det
Systematiske.

b) \emph{Den levende Form}. Endskjøndt den virkelige Grændse, der adskiller det Organiske fra
det Uorganiske, ikke umiddelbart lader sig paapege, er Væsensforskjellen mellem Mineralet og
Planten ikke desto mindre tydelig nok. Medens Krystallen, den døde Form, gaaer op i det
Mathematiske, er det Mathematiske (Tretalsplante, Femtalsplante, Stængelbladenes
Differensvinkel o. s. v.) derimod som mathematisk af underordnet Betydning i de levende
Former. I de forskjellige Arter saavelsom i de enkelte Exemplarer af hver Art gjennemløber
det Organiske en saadan Række af Metamorphoser og tilsvarende Articulationer, at her med
Rette kan tales ikke blot om uophørligt Stofskifte, men ogsaa om lavere og høiere Trin.

Sexualsystemet, det saakaldte kunstige Plantesystem, med dets Klasser og Ordener (Linné) er
et storartet, omfattende, i det Hele praktisk anvendeligt Schema; men først, idet en Fleerhed
af væsentlige Egenskaber og Organforhold tages med i Inddelingens Begrundelse, begynder det
naturlige, d. e. det egentlige, System at danne sig. Streng Fastholden af et metaphysisk
Princip og Vilkaarlighed i Opfatningen af det Givne (Reichenbach) har imidlertid, istedetfor
at bringe Systemet til Fuldendelse, tvertimod viist, hvilke Vanskeligheder her endnu ere at
overvinde.

c) \emph{Den besjælede Form}. Medens den individuelle Eenhed for Plantens Vedkommende enten
slet ikke er, eller dog er saa proteusagtig, at den forsvinder i Phyter og Celler, er Dyret
derimod i Kraft af Begrebets nødvendige Fordringer (Fornemmelse, Vilkaarlighed) oprindelig
sat som bestemt Individ. Den strengere Articulation, hvorved Organismens enkelte Led
uddannes for sig, gaaer i det Hele frem med en stigende Centraliseren, der samler Organernes
% ---- printed p.94 (PDF 102) ----
forskjellige Virksomheder i den fornemmende, sig selv bestemmende Livseenhed.

Det er da en Selvfølge, at Arternes Charakteermærker og den typiske Adskillelse i lavere og
høiere Ordener, — \emph{Bløddyr}, \emph{Leddyr}, \emph{Hvirveldyr} — maa i Dyreriget optræde
med en ganske anden Klarhed og orienterende Gyldighed end analoge Inddelinger i Planteriget.
Streng Empirie (Cuvier) er ogsaa her traadt i Modsætning til naturphilosophisk Speculation
(Geoffroi St. Hilaire). Men, for at Systemet kan bringes til Fuldendelse, maa
Slægtsbegrebets Realitet først afgjøres, og de høiere Klassers Sammenhæng med de lavere
bestemmes saaledes, at Hovedformer og Mellemformer vise sig som naturnødvendige Led i de
besjælede Eenheders opadstigende Kjæde.

Er nu det Almeenvidenskabelige overalt ved Stof og Opgave sammenvoxet med det Særegne, da
maa hver Videnskab ogsaa udelukkende see paa sit Ideal, og kan umulig tage sin
Videnskabelighed til Lehn af Philosophien.

Hvor lidt Philosophien er berettiget til at tilskrive sig den Tænkning og de
Begrebsudviklinger, der tilhøre de lærde Videnskaber, sees udtrykkelig deraf, at Kategorier
som: Substans, Individ, Art, Slægt o. desl. udvikles i Realvidenskaberne paa en ganske anden
Maade, end i en objectiv Logik eller speculativ Naturphilosophie.

Vel ere Tankelovene overalt de samme; men at fastholde Tankelovene for sig i deres rene
Almindelighed er just at holde dem i en betænkelig Eensidighed. Den strenge Tænkning har i
den ene Videnskab en ganske anden Charakteer end i den anden: de særskilte Videnskaber
forudsætte særskilte Talenter.

\begin{center}{\bfseries \S\,21.}\end{center}

Forsaavidt de særskilte Videnskaber i Stof og Formaal have en endelig Begrændsning, ere de
Fagvidenskaber. Hvor eiendommelig de forskjellige Fagvidenskabers Synsmaade dan-
% ---- printed p.95 (PDF 103) ----
nes, fremlyser, naar de ved at behandle et og samme Emne nærme sig til gjensidig Berøring. I
Bestemmelsen af Sjæl, Liv og Livsprincip viser sig en gjennemgribende Forskjel imellem den
strengt empiriske Physiologie, for hvilken Organismen er „et levende Maskinerie“, der ikke
engang behøver en Livskraft til „Maskinmester“ (G. Valentin), og den paa Selviagttagelse
grundede Psychologie, der opfatter Sjælen som constituerende Væsenhed, og lader Livet
opvælde af „en indre Kilde“.

Den Strid, der saaledes kan opkomme mellem virkelige Fagvidenskaber, maa dog for Alting ikke
sammenblandes med Striden mellem Videnskabslæren og en Mirakelvidenskab. Eet er
Videnskabens Strid med det Uvidenskabelige; et Andet er Striden mellem videnskabelige
Eensidigheder. Eet er en Strid mod forvirrende Paahæng af en Uberettiget; et Andet en Strid
for sand Ligevægt mellem Berettigede. Eet er det, at Archæologen og Geognosten, udrustede
hver med sit Apparat af Kjendsgjerningsbeviser, mødes i Strid ved en Kæmpehøi; et ganske
Andet, at en middelalderlig Theolog, uden Kjendskab til de naturlige Mærker, beraaber sig
paa et Overnaturligt, og blander sig i Undersøgelsen om Menneskeracerne og „vore første
Forældre“.

For at den enkelte Videnskab kan skride frem, maa den understøttes af dens nærmest
Beslægtede. Idet Mathematiken er en Hovedvidenskab for sig, er den tillige Hjælpevidenskab
for Mechanik og Astronomie. Chemie, Mineralogie og Geognosie forene sig som
Hjælpevidenskaber i Geologien. De hjælpende Videnskaber hjælpes igjen af dem, de selv have
hjulpet. I hvilken Grad Botaniken og Zoologien fremhjælpe Geologien, og igjen fremhjælpes af
denne, er iøinefaldende. I den store Kreds af sideordnede Opgaver har intet Fag udelukkende
Hævd paa at være Hovedfag: i Videnskabernes Organisme er Midlet Formaal og Formaalet Middel.

% [text to be added: pp. 96--204 — continue II. Videnskabslære (§21 cont., B. De særskilte Videnskaber → C. Philosophien og de særskilte Videnskaber) onward]

\end{document}
